\documentclass[11pt,openright]{book}
% ---------------------------------------------------------------
%  Preamble: default LaTeX fonts (Computer Modern), ISO 80000-2 style
% ---------------------------------------------------------------
\usepackage[b5paper,hmargin=19mm,top=20mm,bottom=22mm,headheight=16pt]{geometry}
\usepackage{emptypage}             % blank pages before odd-page chapter starts
\usepackage{amsmath,amssymb,amsthm,mathtools}
\usepackage{bm}
\usepackage[Euler]{upgreek}        % upright \uppi only (ISO: constants upright)
\usepackage{array,booktabs,multirow,tabularx}
\usepackage{enumitem}
\usepackage{xpatch}               % robust-command patches; loads etoolbox
\usepackage{xcolor}
\usepackage{tikz}
\usetikzlibrary{arrows.meta,calc,decorations.markings,patterns,positioning}
\usepackage{pgfplots}
\pgfplotsset{compat=1.18}
\usepackage[font=small,labelfont=md]{caption}
\usepackage{float}
\usepackage[most]{tcolorbox}
\usepackage{esvect}
\usepackage[colorlinks=true,linkcolor=blue!50!black,citecolor=blue!50!black,urlcolor=blue!60!black]{hyperref}
\usepackage[capitalise,nameinlink]{cleveref}

% Body text: 16 pt baseline spacing at the existing 11 pt class font size.
% Keep the separate small-text and heading sizes and spacing unchanged.
\xpatchcmd{\normalsize}{13.6}{16}{}%
  {\PackageError{lecture-notes}{Could not set 16 pt body line spacing}%
    {Check the 11 pt book class definition of \string\normalsize.}}
\normalsize

\setlist{itemsep=2pt,topsep=4pt}
\setlength{\parskip}{3pt}
\renewcommand{\arraystretch}{1.18}
\setlength{\emergencystretch}{2em}  % the B5 measure is narrow: allow looser lines, not overfull ones
\raggedbottom                       % no stretched vertical gaps to fill a page

% ---------------- page layout for B5 ----------------
% Section headings ragged right, so long titles break instead of overflowing.
\makeatletter
% Regular-weight text for tablet reading. Patch text styles individually so
% mathematical bold (\mat and \one) keeps its matrix/scalar distinction.
\newcommand{\regulartextstyle}[1]{%
  \patchcmd{#1}{\bfseries}{\mdseries}{}%
    {\PackageError{lecture-notes}{Could not set regular-weight text style}%
      {Check the book class definition of \string#1.}}}
\regulartextstyle{\@makechapterhead} % chapter number
\regulartextstyle{\@makechapterhead} % chapter title
\regulartextstyle{\@makeschapterhead}
\regulartextstyle{\l@chapter}
\regulartextstyle{\subsubsection}
\regulartextstyle{\paragraph}
\regulartextstyle{\subparagraph}
\regulartextstyle{\descriptionlabel}
\regulartextstyle{\labelitemii}
\thm@headfont{\mdseries}
\renewcommand\section{\@startsection{section}{1}{\z@}%
  {-3.5ex \@plus -1ex \@minus -.2ex}{2.3ex \@plus.2ex}%
  {\normalfont\Large\mdseries\raggedright}}
\renewcommand\subsection{\@startsection{subsection}{2}{\z@}%
  {-3.25ex\@plus -1ex \@minus -.2ex}{1.5ex \@plus .2ex}%
  {\normalfont\large\mdseries\raggedright}}
\makeatother
% Running heads in a smaller size, so long chapter and section titles fit the line.
\usepackage{fancyhdr}
\pagestyle{fancy}
\fancyhf{}
\fancyhead[LE,RO]{\thepage}
\fancyhead[RE]{\footnotesize\slshape\leftmark}
\fancyhead[LO]{\footnotesize\slshape\rightmark}
\renewcommand{\headrulewidth}{0pt}

% ---------------- ISO 80000-2 conventions ----------------
\newcommand{\dd}{\mathop{}\!\mathrm{d}}      % differential d (upright)
\newcommand{\ee}{\mathrm{e}}                  % Euler's number (upright)
\newcommand{\ii}{\mathrm{i}}                  % imaginary unit (upright)
\newcommand{\cpi}{\uppi}                      % the number pi (upright)
\DeclareRobustCommand{\vec}[1]{\vv{#1}}       % vectors: arrow over italic letter
\newcommand{\mat}[1]{\bm{#1}}                 % matrices: bold italic
\newcommand{\one}{\mat{1}}                    % unit matrix
\newcommand{\T}{^{\mathrm{T}}}                % transpose
\newcommand{\herm}{^{\dagger}}                % hermitian conjugate
\DeclareMathOperator{\tr}{tr}
\DeclareMathOperator{\diag}{diag}
\DeclareMathOperator{\sgn}{sgn}
\DeclareMathOperator{\Ind}{Ind}
\DeclareMathOperator{\Span}{span}
\DeclareMathOperator{\Ker}{Ker}
\DeclareMathOperator{\Imag}{Im}
\DeclareMathOperator{\Real}{Re}
\newcommand{\SO}[1]{\mathrm{SO}(#1)}
\newcommand{\SU}[1]{\mathrm{SU}(#1)}
\newcommand{\Og}[1]{\mathrm{O}(#1)}
\newcommand{\U}[1]{\mathrm{U}(#1)}
\newcommand{\GL}[1]{\mathrm{GL}(#1)}
% Schoenflies point-group names: letter italic, label subscript upright
\newcommand{\Cv}[1]{C_{#1\mathrm{v}}}
\newcommand{\Ch}[1]{C_{#1\mathrm{h}}}
\newcommand{\Dh}[1]{D_{#1\mathrm{h}}}
\newcommand{\Td}{T_{\mathrm{d}}}
\newcommand{\Oh}{O_{\mathrm{h}}}
\newcommand{\Ci}{C_{\mathrm{i}}}
\newcommand{\Cs}{C_{\mathrm{s}}}
% symmetry operations (group elements italic; labels upright)
\newcommand{\sv}{\sigma_{\mathrm{v}}}
\newcommand{\sd}{\sigma_{\mathrm{d}}}
\newcommand{\sh}{\sigma_{\mathrm{h}}}
\newcommand{\Ebar}{\bar{E}}
% irreducible representations (Mulliken labels): upright
\newcommand{\irr}[1]{\mathrm{#1}}
% representations with fixed meaning (upright labels)
\newcommand{\GV}{\Gamma_{\mathrm{V}}}
\newcommand{\GR}{\Gamma_{\mathrm{R}}}
\newcommand{\GM}{\Gamma_{\mathrm{M}}}
\newcommand{\Gvib}{\Gamma_{\mathrm{vib}}}
\newcommand{\chiV}{\chi_{\mathrm{V}}}
\newcommand{\chiR}{\chi_{\mathrm{R}}}
\newcommand{\chiM}{\chi_{\mathrm{M}}}
\newcommand{\Ninv}{N_{\mathrm{inv}}}
% misc
\DeclarePairedDelimiter{\abs}{\lvert}{\rvert}
\DeclarePairedDelimiter{\norm}{\lVert}{\rVert}
\newcommand{\ket}[1]{\lvert #1\rangle}
\newcommand{\bra}[1]{\langle #1\rvert}
\newcommand{\braket}[2]{\langle #1\vert #2\rangle}
\newcommand{\inner}[2]{\langle #1, #2\rangle}
\newcommand{\Tgroup}{\mathcal{T}}
\newcommand{\seitz}[2]{\{#1\,|\,#2\}}
\newcommand{\TR}{\Theta}

% ---------------- theorem-like environments ----------------
\theoremstyle{plain}
\newtheorem{theorem}{Theorem}[chapter]
\newtheorem{lemma}[theorem]{Lemma}
\newtheorem{proposition}[theorem]{Proposition}
\newtheorem{corollary}[theorem]{Corollary}
\theoremstyle{definition}
\newtheorem{definition}[theorem]{Definition}
\newtheorem{example}[theorem]{Example}
\theoremstyle{remark}
\newtheorem{remark}[theorem]{Remark}

% Box titles are white on a dark bar, so links inside a title must be white too.
\newcommand{\titlelinks}{\hypersetup{linkcolor=white,urlcolor=white}}
\tcbset{
  lecbox/.style={colback=blue!4,colframe=blue!45!black,boxrule=0.5pt,arc=1mm,
                 left=6pt,right=6pt,top=4pt,bottom=4pt,fonttitle=\mdseries\small\titlelinks},
  keybox/.style={colback=orange!6,colframe=orange!60!black,boxrule=0.5pt,arc=1mm,
                 left=6pt,right=6pt,top=4pt,bottom=4pt,fonttitle=\mdseries\small\titlelinks},
  gapbox/.style={colback=green!5,colframe=green!45!black,boxrule=0.5pt,arc=1mm,
                 left=6pt,right=6pt,top=4pt,bottom=4pt,fonttitle=\mdseries\small\titlelinks}
}
\newenvironment{watch}[1]{\begin{tcolorbox}[lecbox,title={Video: #1}]\small}{\end{tcolorbox}}
\newenvironment{keyresult}[1][Key result]{\begin{tcolorbox}[keybox,title={#1}]}{\end{tcolorbox}}
\newenvironment{beyond}[1][Beyond the lectures]{\begin{tcolorbox}[gapbox,title={#1}]\small}{\end{tcolorbox}}

% Sections and subsections that go beyond the lectures carry a green star.
\DeclareRobustCommand{\bstar}{\textcolor{green!45!black}{$\bigstar$}}
\newcommand{\beyondmark}{\texorpdfstring{\ \bstar}{ *}}

% ---------------- exercises and solutions ----------------
\newcommand{\exercisesection}{%
  \section*{Exercises and solutions}%
  \addcontentsline{toc}{section}{Exercises and solutions}%
  \markright{\MakeUppercase{Exercises and solutions}}}
\newcounter{exercise}[chapter]
\renewcommand{\theexercise}{\thechapter.\arabic{exercise}}
\newenvironment{exercise}[1][]{\refstepcounter{exercise}\par\medskip\noindent
  \textmd{Exercise~\theexercise}\if\relax\detokenize{#1}\relax\else\ \textit{(#1)}\fi.\ }{\par}
% Each solution follows its exercise, set off by a grey rule on the left.
\newtcolorbox{solution}{breakable,enhanced jigsaw,blanker,left=9pt,
  borderline west={1pt}{0pt}{black!35},before skip=4pt,after skip=10pt,
  before upper={\noindent\textit{Solution.}\ }}
\crefname{exercise}{Exercise}{Exercises}

% ---------------- links to the video lectures ----------------
\newcommand{\lecid}[1]{%
  \ifcase#1\or iHRC-f4SPxs\or 19uuuiJiY8Y\or BUZ-ZwHDl6M\or wTN-b-uhbAE\or
  XYyzpEt\string_Vgk\or vUt\string_IAdrrN8\or oZ5iq\string_\string_MQgk\or rXGrZq4mZAY\fi}
% \ts{lecture}{seconds}{mm:ss}: clickable timestamp (youtu.be short link)
\newcommand{\ts}[3]{\href{https://youtu.be/\lecid{#1}?t=#2}{L#1\,@\,#3}}
\newcommand{\lec}[1]{\href{https://youtu.be/\lecid{#1}}{Lecture~#1}}


\hypersetup{pdftitle={Group Theory in Condensed Matter: Companion Notes},
            pdfsubject={Companion notes to the lecture series by J. L. Manes}}

\begin{document}
\frontmatter
\begin{titlepage}
\centering
\vspace*{3cm}
{\LARGE\bfseries Group Theory in Condensed Matter\par}
\vspace{0.4cm}
{\Large A Practical Introduction\par}
\vspace{1.2cm}
{\large Companion Notes to the Video Lectures of Prof.\ Juan Luis Ma\~nes\par}
\vspace{0.3cm}
{\normalsize (UPV/EHU, recorded Winter 2022; published on the \emph{Topological Matter School} channel)\par}
\vspace{2cm}
\begin{minipage}{0.8\textwidth}\small
These notes were written independently as a study companion. They follow the
order and the examples of the eight video lectures, but the text, derivations,
figures and exercises are our own. They are neither written nor endorsed by the
lecturer; any error is ours. Timestamps such as \ts{2}{2520}{42:00} are clickable
links to the corresponding moment in the video.
\end{minipage}
\vfill
{\small Playlist: \url{https://www.youtube.com/playlist?list=PLygH4pJ_4EhrbbYmYjEYlVRpaSrRRQTri}\par}
\end{titlepage}

\tableofcontents

\chapter*{Preface}
\addcontentsline{toc}{chapter}{Preface}
\markboth{\MakeUppercase{Preface}}{\MakeUppercase{Preface}}

\section*{What these notes are}
The lecture series \emph{Group Theory in Condensed Matter: A Practical Introduction}
teaches group theory \emph{backwards}: first you see what a group-theory computation
produces, then you learn to perform it, and only at the end you learn why it works.
The lecturer states the key theorems without proof and devotes the time to
applications, from molecular vibrations to topological quantum chemistry.

These notes keep that order, so that each chapter can be read next to the
corresponding lecture, and add three things:
\begin{enumerate}
\item a written, self-contained account of every lecture, with all computations
      carried out explicitly and checked numerically;
\item green boxes marked \emph{Beyond the lectures}, and sections marked \bstar{}
      that fill the gaps a theorist will need later: cosets and quotient groups
      (\cref{sec:cosets,sec:normal}), proofs of Schur's lemma and the orthogonality
      theorems together with projection operators
      (\cref{sec:unitarity,sec:proofs,sec:projectors}), the rotation group, $\SU{2}$
      and Clebsch--Gordan series (\cref{sec:su2}), time reversal and Kramers
      degeneracy (\cref{sec:TR,sec:TRcrystal}), and the method of invariants for
      $\vec{k}\cdot\vec{p}$ Hamiltonians (\cref{sec:kp});
\item exercises at the end of every chapter, each followed by its detailed solution.
\end{enumerate}
Each starred section is placed where the lectures first need it; it can be skipped
on a first reading without losing the thread.

\section*{How the chapters map onto the lectures}
\begin{center}
\small
\begin{tabularx}{\textwidth}{@{}lXll@{}}
\toprule
Chapter & Topic & Video & Length\\
\midrule
\ref{ch:intro} & Symmetries, degeneracies and representations & \lec{1} & 89 min\\
\ref{ch:building} & Irreducible representations as building blocks & \lec{2} & 92 min\\
\ref{ch:ops1} & Operations with representations I: tensors, invariants & \lec{3} & 91 min\\
\ref{ch:ops2} & Operations with representations II: spectroscopy & \lec{4} & 84 min\\
\ref{ch:spin} & Spin, double groups, crystal-field splitting & \lec{5} & 96 min\\
\ref{ch:space} & Space groups and their representations & \lec{6} & 90 min\\
\ref{ch:bands} & The group theory of electron bands & \lec{7} & 89 min\\
\ref{ch:topology} & Group theory and topology & \lec{8} & 88 min\\
\bottomrule
\end{tabularx}
\end{center}
Lectures~5 and~8 have no captions; their timestamps in these notes were read off
the slides and are approximate.

\section*{Prerequisites}
Linear algebra (matrices, eigenvalues, change of basis), quantum mechanics at the
level of the hydrogen atom and angular momentum, and an introductory solid-state
course (Bravais lattices, reciprocal lattice, Bloch's theorem, tight binding).
No previous knowledge of group theory is assumed.

\section*{Notation and typographical conventions}
The typesetting follows ISO~80000-2:
\begin{itemize}
\item Variables are italic; mathematical constants and operators with fixed meaning
      are upright: the imaginary unit $\ii$, Euler's number $\ee$, the number $\cpi$,
      the differential $\dd$ (as in $\dd x$), and named functions and operators
      ($\exp$, $\sin$, $\tr$, $\det$).
\item Vectors carry an arrow, $\vec{r}$, $\vec{k}$, $\vec{E}$; their components are
      italic scalars, $r_x$, $k_i$.
\item Matrices are bold italic, $\mat{D}(g)$; their elements are italic scalars,
      $D_{ij}(g)$. The transpose is $\mat{A}\T$ and the hermitian conjugate
      $\mat{A}\herm$.
\item Labels that are not variables are upright: the ``v'' in $\Cv{3}$, the ``d''
      in $\sd$, the ``g'' in $\irr{E}_{\mathrm{g}}$. Irreducible representations
      carry upright Mulliken symbols, $\irr{A}_1$, $\irr{E}$, $\irr{T}_2$, while the
      identity \emph{operation} is the italic $E$.
\end{itemize}

\begin{center}
\begin{tabular}{@{}ll@{}}
\toprule
Symbol & Meaning\\
\midrule
$G$, $\abs{G}$ & a group and its order (number of elements)\\
$g$, $h$ & group elements\\
$\Gamma$, $\Gamma^{(\mu)}$ & a representation; the $\mu$-th irreducible representation (irrep)\\
$\mat{D}(g)$ & matrix representing $g$\\
$\chi(g)=\tr\mat{D}(g)$ & character\\
$d_\mu$ & dimension of irrep $\mu$\\
$m_\mu$ & multiplicity of irrep $\mu$ in a representation\\
$\GV$, $\GR$ & vector and axial-vector (``rotational'') representations\\
$\GM$, $\Gvib$ & mechanical and vibrational representations of a molecule\\
$\oplus$, $\otimes$ & direct sum and (tensor/Kronecker) product\\
$[\Gamma^2]$, $\{\Gamma^2\}$ & symmetric and antisymmetric squares\\
$\Tgroup$ & translation group of a crystal\\
$\seitz{R}{\vec{v}}$ & Seitz symbol: point operation $R$ followed by translation $\vec{v}$\\
$G_{\vec{k}}$ & little group of $\vec{k}$\\
$\TR$ & time-reversal operator\\
\bottomrule
\end{tabular}
\end{center}
In the lectures, general representations are denoted by $T$ and irreducible ones by
$\tau$ (following Lyubarskii). We use $\Gamma$ throughout, because $T$ is also
needed for the translation group and for the tetrahedral group $\Td$.

Symmetry operations always act \emph{actively}: they move the physical object
(a vector, a wave function, a molecule) and leave the coordinate axes fixed
(\cref{sec:active}). A function transforms as $(\hat{O}_g f)(\vec{r})=f(g^{-1}\vec{r})$.

\section*{Recommended textbooks}
The lecturer recommends:
G.\,Ya.\ Lyubarskii, \emph{The Application of Group Theory in Physics}
(Pergamon, 1960), which he learned from;
M.\,S.\ Dresselhaus, G.\ Dresselhaus and A.\ Jorio, \emph{Group Theory: Application
to the Physics of Condensed Matter} (Springer, 2008), rich in applications; and
C.\,J.\ Bradley and A.\,P.\ Cracknell, \emph{The Mathematical Theory of Symmetry in
Solids} (Oxford, 1972), the reference for tables and for magnetic groups but not a
book to learn from. For the topological part: J.\ Cano and B.\ Bradlyn,
\emph{Band representations and topological quantum chemistry},
Annu.\ Rev.\ Condens.\ Matter Phys.\ \textbf{12}, 225 (2021), arXiv:2006.04890.
The tables in these notes use the conventions of the Bilbao Crystallographic Server
(\url{https://www.cryst.ehu.es}) wherever possible.

\mainmatter
\chapter{Symmetries, degeneracies and representations}\label{ch:intro}

\begin{watch}{\lec{1}}
Why group theory \ts{1}{51}{0:51} $\cdot$
continuous vs discrete symmetries \ts{1}{193}{3:13} $\cdot$
methane: what group theory predicts \ts{1}{824}{13:44} $\cdot$
the idea of a representation \ts{1}{1394}{23:14} $\cdot$
Wigner's theorem \ts{1}{2278}{37:58} $\cdot$
the group $\Cv{3}$ and the vector representation \ts{1}{2528}{42:08} $\cdot$
definition of a representation \ts{1}{2784}{46:24} $\cdot$
reducible vs irreducible \ts{1}{4302}{71:42} $\cdot$
equivalent representations \ts{1}{4990}{83:10}
\end{watch}

\section{Why group theory?}

Group theory is the systematic way of extracting \emph{all} the consequences of the
symmetry of a physical system, reliably and efficiently. Its computations are simple
(mostly sums and products of integers) but the concepts need care.

In condensed matter we deal almost exclusively with \emph{space-time} symmetries:
rotations, reflections, inversion, lattice translations and time reversal. Internal
symmetries ($\U{1}$, $\SU{2}$ isospin, $\SU{3}$ colour) belong mostly to high-energy
physics. A second classification is more important for the method:

\begin{description}
\item[Continuous symmetries] such as $\SO{3}$ are Lie groups. A finite rotation is
the exponential of an infinitesimal one,
\begin{equation}\label{eq:rotexp}
  \hat{U}(\vec{\theta}) = \exp\!\Bigl(-\frac{\ii}{\hbar}\,\vec{\theta}\cdot\hat{\vec{L}}\Bigr),
\end{equation}
so invariance of the Hamiltonian, $\hat{U}\hat{H}\hat{U}\herm=\hat{H}$ for all
angles, implies at first order $[\hat{H},\hat{\vec{L}}]=0$: a \emph{conservation law}.
Degeneracies and selection rules then follow from the conservation law and the
algebra of the generators, without any group theory. The $(2l+1)$-fold degeneracy
of the levels of a central potential and the dipole selection rule
$\Delta l=\pm1$ are derived this way in every quantum-mechanics course.
\item[Discrete symmetries] such as the finite set of rotations and reflections that
leave a molecule or a crystal invariant have no infinitesimal elements, hence no
generators and no conservation laws. There is no Noether theorem for them.
To go from symmetry to degeneracies and selection rules we must jump directly,
and the bridge is the representation theory of finite groups: characters, their
orthogonality, projection operators.
\end{description}

\begin{center}
\begin{tikzpicture}[node distance=1.6cm,>=Stealth,every node/.style={font=\small}]
\node (s1) {continuous symmetry};
\node[below=0.9cm of s1] (c1) {conservation law};
\node[below=0.9cm of c1] (d1) {degeneracies, selection rules};
\draw[->] (s1)--node[right]{Noether} (c1);
\draw[->] (c1)--node[right,font=\footnotesize]{algebra of generators} (d1);
\node[right=3.4cm of s1] (s2) {discrete symmetry};
\node[below=2.3cm of s2] (d2) {degeneracies, selection rules};
\draw[->,thick,blue!60!black] (s2)--node[right,text width=2.4cm,font=\footnotesize]{representation theory of finite groups} (d2);
\end{tikzpicture}
\end{center}

In this course we never need Lie-algebra techniques such as roots, Dynkin diagrams
or Young tableaux. From the rotation group we only need its irreducible representations,
which we obtain elementarily in \cref{ch:spin}. All the groups we meet are discrete
subgroups of $\Og{3}$ (rotations, reflections, inversion), lattice translations, and
time reversal.

\section{What group theory predicts: the vibrations of methane}\label{sec:methane}

The methane molecule $\mathrm{CH_4}$ has its four hydrogen atoms at the vertices of a
regular tetrahedron with the carbon atom at the centre. Its symmetry group, the
tetrahedral group $\Td$, has 24 elements.

Without symmetry considerations we can only say the following. Each of the five
atoms can move in three directions, so a distortion of the molecule is a vector with
$15$ components. In the harmonic approximation the frequencies are obtained by
diagonalising a real symmetric $15\times15$ matrix, which gives at most $15$ distinct
frequencies. Six of them vanish: three uniform translations and three rigid
rotations do not stretch any bond (these are the \emph{zero modes} of a non-linear
molecule). So at most nine non-zero frequencies remain.

A group-theory computation, which after \cref{ch:building} will take you five minutes
with pencil and paper, says much more:

\begin{keyresult}[What group theory says about $\mathrm{CH_4}$ and $\mathrm{CH_3D}$]
\begin{itemize}
\item $\mathrm{CH_4}$ has exactly \textbf{four} distinct vibrational frequencies,
      with degeneracies $1$, $2$, $3$ and $3$ ($1+2+3+3=9$).
\item All four are Raman active; only the two triply degenerate ones are infrared
      active; the molecule has no microwave (pure rotational) spectrum.
\item In deuterated methane $\mathrm{CH_3D}$ each triply degenerate frequency splits
      into one non-degenerate and one doubly degenerate frequency: six distinct
      frequencies, three non-degenerate and three doubly degenerate, all of them
      Raman and infrared active; $\mathrm{CH_3D}$ is also microwave active.
\end{itemize}
\end{keyresult}

The result of the computation is written as
\begin{equation}\label{eq:CH4result}
\begin{aligned}
  \GM  &= \irr{A}_1\oplus\irr{E}\oplus\irr{T}_1\oplus 3\,\irr{T}_2 &&(\dim 15),\\
  \GV  &= \irr{T}_2 &&(\dim 3),\qquad \GR = \irr{T}_1 \quad(\dim 3),\\
  \Gvib &= \GM\ominus\GV\ominus\GR = \irr{A}_1\oplus\irr{E}\oplus 2\,\irr{T}_2 &&(\dim 9).
\end{aligned}
\end{equation}
Every symbol here is a \emph{representation}: a rule telling how some kind of object
transforms under the symmetry operations. Loosely:
\begin{itemize}
\item $\GV$, the \emph{vector representation}, tells how ordinary (polar) vectors
      transform: a set of $3\times3$ matrices, one per symmetry operation.
\item $\GR$, the \emph{rotational} or \emph{axial-vector representation}, tells how
      pseudovectors (angular momentum, magnetic field) transform. It coincides with
      $\GV$ on proper rotations and differs by a sign on improper operations
      (reflections, inversion), because $\vec{L}=\vec{r}\times\vec{p}$ does not change
      sign under inversion while $\vec{r}$ and $\vec{p}$ do.
\item $\GM$, the \emph{mechanical representation}, tells how an arbitrary distortion
      of the molecule (a 15-component vector) transforms: $15\times15$ matrices.
\item $\Gvib$, the \emph{vibrational representation}, acts on the 9-dimensional
      space of pure vibrations, obtained by removing uniform translations (which
      transform as vectors) and rigid rotations (which transform as axial vectors).
\end{itemize}
The symbols $\irr{A}_1$, $\irr{A}_2$, $\irr{E}$, $\irr{T}_1$, $\irr{T}_2$ are the five
\emph{irreducible representations} (irreps) of $\Td$. Unlike $\GV$ or $\GM$ they do not
describe any particular object; they are abstract building blocks. Their dimension is
read off the letter: $\irr{A}$ and $\irr{B}$ are one-dimensional, $\irr{E}$ is
two-dimensional, $\irr{T}$ is three-dimensional.

\begin{keyresult}[How to read a decomposition (Wigner's theorem)]
If the representation carried by the eigenvectors of a symmetric operator (here
the dynamical matrix, later the Hamiltonian) decomposes as
$\Gamma=\bigoplus_\mu m_\mu\,\Gamma^{(\mu)}$, then
\begin{itemize}
\item the \textbf{multiplicity} $m_\mu$ is the number of distinct eigenvalues
      associated with the irrep $\Gamma^{(\mu)}$;
\item the \textbf{dimension} $d_\mu$ of the irrep is the degeneracy of each of them.
\end{itemize}
\end{keyresult}
For $\mathrm{CH_4}$, $\Gvib=\irr{A}_1\oplus\irr{E}\oplus2\,\irr{T}_2$ gives one
non-degenerate, one doubly degenerate and two triply degenerate frequencies. For
$\mathrm{CH_3D}$ the symmetry group is the smaller group $\Cv{3}$ and
$\Gvib=3\,\irr{A}_1\oplus3\,\irr{E}$ (\cref{sec:mechanical}). We prove Wigner's
theorem in an explicit example in \cref{sec:wigner}; the general proof rests on
Schur's lemma (\cref{sec:schur}).

\begin{remark}
Group theory tells you which frequencies must be degenerate, how many distinct ones
there are, and which matrix elements must vanish. It never tells you the numerical
values of the frequencies: those need physical input (force constants).
\end{remark}

\section{Groups}\label{sec:groups}

\begin{definition}
A \emph{group} is a set $G$ with a composition law $\circ$ such that
\begin{enumerate}
\item $g_1\circ g_2\in G$ for all $g_1,g_2\in G$ (closure);
\item $(g_1\circ g_2)\circ g_3=g_1\circ(g_2\circ g_3)$ (associativity);
\item there is a neutral element $e$ with $e\circ g=g\circ e=g$;
\item every $g$ has an inverse $g^{-1}$ with $g\circ g^{-1}=g^{-1}\circ g=e$.
\end{enumerate}
The number of elements $\abs{G}$ is the \emph{order} of the group.
\end{definition}
For a group of symmetry operations, $g_1\circ g_2$ means ``apply $g_2$ first, then
$g_1$''. The composition law can be anything: for the translation group it is vector
addition, $\vec{t}_1\circ\vec{t}_2=\vec{t}_1+\vec{t}_2$. Symmetry operations are
denoted as follows (Schoenflies notation):
\begin{itemize}
\item $E$: identity (do nothing). Not to be confused with the irrep $\irr{E}$.
\item $C_n$: rotation by $2\cpi/n$; $C_n^{+}$ counterclockwise, $C_n^{-}$ clockwise
      (seen from the tip of the axis).
\item $\sigma$: reflection in a plane; $\sv$ vertical (containing the main axis),
      $\sh$ horizontal, $\sd$ diagonal.
\item $I$: inversion, $\vec{r}\mapsto-\vec{r}$.
\item $S_n$: rotoreflection, a rotation $C_n$ followed by reflection in the plane
      perpendicular to the axis (\cref{sec:pointops}).
\end{itemize}

\subsection{The group \texorpdfstring{$\Cv{3}$}{C3v}}

Replacing one hydrogen of methane by deuterium leaves a triangular pyramid with the
deuterium on top. Its symmetry group $\Cv{3}$ (crystallographic symbol $3m$) has six
elements (\cref{fig:c3v}): the identity $E$, the rotations $C_3^{\pm}$ by
$\pm120^\circ$ about the vertical $z$ axis, and three vertical mirror planes
$\sigma_{\mathrm{d}1},\sigma_{\mathrm{d}2},\sigma_{\mathrm{d}3}$, each containing the
$z$ axis and one hydrogen atom. We choose the $x$ axis inside
$\sigma_{\mathrm{d}1}$.

\begin{figure}[ht]
\centering
\begin{tikzpicture}[scale=2.0,>=Stealth]
  \foreach \a/\n in {0/1,120/2,240/3}{
    \draw[dashed,gray] (\a:-1.45)--(\a:1.45);
    \node[font=\small] at (\a:1.72) {$\sigma_{\mathrm{d}\n}$};
    \fill[white,draw=black] (\a:1) circle (0.09);
    \node[font=\scriptsize] at ($(\a:1)+(\a+90:0.2)$) {H$_\n$};
  }
  \draw[thick] (0:1)--(120:1)--(240:1)--cycle;
  \fill (0,0) circle (0.06);
  \node[font=\scriptsize,below right] at (0.02,-0.02) {C, D};
  \draw[->] (0,0)--(0.7,0) node[below right,font=\small]{$x$};
  \draw[->] (0,0)--(0,0.7) node[above,font=\small]{$y$};
  \draw[->,blue!60!black] (200:0.4) arc (200:280:0.4);
  \node[blue!60!black,font=\small] at (240:0.6) {$C_3^+$};
\end{tikzpicture}
\caption{The group $\Cv{3}$ seen from above ($z$ out of the page). The deuterium
atom sits above the carbon on the $z$ axis; the three hydrogens lie below. The
mirror $\sigma_{\mathrm{d}i}$ contains the $z$ axis and hydrogen H$_i$.}
\label{fig:c3v}
\end{figure}

Its multiplication table (row $g$, column $h$, entry $g h$, i.e.\ $h$ applied first)
is
\begin{equation}\label{eq:c3vtable}
\begin{array}{c|cccccc}
 & E & C_3^+ & C_3^- & \sigma_{\mathrm{d}1} & \sigma_{\mathrm{d}2} & \sigma_{\mathrm{d}3}\\\hline
E & E & C_3^+ & C_3^- & \sigma_{\mathrm{d}1} & \sigma_{\mathrm{d}2} & \sigma_{\mathrm{d}3}\\
C_3^+ & C_3^+ & C_3^- & E & \sigma_{\mathrm{d}3} & \sigma_{\mathrm{d}1} & \sigma_{\mathrm{d}2}\\
C_3^- & C_3^- & E & C_3^+ & \sigma_{\mathrm{d}2} & \sigma_{\mathrm{d}3} & \sigma_{\mathrm{d}1}\\
\sigma_{\mathrm{d}1} & \sigma_{\mathrm{d}1} & \sigma_{\mathrm{d}2} & \sigma_{\mathrm{d}3} & E & C_3^+ & C_3^-\\
\sigma_{\mathrm{d}2} & \sigma_{\mathrm{d}2} & \sigma_{\mathrm{d}3} & \sigma_{\mathrm{d}1} & C_3^- & E & C_3^+\\
\sigma_{\mathrm{d}3} & \sigma_{\mathrm{d}3} & \sigma_{\mathrm{d}1} & \sigma_{\mathrm{d}2} & C_3^+ & C_3^- & E
\end{array}
\end{equation}
The group is \emph{not} commutative: $C_3^+\sigma_{\mathrm{d}1}=\sigma_{\mathrm{d}3}$
but $\sigma_{\mathrm{d}1}C_3^+=\sigma_{\mathrm{d}2}$. Each row and each column
contains every element exactly once (the \emph{rearrangement theorem}).

\section{Subgroups, cosets and Lagrange's theorem\beyondmark}\label{sec:cosets}

The lectures skip the structural part of group theory. Space groups, little groups and
induced representations use it constantly, so the starred sections supply it where it
first becomes useful.

A subset $H\subseteq G$ that is itself a group under the same law is a
\emph{subgroup}. For a finite subset it suffices to check closure.

\begin{definition}
For $g\in G$ the \emph{left coset} of $H$ is $gH=\{gh:h\in H\}$; the right coset is
$Hg$.
\end{definition}
Two left cosets are either identical or disjoint ($g_1H\cap g_2H\neq\emptyset$ implies
$g_2^{-1}g_1\in H$, hence $g_1H=g_2H$), and each contains $\abs{H}$ elements. So $G$ is a
disjoint union of cosets,
\[
G=g_1H\cup g_2H\cup\dots\cup g_nH,\qquad n=\abs{G}/\abs{H}.
\]
\begin{theorem}[Lagrange]
The order of a subgroup divides the order of the group. The integer
$[G:H]=\abs{G}/\abs{H}$ is the \emph{index} of $H$.
\end{theorem}
Consequences: the order of every element (the smallest $n$ with $g^n=e$) divides
$\abs{G}$; a group of prime order is cyclic. $\Cv{3}$ ($\abs{G}=6$) can only have
subgroups of order $1,2,3,6$: $\{E\}$, the three $\{E,\sigma_{\mathrm{d}i}\}$, $C_3$, and
$\Cv{3}$ itself.

\section{Representations}\label{sec:reps}

\subsection{The vector representation of \texorpdfstring{$\Cv{3}$}{C3v}}

A rotation or reflection acts linearly on vectors, so it is described by a
$3\times3$ matrix acting on the column of components $(v_x,v_y,v_z)\T$. With the axes
of \cref{fig:c3v}:
\begin{equation}\label{eq:VC3v}
\begin{gathered}
\mat{V}(E)=\one,\qquad
\mat{V}(\sigma_{\mathrm{d}1})=\begin{pmatrix}1&0&0\\0&-1&0\\0&0&1\end{pmatrix},\\[4pt]
\mat{V}(C_3^{+})=\begin{pmatrix}-\frac12&-\frac{\sqrt3}{2}&0\\ \frac{\sqrt3}{2}&-\frac12&0\\0&0&1\end{pmatrix},\qquad
\mat{V}(C_3^{-})=\begin{pmatrix}-\frac12&\frac{\sqrt3}{2}&0\\ -\frac{\sqrt3}{2}&-\frac12&0\\0&0&1\end{pmatrix},\\[4pt]
\mat{V}(\sigma_{\mathrm{d}2})=\begin{pmatrix}-\frac12&-\frac{\sqrt3}{2}&0\\ -\frac{\sqrt3}{2}&\frac12&0\\0&0&1\end{pmatrix},\qquad
\mat{V}(\sigma_{\mathrm{d}3})=\begin{pmatrix}-\frac12&\frac{\sqrt3}{2}&0\\ \frac{\sqrt3}{2}&\frac12&0\\0&0&1\end{pmatrix}.
\end{gathered}
\end{equation}
The mirror $\sigma_{\mathrm{d}1}$ contains the $x$ and $z$ axes, so it only reverses
$v_y$. How to construct the other matrices without memorising formulas is explained
in \cref{sec:active}. The matrices reproduce the multiplication table; for example
$\mat{V}(C_3^+)\mat{V}(C_3^+)=\mat{V}(C_3^-)$, because two rotations by $120^\circ$
make a rotation by $240^\circ=-120^\circ$.

\subsection{Definition}

\begin{definition}
A (linear) \emph{representation} $\Gamma$ of a group $G$ on a vector space $L$
assigns to each $g\in G$ an invertible linear operator $\mat{D}(g)$ on $L$ such that
\begin{equation}\label{eq:homom}
  \mat{D}(g_1)\,\mat{D}(g_2)=\mat{D}(g_1 g_2)\qquad\text{for all } g_1,g_2\in G.
\end{equation}
The dimension of $L$ is the \emph{dimension} of the representation. Choosing a basis
of $L$, each $\mat{D}(g)$ becomes a matrix.
\end{definition}
Whatever the composition law of the group (vector addition for translations), the
operators of a representation are always composed by \emph{matrix multiplication}.
From \eqref{eq:homom} it follows that $\mat{D}(e)=\one$ and
$\mat{D}(g^{-1})=\mat{D}(g)^{-1}$.

\begin{itemize}
\item The assignment $\mat{D}(g)=1$ (a $1\times1$ matrix) for every $g$ satisfies
      \eqref{eq:homom}. It is the \emph{identity} or \emph{trivial} representation,
      called $\irr{A}_1$ in $\Cv{3}$ and $\Td$. It is not useless: it describes
      invariant quantities.
\item A representation is \emph{faithful} if different elements get different
      matrices. The vector representation of $\Cv{3}$ is faithful; the trivial one is
      not. Faithfulness is not required.
\item There are infinitely many representations (any dimension, any basis). The
      central task of the theory is to classify them and to reduce any given one to
      a few standard pieces.
\end{itemize}

\section{Reducible and irreducible representations}\label{sec:reducible}

All six matrices \eqref{eq:VC3v} have the same block-diagonal structure: a $2\times2$
block acting on $(v_x,v_y)$ and a $1\times1$ block acting on $v_z$. Because matrix
multiplication preserves this structure, each set of blocks is itself a
representation. The $1\times1$ blocks are all equal to $1$: that is $\irr{A}_1$. The
$2\times2$ blocks form a two-dimensional representation called $\irr{E}$. We write
\begin{equation}
  \GV=\irr{A}_1\oplus\irr{E}.
\end{equation}
The symbol $\oplus$ (\emph{direct sum}) does not add matrices of different sizes; it
means that in a suitable basis all matrices take block-diagonal form with the given
blocks.

Geometrically: the $z$ axis and the $xy$ plane are \emph{invariant subspaces}. No
operation of the group takes a vector on the $z$ axis off the axis, nor a vector in
the plane out of the plane.

\begin{definition}
A representation on $L$ is \emph{reducible} if $L$ has a proper invariant subspace
($\neq\{0\},L$) and \emph{irreducible} otherwise.
\end{definition}
The representation $\irr{E}$ is irreducible: every line in the $xy$ plane is moved by
some $C_3$ rotation or some mirror.

For finite groups every representation is equivalent to a unitary one
(\cref{sec:unitarity}), and for unitary representations the orthogonal complement of an
invariant subspace is also invariant. Hence a reducible representation always
splits completely into a direct sum of irreducible ones. This is why irreps are
\emph{building blocks}.

\section{Equivalent representations}\label{sec:equiv}

The block structure depends on the basis. Rotate the $y$ and $z$ axes by $45^\circ$
about $x$, i.e.\ change basis with
\[
\mat{A}=\begin{pmatrix}1&0&0\\0&\tfrac{1}{\sqrt2}&-\tfrac{1}{\sqrt2}\\0&\tfrac{1}{\sqrt2}&\tfrac{1}{\sqrt2}\end{pmatrix},
\qquad \mat{W}(g)=\mat{A}\,\mat{V}(g)\,\mat{A}^{-1}.
\]
Then, for instance,
\[
\mat{W}(C_3^+)=\begin{pmatrix}-\frac12&-\frac{\sqrt6}{4}&-\frac{\sqrt6}{4}\\[2pt]
\frac{\sqrt6}{4}&\frac14&-\frac34\\[2pt]\frac{\sqrt6}{4}&-\frac34&\frac14\end{pmatrix},
\qquad
\mat{W}(\sigma_{\mathrm{d}1})=\begin{pmatrix}1&0&0\\0&0&-1\\0&-1&0\end{pmatrix},
\]
and no block structure is visible, although $\mat{W}$ describes exactly the same
physics in rotated axes.

\begin{definition}
Two representations $\mat{D}_1$ and $\mat{D}_2$ of $G$ are \emph{equivalent} if there
is an invertible matrix $\mat{A}$ such that
$\mat{D}_2(g)=\mat{A}\,\mat{D}_1(g)\,\mat{A}^{-1}$ for all $g\in G$.
\end{definition}
Equivalent representations are the same representation written in different bases.
This raises two practical questions: given a set of matrices, how do we know whether
it is reducible, and how do we find its irreducible components, without trying
infinitely many changes of basis? The answer is the theory of characters
(\cref{ch:building}). It also shows that a finite group has only finitely many
inequivalent irreps, as many as it has conjugacy classes: $\Cv{3}$ has three and
$\Td$ has five.

\section{Unitarity and complete reducibility\beyondmark}\label{sec:unitarity}

The lectures state the main theorems of representation theory without proof. The
starred sections supply the proofs, starting here with the two facts used in
\cref{sec:reducible}. $G$ is a finite group throughout; sums over $g$ run over all
elements.

\begin{theorem}
Every representation of a finite group is equivalent to a unitary one.
\end{theorem}
\begin{proof}
Given $\mat{D}(g)$ on $\mathbb{C}^n$ with the standard inner product $(x,y)=x\herm y$,
define the averaged inner product
\[\{x,y\}=\sum_g\bigl(\mat{D}(g)x,\mat{D}(g)y\bigr).\] It is positive
definite and invariant: $\{\mat{D}(h)x,\mat{D}(h)y\}=\sum_g(\mat{D}(gh)x,\mat{D}(gh)y)=\{x,y\}$
by the rearrangement theorem. In a basis orthonormal for $\{\cdot,\cdot\}$ the matrices
are unitary.
\end{proof}

\begin{theorem}[Maschke]
If $W$ is an invariant subspace of a unitary representation, so is its orthogonal
complement $W^\perp$. Hence every representation is a direct sum of irreps.
\end{theorem}
\begin{proof}
If $y\in W^\perp$ and $x\in W$, then $(x,\mat{D}(g)y)=(\mat{D}(g)^{-1}x,y)=0$, because
$\mat{D}(g)^{-1}x=\mat{D}(g^{-1})x\in W$. Iterate on $W$ and $W^\perp$ until the pieces are
irreducible (the dimension decreases at each step).
\end{proof}

\exercisesection

\begin{exercise}\label{ex:1:table}
Verify three entries of the multiplication table \eqref{eq:c3vtable} by multiplying
the matrices \eqref{eq:VC3v}: $C_3^+\sigma_{\mathrm{d}1}$,
$\sigma_{\mathrm{d}1}C_3^+$ and $\sigma_{\mathrm{d}2}\sigma_{\mathrm{d}1}$. Interpret
the last one geometrically.
\end{exercise}

\begin{solution}
With the matrices \eqref{eq:VC3v} (only the upper-left $2\times2$ blocks matter; the
$zz$ entry is always $1$):
\[
\mat{V}(C_3^+)\mat{V}(\sigma_{\mathrm{d}1})=
\begin{pmatrix}-\frac12&-\frac{\sqrt3}{2}\\ \frac{\sqrt3}{2}&-\frac12\end{pmatrix}
\begin{pmatrix}1&0\\0&-1\end{pmatrix}=
\begin{pmatrix}-\frac12&\frac{\sqrt3}{2}\\ \frac{\sqrt3}{2}&\frac12\end{pmatrix}
=\mat{V}(\sigma_{\mathrm{d}3}),
\]
\[
\mat{V}(\sigma_{\mathrm{d}1})\mat{V}(C_3^+)=
\begin{pmatrix}-\frac12&-\frac{\sqrt3}{2}\\ -\frac{\sqrt3}{2}&\frac12\end{pmatrix}
=\mat{V}(\sigma_{\mathrm{d}2}),
\]
\[
\mat{V}(\sigma_{\mathrm{d}2})\mat{V}(\sigma_{\mathrm{d}1})=
\begin{pmatrix}-\frac12&\frac{\sqrt3}{2}\\ -\frac{\sqrt3}{2}&-\frac12\end{pmatrix}
=\mat{V}(C_3^-).
\]
Geometrically, reflecting first in a line and then in a second line making the angle
$\alpha$ with it is a rotation by $2\alpha$. From $\sigma_{\mathrm{d}1}$ ($0^\circ$) to
$\sigma_{\mathrm{d}2}$ ($120^\circ$) the angle is $120^\circ$, so the product is a rotation
by $240^\circ=-120^\circ$, i.e.\ $C_3^-$.
\end{solution}

\begin{exercise}\label{ex:1:subgroups}
List all subgroups of $\Cv{3}$. Which of them are commutative?
\end{exercise}

\begin{solution}
By Lagrange's theorem the orders can be $1,2,3,6$: $\{E\}$; the three groups
$\{E,\sigma_{\mathrm{d}i}\}$; $C_3=\{E,C_3^+,C_3^-\}$; and $\Cv{3}$. All proper subgroups have
prime or unit order, hence are cyclic and commutative; $\Cv{3}$ itself is not
commutative.
\end{solution}

\begin{exercise}\label{ex:1:A2}
Show that assigning $+1$ to $E,C_3^\pm$ and $-1$ to the three mirrors is a
representation of $\Cv{3}$ (called $\irr{A}_2$). Is it faithful? Give a physical
quantity that transforms according to $\irr{A}_2$.
\end{exercise}

\begin{solution}
The assignment is $g\mapsto\det\mat{V}(g)$, and $\det$ is multiplicative:
$\det(\mat{V}(g_1)\mat{V}(g_2))=\det\mat{V}(g_1)\det\mat{V}(g_2)$. (Equivalently: rotation
times rotation and mirror times mirror are rotations, rotation times mirror is a mirror,
as the table \eqref{eq:c3vtable} confirms.) It is not faithful: its kernel is $C_3$.
Quantities transforming as $\irr{A}_2$: the $z$ component of an axial vector ($R_z$,
$L_z$, $B_z$), which is invariant under the rotations about $z$ and reversed by vertical
mirrors.
\end{solution}

\begin{exercise}\label{ex:1:count}
A planar molecule $\mathrm{BF_3}$ has four atoms. How many vibrational degrees of
freedom does it have? What about the linear molecule $\mathrm{CO_2}$? (For a linear
molecule, rotation about the molecular axis does not move the atoms.)
\end{exercise}

\begin{solution}
$\mathrm{BF_3}$: $3\times4=12$ degrees of freedom, minus 3 translations and 3 rotations:
$6$ vibrational. $\mathrm{CO_2}$ (linear): $9-3-2=4$.
\end{solution}

\begin{exercise}\label{ex:1:equiv}
Show that the matrices $\mat{W}(g)$ of \cref{sec:equiv} have the same traces as the
matrices $\mat{V}(g)$. Then prove that $\tr(\mat{A}\mat{B}\mat{A}^{-1})=\tr\mat{B}$
for any invertible $\mat{A}$.
\end{exercise}

\begin{solution}
$\tr\mat{W}(C_3^+)=-\tfrac12+\tfrac14+\tfrac14=0=\tr\mat{V}(C_3^+)$ and
$\tr\mat{W}(\sigma_{\mathrm{d}1})=1+0+0=1=\tr\mat{V}(\sigma_{\mathrm{d}1})$. In general
$\tr(\mat{A}\mat{B}\mat{A}^{-1})=\tr(\mat{A}^{-1}\mat{A}\mat{B})=\tr\mat{B}$ by the cyclic
property $\tr(\mat{X}\mat{Y})=\sum_{ij}X_{ij}Y_{ji}=\tr(\mat{Y}\mat{X})$.
\end{solution}

\chapter{Irreducible representations as building blocks}\label{ch:building}

\begin{watch}{\lec{2}}
Recap and the problem of finding invariant subspaces \ts{2}{498}{8:18} $\cdot$
traces and characters \ts{2}{1129}{18:49} $\cdot$
classes \ts{2}{1447}{24:07} $\cdot$
number of irreps and $\sum d_\mu^2=\abs{G}$ \ts{2}{2017}{33:37} $\cdot$
orthogonality of characters \ts{2}{2204}{36:44} $\cdot$
the ``magic formula'' \ts{2}{2461}{41:01} $\cdot$
computing characters without matrices \ts{2}{3441}{57:21} $\cdot$
rotoreflections and rotoinversions \ts{2}{3691}{61:31} $\cdot$
axial vectors \ts{2}{4330}{72:10} $\cdot$
the mechanical representation \ts{2}{4453}{74:13} $\cdot$
$\mathrm{CH_3D}$ worked out \ts{2}{5215}{86:55}
\end{watch}

The aim of this chapter is to derive, with nothing but a character table and a few
multiplications, the decomposition $\Gvib(\mathrm{CH_3D})=3\,\irr{A}_1\oplus3\,\irr{E}$.

\section{Characters}

Matrices change under a change of basis, but traces do not:
$\tr(\mat{A}\mat{D}\mat{A}^{-1})=\tr\mat{D}$. Hence equivalent representations have
the same traces. The converse is far from obvious, and is one of the central
results of the theory (proved in \cref{sec:proofs}):

\begin{theorem}\label{thm:charequiv}
Two representations of a finite group are equivalent if and only if
$\tr\mat{D}_1(g)=\tr\mat{D}_2(g)$ for \emph{every} $g\in G$.
\end{theorem}

\begin{definition}
The \emph{character} of a representation $\Gamma$ is the function
$\chi_\Gamma\colon G\to\mathbb{C}$, $\chi_\Gamma(g)=\tr\mat{D}(g)$.
\end{definition}
A character is the whole list of traces, one per group element; the trace of a
single element is not enough to characterise a representation. Two consequences:
\begin{itemize}
\item $\chi(E)=\dim\Gamma$, because $\mat{D}(E)$ is the unit matrix.
\item If $\Gamma=\Gamma_1\oplus\Gamma_2$ then $\chi_\Gamma=\chi_{\Gamma_1}+\chi_{\Gamma_2}$:
      the trace of a block-diagonal matrix is the sum of the traces of the blocks.
\end{itemize}
For the vector representation of $\Cv{3}$ \eqref{eq:VC3v}, the traces of the full
matrices, of the $2\times2$ blocks and of the $1\times1$ blocks are
\[
\begin{array}{c|cccccc}
 & E & C_3^+ & C_3^- & \sigma_{\mathrm{d}1} & \sigma_{\mathrm{d}2} & \sigma_{\mathrm{d}3}\\\hline
\chiV & 3 & 0 & 0 & 1 & 1 & 1\\
\chi_{\irr{E}} & 2 & -1 & -1 & 0 & 0 & 0\\
\chi_{\irr{A}_1} & 1 & 1 & 1 & 1 & 1 & 1
\end{array}
\]
and indeed $\chiV=\chi_{\irr{A}_1}+\chi_{\irr{E}}$.

\section{Conjugacy classes}

In the table above $C_3^+$ and $C_3^-$ have equal traces, and so do the three mirrors.
This is no accident.

\begin{definition}
Two elements $a,b\in G$ are \emph{conjugate} if $b=g a g^{-1}$ for some $g\in G$.
Conjugacy is an equivalence relation; its equivalence classes are the
\emph{(conjugacy) classes} of $G$.
\end{definition}
Since $\tr\mat{D}(gag^{-1})=\tr(\mat{D}(g)\mat{D}(a)\mat{D}(g)^{-1})=\tr\mat{D}(a)$,
\emph{characters are class functions}: constant on each class.

Intuitively, a class collects the elements ``of the same kind'' that are related by
some symmetry of the group. In $\Cv{3}$:
\begin{itemize}
\item $\{E\}$: the identity is always a class by itself.
\item $\{C_3^+,C_3^-\}$: a vertical mirror turns a counterclockwise rotation into a
      clockwise one (a rotation seen in a mirror parallel to its axis reverses its
      sense), $\sigma_{\mathrm{d}1}C_3^+\sigma_{\mathrm{d}1}^{-1}=C_3^-$.
\item $\{\sigma_{\mathrm{d}1},\sigma_{\mathrm{d}2},\sigma_{\mathrm{d}3}\}$: a $C_3$
      rotation carries each mirror plane into another,
      $C_3^+\sigma_{\mathrm{d}1}C_3^-=\sigma_{\mathrm{d}2}$.
\end{itemize}
Being ``of the same kind'' is not enough: there must be a group element relating
them. In a group containing only rotations about $z$ (such as $C_3$), $C_3^+$ and
$C_3^-$ are \emph{not} conjugate, since nothing reverses the sense of rotation. In
a commutative (Abelian) group every element is a class by itself.

\subsection{Class sizes and rules of thumb\beyondmark}
Conjugation $a\mapsto gag^{-1}$ is the change of ``point of view'' by the symmetry $g$.
The number of elements in the class of $a$ is $\abs{G}/\abs{Z(a)}$, where
$Z(a)=\{g:ga=ag\}$ is the centraliser; in particular class sizes divide $\abs{G}$.
Rules of thumb for point groups:
\begin{itemize}
\item rotations by the same angle about axes related by a group operation are
      conjugate;
\item $C_n^{+}$ and $C_n^{-}$ about the same axis are conjugate if and only if the group
      contains a $C_2$ perpendicular to the axis or a mirror containing the axis;
\item mirrors related by a group operation are conjugate.
\end{itemize}

\section{Normal subgroups, quotient groups and homomorphisms\beyondmark}\label{sec:normal}

\begin{definition}
A subgroup $N$ is \emph{normal} ($N\trianglelefteq G$) if $gNg^{-1}=N$ for all $g$,
i.e.\ if it is a union of whole classes. Equivalently, left and right cosets coincide,
$gN=Ng$.
\end{definition}
For a normal subgroup the product of cosets (\cref{sec:cosets}) is well defined,
$(g_1N)(g_2N)=g_1g_2N$, and the cosets form the \emph{quotient group} $G/N$ of order
$\abs{G}/\abs{N}$.

\begin{example}
$C_3=\{E,C_3^+,C_3^-\}$ is normal in $\Cv{3}$ (it is the union of the classes $\{E\}$
and $\{C_3^\pm\}$). The quotient $\Cv{3}/C_3$ has two elements, ``rotations'' and
``reflections''. It is isomorphic to $\{+1,-1\}$ under multiplication, which is
precisely the irrep $\irr{A}_2$ viewed as a homomorphism.
\end{example}
The most important example in this course is the translation subgroup of a space
group, whose quotient is the point group (\cref{sec:semidirect}).

\subsection{Homomorphisms}
A map $f\colon G\to G'$ with $f(g_1g_2)=f(g_1)f(g_2)$ is a \emph{homomorphism}. Its
kernel $\Ker f=\{g:f(g)=e'\}$ is a normal subgroup, and
\begin{theorem}[First isomorphism theorem]
$G/\Ker f\cong f(G)$.
\end{theorem}
A representation is a homomorphism $G\to\GL{n,\mathbb{C}}$. It is faithful exactly when
its kernel is $\{e\}$. A non-faithful representation of $G$ is a faithful one of
$G/\Ker$: the irrep $\irr{A}_2$ of $\Cv{3}$ ``lives'' on $\Cv{3}/C_3$. The map
$\SU{2}\to\SO{3}$ is a homomorphism with kernel $\{\one,-\one\}$, so
$\SO{3}\cong\SU{2}/\{\pm\one\}$ (\cref{sec:su2}).

\section{Character tables}

Because characters are class functions, they are tabulated by class, with the number
of elements of each class written in front of its symbol:
\begin{equation}\label{eq:C3vchar}
\begin{array}{c|ccc|l}
\Cv{3} & E & 2C_3 & 3\sd & \text{basis functions}\\\hline
\irr{A}_1 & 1 & 1 & 1 & z,\ x^2+y^2,\ z^2\\
\irr{A}_2 & 1 & 1 & -1 & R_z\\
\irr{E} & 2 & -1 & 0 & (x,y),\ (R_x,R_y),\ (x^2-y^2,\,xy),\ (xz,\,yz)
\end{array}
\end{equation}
The last column lists functions or quantities that transform according to each irrep;
$R_i$ denote components of an axial vector.

\begin{theorem}\label{thm:classcount}
For a finite group $G$:
\begin{enumerate}
\item the number of inequivalent irreps equals the number of conjugacy classes;
\item $\sum_\mu d_\mu^2=\abs{G}$ (Burnside).
\end{enumerate}
\end{theorem}
For $\Cv{3}$: three classes, hence three irreps, with $d_1^2+d_2^2+d_3^2=6$. The only
solution in positive integers is $1,1,2$. Infinitely many representations of any
dimension, and all of them are built from just these three blocks. For $\Td$
($24$ elements, five classes): $1+1+4+9+9=24$, dimensions $1,1,2,3,3$.

You never need to construct character tables yourself: they are tabulated in
Bradley and Cracknell and in the \texttt{POINT} program of the Bilbao
Crystallographic Server. The table of $\Td$ is
\begin{equation}\label{eq:Tdchar}
\begin{array}{c|ccccc|l}
\Td & E & 8C_3 & 3C_2 & 6S_4 & 6\sd & \text{basis functions}\\\hline
\irr{A}_1 & 1 & 1 & 1 & 1 & 1 & x^2+y^2+z^2\\
\irr{A}_2 & 1 & 1 & 1 & -1 & -1 & \\
\irr{E} & 2 & -1 & 2 & 0 & 0 & (2z^2-x^2-y^2,\ x^2-y^2)\\
\irr{T}_1 & 3 & 0 & -1 & 1 & -1 & (R_x,R_y,R_z)\\
\irr{T}_2 & 3 & 0 & -1 & -1 & 1 & (x,y,z),\ (yz,xz,xy)
\end{array}
\end{equation}

\section{Orthogonality and the reduction formula}

Characters of a group with $\abs{G}$ elements can be regarded as vectors with
$\abs{G}$ components, and we give them the natural inner product
\begin{equation}\label{eq:inner}
  \inner{\chi}{\chi'}=\frac{1}{\abs{G}}\sum_{g\in G}\chi(g)^{*}\,\chi'(g).
\end{equation}

\begin{theorem}[Orthogonality of characters]\label{thm:orth}
If $\chi^{(\mu)}$ and $\chi^{(\nu)}$ are characters of irreps, then
$\inner{\chi^{(\mu)}}{\chi^{(\nu)}}=\delta_{\mu\nu}$.
\end{theorem}
The irreducible characters are therefore an orthonormal set, and since their number
equals the number of classes they are a basis of the space of class functions. If
$\Gamma=\bigoplus_\nu m_\nu\Gamma^{(\nu)}$, then $\chi_\Gamma=\sum_\nu m_\nu\chi^{(\nu)}$,
and projecting on $\chi^{(\mu)}$ gives the multiplicities:

\begin{keyresult}[Reduction formula (the ``magic formula'')]
\begin{equation}\label{eq:magic}
  m_\mu=\inner{\chi^{(\mu)}}{\chi_\Gamma}
      =\frac{1}{\abs{G}}\sum_{g\in G}\chi^{(\mu)}(g)^{*}\chi_\Gamma(g)
      =\frac{1}{\abs{G}}\sum_{\text{classes }\mathcal{C}}N_{\mathcal{C}}\,
        \chi^{(\mu)}(\mathcal{C})^{*}\chi_\Gamma(\mathcal{C}),
\end{equation}
where $N_{\mathcal{C}}$ is the number of elements of class $\mathcal{C}$. A useful
corollary: $\Gamma$ is irreducible if and only if $\inner{\chi_\Gamma}{\chi_\Gamma}=1$;
in general $\inner{\chi_\Gamma}{\chi_\Gamma}=\sum_\mu m_\mu^2$.
\end{keyresult}

The multiplicities come out as non-negative integers: $m_\mu$ counts diagonal blocks,
and there is no such thing as half a block or a negative number of blocks. A
fractional result is a sure sign of an arithmetic mistake (or, as in
\cref{sec:doublevalued}, of using the wrong group).

\begin{remark}[The most common mistake]
The sum in \eqref{eq:magic} runs over \emph{elements}, not classes. When using a
table organised by classes, multiply each term by the class size $N_{\mathcal{C}}$.
A practical trick: write the class sizes above the columns before starting.
\end{remark}

\begin{example}[Decomposing $\GV$ of $\Cv{3}$]
With $\chiV=(3,0,1)$ on the classes $(E,2C_3,3\sd)$ and the table \eqref{eq:C3vchar}:
\begin{align*}
m_{\irr{A}_1}&=\tfrac16\,(1\cdot3\cdot1+2\cdot0\cdot1+3\cdot1\cdot1)=1,\\
m_{\irr{A}_2}&=\tfrac16\,(1\cdot3\cdot1+2\cdot0\cdot1+3\cdot1\cdot(-1))=0,\\
m_{\irr{E}}&=\tfrac16\,(1\cdot3\cdot2+2\cdot0\cdot(-1)+3\cdot1\cdot0)=1.
\end{align*}
So $\GV=\irr{A}_1\oplus\irr{E}$, without ever looking for a good basis. Checks:
$\inner{\chiV}{\chiV}=\frac16(9+0+3)=2=1^2+1^2$.
\end{example}

\section{Schur's lemma and the orthogonality theorems\beyondmark}\label{sec:proofs}

The lecture states \cref{thm:charequiv,thm:classcount,thm:orth} without proof. All of
them rest on Schur's lemma. By \cref{sec:unitarity} we may take every representation
unitary.

\subsection{Schur's lemma}\label{sec:schur}

\begin{lemma}[Schur I]
Let $\mat{D}^{(1)}$ ($d_1$-dimensional) and $\mat{D}^{(2)}$ ($d_2$-dimensional) be irreps and
let $\mat{A}$ be a $d_1\times d_2$ matrix with
$\mat{D}^{(1)}(g)\mat{A}=\mat{A}\mat{D}^{(2)}(g)$ for all $g$ (an \emph{intertwiner}).
Then either $\mat{A}=0$, or $\mat{A}$ is invertible and the irreps are equivalent.
\end{lemma}
\begin{proof}
$\Ker\mat{A}$ is invariant under $\mat{D}^{(2)}$ (if $\mat{A}x=0$ then
$\mat{A}\mat{D}^{(2)}x=\mat{D}^{(1)}\mat{A}x=0$). By irreducibility it is $\{0\}$ or everything.
Likewise the image of $\mat{A}$ is invariant under $\mat{D}^{(1)}$, so it is $\{0\}$ or
everything. Either $\mat{A}=0$, or $\mat{A}$ is injective and surjective.
\end{proof}

\begin{lemma}[Schur II]
A matrix commuting with all matrices of an irrep (over $\mathbb{C}$) is a multiple of
the identity.
\end{lemma}
\begin{proof}
Let $\lambda$ be an eigenvalue of $\mat{A}$. Then $\mat{A}-\lambda\one$ commutes with the
irrep and is not invertible, so by Schur I it vanishes.
\end{proof}
Schur II is the heart of Wigner's theorem: an invariant Hamiltonian restricted to one
copy of an irrep is proportional to the identity, so all its $d_\mu$ partner states have
the same energy.

\subsection{The great orthogonality theorem}

\begin{theorem}\label{thm:GOT}
For unitary irreps $\mat{D}^{(\mu)}$, $\mat{D}^{(\nu)}$ (chosen identical if equivalent),
\begin{equation}\label{eq:GOT}
\sum_{g\in G}D^{(\mu)}_{ij}(g)^{*}\,D^{(\nu)}_{kl}(g)=\frac{\abs{G}}{d_\mu}\,
\delta_{\mu\nu}\,\delta_{ik}\,\delta_{jl}.
\end{equation}
\end{theorem}
\begin{proof}
For any $d_\nu\times d_\mu$ matrix $\mat{X}$, the average
$\mat{A}=\sum_g\mat{D}^{(\nu)}(g)\,\mat{X}\,\mat{D}^{(\mu)}(g)^{-1}$ is an intertwiner:
$\mat{D}^{(\nu)}(h)\mat{A}=\mat{A}\mat{D}^{(\mu)}(h)$ by rearrangement. For $\mu\neq\nu$,
$\mat{A}=0$ (Schur I). For $\mu=\nu$, $\mat{A}=\lambda\one$ (Schur II), with $\lambda$
fixed by the trace: $d_\mu\lambda=\tr\mat{A}=\abs{G}\tr\mat{X}$. Now choose $\mat{X}$ with a
single entry $1$ in position $(l,j)$. Using unitarity,
$[\mat{D}^{(\mu)}(g)^{-1}]_{ji}=D^{(\mu)}_{ij}(g)^{*}$, the $(k,i)$ element of $\mat{A}$ is
$\sum_gD^{(\nu)}_{kl}(g)D^{(\mu)}_{ij}(g)^{*}$. It vanishes for $\mu\neq\nu$ and equals
$\lambda\delta_{ki}=(\abs{G}/d_\mu)\,\delta_{lj}\delta_{ki}$ for $\mu=\nu$, which is
\eqref{eq:GOT}.
\end{proof}
Interpretation: the $\sum_\mu d_\mu^2$ functions $g\mapsto\sqrt{d_\mu/\abs{G}}\,D^{(\mu)}_{ij}(g)$
are orthonormal in the $\abs{G}$-dimensional space of functions on $G$. Hence
$\sum_\mu d_\mu^2\le\abs{G}$.

\subsection{Proofs of the character theorems}

Setting $i=j$, $k=l$ in \eqref{eq:GOT} and summing gives the orthogonality of
characters, \cref{thm:orth}:
$\frac{1}{\abs{G}}\sum_g\chi^{(\mu)}(g)^{*}\chi^{(\nu)}(g)=\delta_{\mu\nu}$.
From it follow the reduction formula \eqref{eq:magic} and the irreducibility test
$\inner{\chi}{\chi}=1$.

\begin{proof}[Proof of \cref{thm:charequiv}]
By complete reducibility every representation is $\bigoplus m_\mu\Gamma^{(\mu)}$, and
the multiplicities are determined by the character through \eqref{eq:magic}. Two
representations with the same character have the same multiplicities, hence are both
equivalent to the same direct sum.
\end{proof}

\paragraph{The regular representation.} Let $G$ act on the $\abs{G}$-dimensional space
with basis $\{e_h\}_{h\in G}$ by $g\,e_h=e_{gh}$. Since $gh\neq h$ unless $g=e$, its
character is $\chi_{\mathrm{reg}}(e)=\abs{G}$ and $\chi_{\mathrm{reg}}(g)=0$ otherwise. By
\eqref{eq:magic}, $m_\mu=\frac{1}{\abs{G}}\,\abs{G}\,d_\mu=d_\mu$: \emph{every irrep occurs
in the regular representation as many times as its dimension}. Counting dimensions,
$\sum_\mu d_\mu^2=\abs{G}$ (Burnside), which proves part 2 of \cref{thm:classcount}.

\paragraph{Number of irreps.} The characters are class functions, orthonormal, so their
number is at most the number of classes $n_{\mathrm{c}}$. Completeness of the matrix
elements (they span all functions on $G$, by Burnside) implies that the characters span
all class functions, so the number of irreps is exactly $n_{\mathrm{c}}$. An equivalent
statement is the \emph{second orthogonality relation},
\begin{equation}
\sum_\mu\chi^{(\mu)}(\mathcal{C})^{*}\chi^{(\mu)}(\mathcal{C}')
=\frac{\abs{G}}{N_{\mathcal{C}}}\,\delta_{\mathcal{C}\mathcal{C}'},
\end{equation}
orthogonality of the \emph{columns} of a character table. It is a quick check of a
table: for $\Cv{3}$ the column of $\sd$ gives $1+1+0=2=6/3$.

\section{Computing characters without matrices}\label{sec:pointops}

To use the reduction formula we still need the character of the representation of
interest. Two observations make this almost free:
\begin{enumerate}
\item Characters are class functions: compute one element per class, choosing the
      easiest.
\item Traces do not depend on the basis: for each element, use the most convenient
      axes.
\end{enumerate}

\subsection{The operations of a point group}
Besides the identity, a point group contains four kinds of operations:
\begin{itemize}
\item proper rotations $C_n^{\pm}$ by $\theta=\pm2\cpi/n$;
\item reflections $\sigma$;
\item the inversion $I$;
\item \emph{rotoreflections} $S_n^{\pm}=\sh C_n^{\pm}$: rotate by $\pm2\cpi/n$, then
      reflect in the plane perpendicular to the axis.
\end{itemize}
Crystallographers (and the Bilbao server) use \emph{rotoinversions} $\bar{n}$ instead:
rotate by $2\cpi/n$, then invert. The two are related through $\sh=I C_2$
(reflection in a plane equals inversion times a half-turn about the normal):
\[
S_n^{\pm}=\sh C_n^{\pm}=I\,C_2\,C_n^{\pm}=I\,C(\cpi\pm2\cpi/n).
\]
For example $S_3^{+}=I\,C_6^{-}=\bar6^{-}$ and $S_4^{+}=\bar4^{-}$: a rotoreflection
corresponds to a rotoinversion about the same axis but with a different angle and the
\emph{opposite} sense. Tables in Schoenflies notation (Bradley--Cracknell) use $S_n$;
tables in international notation use $\bar n$. Check which convention a table uses.

Every improper operation (determinant $-1$) is the inversion times a proper rotation.

\subsection{Character of the vector representation}
Take the rotation axis as $z$. A rotation by $\theta$ has diagonal elements
$\cos\theta,\cos\theta,1$; a rotoreflection $S(\theta)=\sh C(\theta)$ has
$\cos\theta,\cos\theta,-1$. Hence
\begin{equation}\label{eq:chiV}
\boxed{\;\chiV\bigl(C(\theta)\bigr)=1+2\cos\theta,\qquad
\chiV\bigl(S(\theta)\bigr)=-1+2\cos\theta\;}
\end{equation}
with the special cases $\chiV(E)=3$, $\chiV(C_2)=-1$, $\chiV(C_3)=0$, $\chiV(C_4)=1$,
$\chiV(C_6)=2$, $\chiV(\sigma)=\chiV(S(0))=1$, $\chiV(I)=\chiV(S(\cpi))=-3$,
$\chiV(S_4)=-1$, $\chiV(S_6)=0$, $\chiV(S_3)=-2$.

\subsection{Axial vectors}
An axial vector (pseudovector) transforms like a vector under proper rotations but
does not change sign under inversion, $\vec{L}=\vec{r}\times\vec{p}\mapsto
(-\vec{r})\times(-\vec{p})=\vec{L}$. Since every improper operation is $I$ times a
rotation, the axial-vector matrices are $\mat{R}(g)=\det\bigl(\mat{V}(g)\bigr)\,\mat{V}(g)$
and
\begin{equation}\label{eq:chiR}
\chiR(g)=\det\bigl(\mat{V}(g)\bigr)\,\chiV(g).
\end{equation}
A counter-intuitive consequence: a mirror \emph{reverses} the components of a
pseudovector parallel to the mirror and \emph{preserves} the perpendicular one,
exactly the opposite of an ordinary vector. Think of a current loop and its mirror
image: a magnetic field lying in the mirror plane flips, while the component normal
to the mirror does not.

For $\Cv{3}$: $\chiR=(3,0,-1)$, so $\GR=\irr{A}_2\oplus\irr{E}$ ($R_z$ spans $\irr{A}_2$,
which explains the entry in \eqref{eq:C3vchar}).

\section{The mechanical representation}\label{sec:mechanical}

Label the atoms of a molecule $a=1,\dots,N$ and collect their displacements in the
$3N$-component vector $\mat{u}=(\vec{u}_1,\dots,\vec{u}_N)$. A symmetry operation $g$
does two things to a distortion:
\begin{enumerate}
\item it \emph{permutes} atoms of the same species: atom $a$ is carried to the position
      of atom $g(a)$;
\item it \emph{rotates} each displacement vector by $\mat{V}(g)$.
\end{enumerate}
Hence $\vec{u}'_{g(a)}=\mat{V}(g)\,\vec{u}_a$, and the $3N\times3N$ matrix $\mat{M}(g)$
consists of $3\times3$ blocks equal to $\mat{V}(g)$, placed in block-row $g(a)$ and
block-column $a$. For $\mathrm{CH_3D}$ with atoms ordered
$(\mathrm{H}_1,\mathrm{H}_2,\mathrm{H}_3,\mathrm{D},\mathrm{C})$, the rotation $C_3^+$
sends $\mathrm{H}_1\to\mathrm{H}_2\to\mathrm{H}_3\to\mathrm{H}_1$ and fixes D and C:
\[
\mat{M}(C_3^+)=\begin{pmatrix}
0&0&\mat{V}&0&0\\
\mat{V}&0&0&0&0\\
0&\mat{V}&0&0&0\\
0&0&0&\mat{V}&0\\
0&0&0&0&\mat{V}
\end{pmatrix},\qquad \mat{V}\equiv\mat{V}(C_3^+).
\]
Only blocks on the diagonal contribute to the trace, and a diagonal block appears
exactly when an atom is left in place. Therefore
\begin{keyresult}
\begin{equation}\label{eq:chiM}
  \chiM(g)=\Ninv(g)\,\chiV(g),
\end{equation}
where $\Ninv(g)$ is the number of atoms left invariant by $g$.
\end{keyresult}
There is no need to construct any $15\times15$ matrix: look at the molecule and count.

\subsection{Removing translations and rotations}
Uniform translations $\vec{u}_a=\vec{t}$ span a 3-dimensional invariant subspace that
transforms as $\GV$. Rigid rotations about the centre of mass span another one that
transforms as $\GR$. What is left is the space of pure vibrations:
\begin{equation}
  \Gvib=\GM\ominus\GV\ominus\GR,
\end{equation}
a formal subtraction of multiplicities (\cref{ch:ops1} shows that the three
subspaces are orthogonal with respect to the mass-weighted inner product, so the
subtraction is legitimate).

\subsection{Worked example: \texorpdfstring{$\mathrm{CH_3D}$}{CH3D}}\label{sec:CH3D}
\[
\begin{array}{l|ccc}
\Cv{3} & E & 2C_3 & 3\sd\\\hline
\Ninv & 5 & 2 & 3\\
\chiV & 3 & 0 & 1\\
\chiM=\Ninv\chiV & 15 & 0 & 3\\
\chiR & 3 & 0 & -1\\
\chi_{\mathrm{vib}}=\chiM-\chiV-\chiR & 9 & 0 & 3
\end{array}
\]
The identity leaves all five atoms in place; a $C_3$ rotation leaves C and D; each
mirror contains C, D and one H. The reduction formula gives
\begin{align*}
m_{\irr{A}_1}&=\tfrac16(9+0+9)=3,\quad
m_{\irr{A}_2}=\tfrac16(9+0-9)=0,\quad
m_{\irr{E}}=\tfrac16(18+0+0)=3,
\end{align*}
so $\Gvib=3\,\irr{A}_1\oplus3\,\irr{E}$: three non-degenerate and three doubly
degenerate frequencies. (Alternatively, $\GM=4\,\irr{A}_1\oplus\irr{A}_2\oplus5\,\irr{E}$,
and subtracting $\GV=\irr{A}_1\oplus\irr{E}$ and $\GR=\irr{A}_2\oplus\irr{E}$ gives the
same result.)

\begin{keyresult}[Summary: the recipe for molecular vibrations]
\begin{enumerate}
\item Find the point group and its character table.
\item For one element per class, write $\chiV$ from \eqref{eq:chiV} and count $\Ninv$.
\item $\chiM=\Ninv\chiV$, $\chiR=\det\cdot\chiV$, $\chi_{\mathrm{vib}}=\chiM-\chiV-\chiR$.
\item Reduce with \eqref{eq:magic}; read degeneracies from the dimensions and the
      number of distinct frequencies from the multiplicities.
\end{enumerate}
\end{keyresult}

\exercisesection

\begin{exercise}[methane]\label{ex:2:CH4}
Using the table \eqref{eq:Tdchar}, derive
$\GM(\mathrm{CH_4})=\irr{A}_1\oplus\irr{E}\oplus\irr{T}_1\oplus3\,\irr{T}_2$ and
$\Gvib=\irr{A}_1\oplus\irr{E}\oplus2\,\irr{T}_2$. You will need the number of atoms
left invariant by each class of $\Td$.
\end{exercise}

\begin{solution}
Classes $E,8C_3,3C_2,6S_4,6\sd$. The $C_3$ axes pass through C and one H; the $C_2$ and
$S_4$ axes pass only through C; each $\sd$ contains C and two H. Hence
\[
\begin{array}{l|ccccc}
 & E & 8C_3 & 3C_2 & 6S_4 & 6\sd\\\hline
\Ninv & 5&2&1&1&3\\
\chiV & 3&0&-1&-1&1\\
\chiM & 15&0&-1&-1&3\\
\chiR & 3&0&-1&1&-1
\end{array}
\]
with $\chiV$ from \eqref{eq:chiV} ($S_4$: $-1+2\cos90^\circ=-1$). The reduction formula
gives
\begin{align*}
m_{\irr{A}_1}&=\tfrac1{24}(15+0-3-6+18)=1, &
m_{\irr{A}_2}&=\tfrac1{24}(15+0-3+6-18)=0,\\
m_{\irr{E}}&=\tfrac1{24}(30+0-6+0+0)=1, &
m_{\irr{T}_1}&=\tfrac1{24}(45+0+3-6-18)=1,\\
m_{\irr{T}_2}&=\tfrac1{24}(45+0+3+6+18)=3,
\end{align*}
so $\GM=\irr{A}_1\oplus\irr{E}\oplus\irr{T}_1\oplus3\,\irr{T}_2$. Since
$\chiV=\chi_{\irr{T}_2}$ and $\chiR=\chi_{\irr{T}_1}$,
$\Gvib=\irr{A}_1\oplus\irr{E}\oplus2\,\irr{T}_2$.
\end{solution}

\begin{exercise}[ammonia]\label{ex:2:NH3}
Find $\Gvib$ for the pyramidal molecule $\mathrm{NH_3}$ (point group $\Cv{3}$). How
many distinct vibrational frequencies are there?
\end{exercise}

\begin{solution}
$\Ninv=(4,1,2)$ (each mirror contains N and one H), $\chiM=(12,0,2)$, so
$\GM=3\,\irr{A}_1\oplus\irr{A}_2\oplus4\,\irr{E}$. Subtracting
$\GV=\irr{A}_1\oplus\irr{E}$ and $\GR=\irr{A}_2\oplus\irr{E}$:
$\Gvib=2\,\irr{A}_1\oplus2\,\irr{E}$. Four distinct frequencies: two non-degenerate, two
doubly degenerate ($2+4=6=3\cdot4-6$).
\end{solution}

\begin{exercise}[water]\label{ex:2:H2O}
The water molecule lies in the $yz$ plane and has point group $\Cv{2}$, with elements
$E$, $C_2$, $\sigma_{\mathrm{v}}(xz)$ and $\sigma'_{\mathrm{v}}(yz)$. Its character table is
\[
\begin{array}{c|cccc}
\Cv{2} & E & C_2 & \sigma_{\mathrm{v}}(xz) & \sigma'_{\mathrm{v}}(yz)\\\hline
\irr{A}_1 & 1&1&1&1\\ \irr{A}_2&1&1&-1&-1\\ \irr{B}_1&1&-1&1&-1\\ \irr{B}_2&1&-1&-1&1
\end{array}
\]
Show that $\Gvib=2\,\irr{A}_1\oplus\irr{B}_2$.
\end{exercise}

\begin{solution}
$\chiV=(3,-1,1,1)$, $\Ninv=(3,1,1,3)$ (the $xz$ mirror contains only O),
$\chiM=(9,-1,1,3)$, $\chiR=\det\cdot\chiV=(3,-1,-1,-1)$,
$\chi_{\mathrm{vib}}=\chiM-\chiV-\chiR=(3,1,1,3)$. Then
$m_{\irr{A}_1}=\tfrac14(3+1+1+3)=2$, $m_{\irr{A}_2}=\tfrac14(3+1-1-3)=0$,
$m_{\irr{B}_1}=\tfrac14(3-1+1-3)=0$, $m_{\irr{B}_2}=\tfrac14(3-1-1+3)=1$:
$\Gvib=2\,\irr{A}_1\oplus\irr{B}_2$ (symmetric stretch, bend, antisymmetric stretch).
\end{solution}

\begin{exercise}[boron trifluoride]\label{ex:2:BF3}
The planar molecule $\mathrm{BF_3}$ has point group $\Dh{3}$ with classes
$E,\ 2C_3,\ 3C_2',\ \sh,\ 2S_3,\ 3\sv$ and character table
\[
\begin{array}{c|cccccc}
\Dh{3}& E & 2C_3 & 3C_2' & \sh & 2S_3 & 3\sv\\\hline
\irr{A}_1' &1&1&1&1&1&1\\
\irr{A}_2' &1&1&-1&1&1&-1\\
\irr{E}' &2&-1&0&2&-1&0\\
\irr{A}_1''&1&1&1&-1&-1&-1\\
\irr{A}_2''&1&1&-1&-1&-1&1\\
\irr{E}''&2&-1&0&-2&1&0
\end{array}
\]
Compute $\GV$, $\GR$, $\GM$ and $\Gvib$.
\end{exercise}

\begin{solution}
With $\chiV$ from \eqref{eq:chiV} ($S_3$: $-1+2\cos120^\circ=-2$) and the invariant atoms
(each $C_2'$ and each $\sv$ contains B and one F; $\sh$ contains all atoms):
\[
\begin{array}{l|cccccc}
 & E&2C_3&3C_2'&\sh&2S_3&3\sv\\\hline
\Ninv &4&1&2&4&1&2\\
\chiV &3&0&-1&1&-2&1\\
\chiM &12&0&-2&4&-2&2\\
\chiR &3&0&-1&-1&2&-1
\end{array}
\]
Reduction: $\GV=\irr{E}'\oplus\irr{A}_2''$, $\GR=\irr{A}_2'\oplus\irr{E}''$,
$\GM=\irr{A}_1'\oplus\irr{A}_2'\oplus3\,\irr{E}'\oplus2\,\irr{A}_2''\oplus\irr{E}''$, and
therefore $\Gvib=\irr{A}_1'\oplus2\,\irr{E}'\oplus\irr{A}_2''$ ($1+4+1=6$). For instance
$m_{\irr{E}'}(\GM)=\tfrac1{12}(24+0+0+8+4+0)=3$.
\end{solution}

\begin{exercise}\label{ex:2:irredtest}
Use $\inner{\chi}{\chi}$ to show that the vector representation of $\Td$ is
irreducible, while that of $\Cv{3}$ is not.
\end{exercise}

\begin{solution}
$\Td$: $\inner{\chiV}{\chiV}=\tfrac1{24}(9+0+3\cdot1+6\cdot1+6\cdot1)=1$: irreducible.
$\Cv{3}$: $\tfrac16(9+0+3)=2=1^2+1^2$: the sum of two different irreps.
\end{solution}

\chapter{Operations with representations I: tensors and invariants}\label{ch:ops1}
\chaptermark{Tensors and invariants}

\begin{watch}{\lec{3}}
Components vs vectors, active vs passive \ts{3}{259}{4:19} $\cdot$
the column recipe for matrices \ts{3}{710}{11:50} $\cdot$
pseudovectors and mirrors \ts{3}{969}{16:09} $\cdot$
molecular polarizability \ts{3}{1919}{31:59} $\cdot$
invariants from irreps \ts{3}{2359}{39:19} $\cdot$
molecular vibrations: Lagrangian and dynamical matrix \ts{3}{3246}{54:06} $\cdot$
the potential energy of $\mathrm{CH_4}$ in symmetry-adapted coordinates \ts{3}{4192}{69:52} $\cdot$
Wigner's theorem seen explicitly \ts{3}{5017}{83:37} $\cdot$
product representations \ts{3}{5205}{86:45}
\end{watch}

\section{How to build representation matrices}\label{sec:active}

\subsection{Active and passive}
Two opposite conventions exist. In the \emph{active} view, which we always use, the
symmetry operation moves the object (vector, molecule, wave function) and the
coordinate axes stay fixed. In the \emph{passive} view the axes are moved and the
object stays, which is equivalent to moving the object by the inverse operation.
Some classic books (for example Bradley--Cracknell for translations) use the passive
convention; mixing the two is a common source of sign errors.

Do not confuse the \emph{coordinate system} (the axes, fixed) with the
\emph{coordinates} of a vector (its components, which change when the vector is
moved).

\subsection{The column recipe}
Representation matrices act on columns of components,
\[
(v_x',v_y',v_z')\T=\mat{V}(g)\,(v_x,v_y,v_z)\T.
\]
Applying $\mat{V}(g)$ to the basis vector $\vec{e}_1=(1,0,0)\T$ returns the first
column. Hence:
\begin{keyresult}[Column recipe]
The $j$-th column of $\mat{D}(g)$ contains the components of the transformed basis
vector $g\,\vec{e}_j$.
\end{keyresult}
This works for every representation and every basis, and avoids general formulas
for rotation or reflection matrices. As an example, construct
$\mat{V}(\sigma_{\mathrm{d}2})$ for $\Cv{3}$ (\cref{fig:c3v}); the mirror line makes
$120^\circ$ with the $x$ axis.
\begin{itemize}
\item $\vec{e}_x$ (angle $0^\circ$) is reflected to the direction at $240^\circ$:
      $g\vec{e}_x=(-\tfrac12,-\tfrac{\sqrt3}{2},0)$.
\item $\vec{e}_y$ (angle $90^\circ$) is reflected to $150^\circ$:
      $g\vec{e}_y=(-\tfrac{\sqrt3}{2},\tfrac12,0)$.
\item $\vec{e}_z$ lies in the mirror and is unchanged.
\end{itemize}
Placing these as columns reproduces $\mat{V}(\sigma_{\mathrm{d}2})$ in \eqref{eq:VC3v}.

\subsection{Pseudovectors, again}
The axial-vector matrices are $\mat{R}(g)=\det\mat{V}(g)\,\mat{V}(g)$: identical to
$\mat{V}(g)$ for proper rotations, opposite for improper operations. For instance
$\mat{R}(\sigma_{\mathrm{d}1})=\diag(-1,1,-1)$ reverses the components of a
pseudovector lying in the mirror ($x$, $z$) and keeps the normal one ($y$).

\section{Tensors and their invariance}

A physical property of a molecule or crystal in equilibrium is typically a tensor.
Our first example is the electric \emph{polarizability} $\alpha_{ij}$, which relates
the induced dipole to an applied field:
\begin{equation}
  p_i=\alpha_{ij}E_j\qquad\text{(sum over repeated indices).}
\end{equation}
A rank-2 tensor transforms like a product of coordinates, $x_ix_j$: one factor of
$\mat{V}$ for each index,
\begin{equation}\label{eq:tensortransf}
  \alpha'_{ij}=V_{ik}V_{jl}\,\alpha_{kl}\qquad\Longleftrightarrow\qquad
  \mat{\alpha}'=\mat{V}\mat{\alpha}\mat{V}\T=\mat{V}\mat{\alpha}\mat{V}^{-1},
\end{equation}
where we used that $\mat{V}(g)$ is orthogonal. If $g$ is a symmetry, the molecule looks
the same after the operation and its polarizability cannot change,
$\mat{\alpha}'=\mat{\alpha}$. Equivalently $[\mat{\alpha},\mat{V}(g)]=0$ for all
$g\in G$.

The brute-force approach writes the most general symmetric $3\times3$ matrix and
imposes these commutation relations. That is feasible for $\mat{\alpha}$, tedious for
the $6\times6$ elastic tensor, and hopeless for the $15\times15$ dynamical matrix of
methane. Group theory offers a shortcut.

\section{Invariants}\label{sec:invariants}

Contract the tensor into a scalar. The energy of the molecule in a static field is
\begin{equation}
  U=-\tfrac12\,\alpha_{ij}E_iE_j .
\end{equation}
Being a property of the molecule, $U$ must be invariant under $G$; being quadratic in
$\vec{E}$, it must be a linear combination of the independent \emph{quadratic
invariants} that can be built from the components of $\vec{E}$. Group theory tells us
all of them.

\begin{keyresult}[Quadratic invariants]
Decompose the vector space into irreps with symmetry-adapted coordinates.
Let all irreps be unitary and \emph{real} (orthogonal matrices). Then:
\begin{enumerate}
\item within each irrep the squared norm $\sum_i x_i^2$ is invariant, and it is the
      only quadratic invariant built from that irrep alone;
\item if an irrep occurs twice, with coordinates $x_i$ and $y_i$ written \emph{in the
      same basis} (same matrices), the scalar product $\sum_i x_iy_i$ is an additional
      invariant;
\item products of coordinates belonging to inequivalent irreps are never invariant.
\end{enumerate}
\end{keyresult}
Point 1 follows from unitarity: an orthogonal transformation preserves the norm.
That the norm is the \emph{only} invariant, and that inequivalent irreps give none,
is a consequence of Schur's lemma (\cref{sec:schur}), as shown in \cref{sec:bilinear}.
For complex irreps the recipe changes slightly (\cref{sec:complexinv}).

\begin{example}[Polarizability of $\mathrm{CH_3D}$]
In $\Cv{3}$, $\GV=\irr{A}_1\oplus\irr{E}$ with coordinates $E_z$ ($\irr{A}_1$) and
$(E_x,E_y)$ ($\irr{E}$). The quadratic invariants are $E_z^2$ and $E_x^2+E_y^2$, so
\[
U=-\tfrac12\bigl[a(E_x^2+E_y^2)+bE_z^2\bigr],\qquad
\mat{\alpha}=\begin{pmatrix}a&0&0\\0&a&0\\0&0&b\end{pmatrix}.
\]
Any molecule or crystal with symmetry $\Cv{3}$ has a polarizability determined by two
constants. Group theory cannot give their values.
\end{example}

\begin{example}[Polarizability of $\mathrm{CH_4}$]
In $\Td$, $\GV=\irr{T}_2$ is irreducible. The only quadratic invariant is
$E_x^2+E_y^2+E_z^2$, so $\mat{\alpha}=a\one$: the polarizability of methane is
isotropic, although the molecule is not spherical. (This fact decides its rotational
Raman spectrum, \cref{sec:raman}.)
\end{example}

\section{Molecular vibrations}\label{sec:vibrations}

Let $\vec{R}_a^{0}$ be the equilibrium position of atom $a$ (measured from the
centre of mass) and $\vec{u}_a$ its displacement. Collect the displacements in the
$3N$-vector $\mat{u}$. In the harmonic approximation
\begin{equation}
  L=\tfrac12\,\dot{\mat{u}}\T\mat{M}\,\dot{\mat{u}}-\tfrac12\,\mat{u}\T\mat{\Phi}\,\mat{u},
\end{equation}
with the diagonal mass matrix $\mat{M}=\diag(m_1,m_1,m_1,m_2,\dots)$ and the matrix of
force constants $\Phi_{ai,bj}=\partial^2V/\partial u_{ai}\partial u_{bj}$ at
equilibrium. The equations of motion $\mat{M}\ddot{\mat{u}}+\mat{\Phi}\mat{u}=0$ with
$\mat{u}(t)=\Real(\mat{n}\,\ee^{-\ii\omega t})$ give the generalised eigenvalue
problem $\mat{\Phi}\mat{n}=\omega^2\mat{M}\mat{n}$. With mass-weighted coordinates
$\mat{w}=\mat{M}^{1/2}\mat{u}$ it becomes an ordinary one,
\begin{equation}
  \mat{D}\,\mat{w}=\omega^2\mat{w},\qquad \mat{D}=\mat{M}^{-1/2}\mat{\Phi}\mat{M}^{-1/2}.
\end{equation}
The real symmetric \emph{dynamical matrix} $\mat{D}$ has $3N$ eigenvalues $\omega^2$
and orthonormal eigenvectors (the \emph{normal modes}). In the original variables
the modes are orthogonal with respect to the metric $\mat{M}$.

The $3N$-dimensional space splits into three $\mat{M}$-orthogonal invariant subspaces:
\begin{itemize}
\item translations, $\vec{u}_a=\vec{t}$ for all $a$ (transform as $\GV$);
\item rigid rotations, $\vec{u}_a=\vec{\varphi}\times\vec{R}_a^0$ (transform as $\GR$);
\item pure vibrations, defined by the six constraints
      $\sum_am_a\vec{u}_a=0$ (no displacement of the centre of mass) and
      $\sum_am_a\vec{R}_a^0\times\vec{u}_a=0$ (no angular momentum at linear order).
\end{itemize}
The first two cost no potential energy, $\omega=0$. The third is awkward to describe
explicitly, and group theory never needs to: it only uses its character
$\chiM-\chiV-\chiR$.

A symmetry operation only exchanges atoms of the same mass, so $\mat{M}(g)$ commutes
with the mass matrix, $\mat{M}(g)\mat{M}=\mat{M}\,\mat{M}(g)$. By the same argument as
for the polarizability, $\mat{D}=\mat{M}(g)\,\mat{D}\,\mat{M}(g)\T$: the dynamical matrix
is an invariant rank-2 tensor on the space of distortions.

\section{Wigner's theorem, seen explicitly}\label{sec:wigner}

Take methane, $\Gvib=\irr{A}_1\oplus\irr{E}\oplus2\,\irr{T}_2$, and assume we use
symmetry-adapted vibrational coordinates:
\[
\underbrace{q_1}_{\irr{A}_1},\qquad
\underbrace{(q_2,q_3)}_{\irr{E}},\qquad
\underbrace{(q_4,q_5,q_6)}_{\irr{T}_2},\qquad
\underbrace{(q_7,q_8,q_9)}_{\irr{T}_2},
\]
where the two $\irr{T}_2$ triplets transform with \emph{identical} matrices. (How to
find such coordinates with projection operators is explained in \cref{sec:projectors}.) The
potential energy is an invariant quadratic form, so by \cref{sec:invariants} it must
be
\begin{equation}\label{eq:VCH4}
  2V=a\,q_1^2+b\,(q_2^2+q_3^2)+c\,(q_4^2+q_5^2+q_6^2)+d\,(q_7^2+q_8^2+q_9^2)
     +2e\,(q_4q_7+q_5q_8+q_6q_9).
\end{equation}
Five constants, instead of the $45$ independent entries of a general symmetric
$9\times9$ matrix. Now reorder the coordinates, pairing the first component of the
first triplet with the first component of the second, and so on:
$(q_1;\,q_2;\,q_3;\,q_4,q_7;\,q_5,q_8;\,q_6,q_9)$. In this order the matrix of the
quadratic form is block diagonal,
\begin{equation}
\mat{D}_{\mathrm{vib}}=
\diag\Bigl(a,\ b,\ b,\
\begin{pmatrix}c&e\\e&d\end{pmatrix},
\begin{pmatrix}c&e\\e&d\end{pmatrix},
\begin{pmatrix}c&e\\e&d\end{pmatrix}\Bigr).
\end{equation}
Its eigenvalues are $a$ (once), $b$ (twice) and the two eigenvalues
$\omega_\pm^2$ of the $2\times2$ block (three times each): four distinct frequencies
with degeneracies $1,2,3,3$. The mechanism is general:

\begin{keyresult}[Wigner's theorem, block form]
If an invariant hermitian operator acts on a space carrying
$\bigoplus_\mu m_\mu\Gamma^{(\mu)}$, then in a symmetry-adapted basis (ordered as
above) it consists, for each irrep $\mu$, of one $m_\mu\times m_\mu$ block repeated
$d_\mu$ times. Hence $m_\mu$ distinct eigenvalues, each $d_\mu$-fold degenerate.
\end{keyresult}
The same statement with ``Hamiltonian'' for ``dynamical matrix'' is the basis of all
quantum-mechanical applications: energy levels are labelled by irreps, and the
degeneracy of a level equals the dimension of its irrep, unless two levels coincide
\emph{accidentally}.

\begin{remark}[Symmetry-adapted vs normal coordinates]
In \eqref{eq:VCH4}, $q_1,q_2,q_3$ are already normal coordinates (their blocks are
$1\times1$). The $\irr{T}_2$ normal coordinates are linear combinations such as
$Q_4=\cos\phi\,q_4+\sin\phi\,q_7$ that diagonalise the $2\times2$ block; symmetry
alone cannot fix the mixing angle $\phi$.
\end{remark}

\section{Projection operators\beyondmark}\label{sec:projectors}

Wigner's theorem is used in a symmetry-adapted basis. Projection operators construct
such a basis from any starting vector. Let $\Gamma$ act on a space $L$ with operators
$\hat{O}_g$, and define
\begin{equation}\label{eq:projector}
\hat{P}^{(\mu)}=\frac{d_\mu}{\abs{G}}\sum_g\chi^{(\mu)}(g)^{*}\,\hat{O}_g,\qquad
\hat{P}^{(\mu)}_{ij}=\frac{d_\mu}{\abs{G}}\sum_gD^{(\mu)}_{ij}(g)^{*}\,\hat{O}_g .
\end{equation}
\begin{theorem}
$\hat{P}^{(\mu)}$ is the projector onto the subspace of $L$ that transforms as
$\Gamma^{(\mu)}$ (all its copies). If $\phi$ is any vector, $\phi_i=\hat{P}^{(\mu)}_{ii}\phi$
(no sum) is either zero or the $i$-th partner of a basis transforming exactly with the
matrices $\mat{D}^{(\mu)}$, and the other partners are $\phi_j=\hat{P}^{(\mu)}_{ji}\phi$.
\end{theorem}
The proof applies $\hat{P}^{(\mu)}_{ij}$ to a symmetry-adapted basis and uses the great
orthogonality theorem \eqref{eq:GOT}. In practice:
\begin{enumerate}
\item Take any trial function (an atomic orbital, a displacement of one atom).
\item Apply $\hat{P}^{(\mu)}$ to get its $\mu$ component. Repeat with other trial
      functions until you have $m_\mu d_\mu$ independent results.
\item For multidimensional irreps, use $\hat{P}^{(\mu)}_{ji}$ (needs the matrices) to
      generate partners transforming with the tabulated matrices. This is required for
      the invariant constructions of \cref{sec:invariants}.
\end{enumerate}

\begin{example}[Three hydrogen orbitals in $\Cv{3}$]
Let $h_1,h_2,h_3$ be the $1s$ orbitals of the three hydrogens of $\mathrm{CH_3D}$. The
operations permute them: $C_3^+\colon h_1\to h_2\to h_3\to h_1$, and
$\sigma_{\mathrm{d}1}$ fixes $h_1$ and swaps $h_2\leftrightarrow h_3$. The permutation
character is $(3,0,1)$, so the span is $\irr{A}_1\oplus\irr{E}$. Projecting $h_1$:
\begin{align*}
\hat{P}^{(\irr{A}_1)}h_1&=\tfrac16\,(h_1+h_2+h_3+h_1+h_3+h_2)
    =\tfrac13(h_1+h_2+h_3),\\
\hat{P}^{(\irr{E})}h_1&=\tfrac26\,\bigl(2h_1-h_2-h_3+0\bigr)
    =\tfrac13(2h_1-h_2-h_3).
\end{align*}
The second $\irr{E}$ partner is the orthogonal combination $h_2-h_3$ (odd under
$\sigma_{\mathrm{d}1}$, like $y$). Normalised symmetry-adapted combinations:
$\tfrac{1}{\sqrt3}(h_1+h_2+h_3)$ for $\irr{A}_1$, and
$\bigl(\tfrac{1}{\sqrt6}(2h_1-h_2-h_3),\ \tfrac{1}{\sqrt2}(h_2-h_3)\bigr)$ for
$\irr{E}$, which transform like $(x,y)$.
\end{example}

\section{Product representations}\label{sec:product}

The nine products $x_ix_j'$ of the components of two vectors span the tensor product
space $\mathbb{R}^3\otimes\mathbb{R}^3$, and transform as
$(x_ix_j')'=V_{ik}V_{jl}\,x_kx_l'$. Treating the pair $(ij)$ as a single index, the
$9\times9$ matrix
\begin{equation}
  \bigl[\mat{V}\otimes\mat{V}\bigr]_{(ij),(kl)}=V_{ik}V_{jl}
\end{equation}
(the \emph{Kronecker product}) defines the \emph{product representation}
$\GV\otimes\GV$. More generally, if $\Gamma_1$ acts on $L_1$ and $\Gamma_2$ on
$L_2$, then $\Gamma_1\otimes\Gamma_2$ acts on $L_1\otimes L_2$ with matrices
$\mat{D}_1(g)\otimes\mat{D}_2(g)$, of dimension $d_1d_2$. Rank-2 tensors
``belong to'' $\GV\otimes\GV$. As with the direct sum, the matrices of the product
representation are not products of the original matrices. Their traces, however,
are products of the traces (\cref{ch:ops2}):
\begin{equation}\label{eq:prodchar}
  \chi_{\Gamma_1\otimes\Gamma_2}(g)=\chi_{\Gamma_1}(g)\,\chi_{\Gamma_2}(g).
\end{equation}

\section{Bilinear invariants: the general statement}\label{sec:bilinear}

\begin{theorem}\label{thm:bilinear}
Let $\Gamma^{(\mu)}$ and $\Gamma^{(\nu)}$ be real orthogonal irreps with coordinates
$x_i$ and $y_j$. A bilinear form $B=\sum_{ij}b_{ij}x_iy_j$ is invariant under $G$ if and
only if $\mat{b}\,\mat{D}^{(\nu)}(g)=\mat{D}^{(\mu)}(g)\,\mat{b}$ for all $g$. By
Schur's lemma (\cref{sec:schur}):
\begin{enumerate}
\item if $\Gamma^{(\mu)}$ and $\Gamma^{(\nu)}$ are inequivalent, then $\mat{b}=0$: no
      invariant;
\item if they are identical (same matrices), then $\mat{b}=\lambda\one$: the only
      invariant is the scalar product $\sum_ix_iy_i$.
\end{enumerate}
\end{theorem}
\begin{proof}
Invariance of $B$ under $x\mapsto\mat{D}^{(\mu)}x$, $y\mapsto\mat{D}^{(\nu)}y$ means
$\mat{D}^{(\mu)\mathrm{T}}\mat{b}\,\mat{D}^{(\nu)}=\mat{b}$, which for orthogonal
matrices is the intertwining condition above. By Schur I, $\mat{b}=0$ for
inequivalent irreps. For identical irreps Schur II gives $\mat{b}=\lambda\one$. For real
irreps Schur II holds over $\mathbb{R}$ as well, because a real matrix commuting with a
real, absolutely irreducible representation is still a multiple of the identity.
\end{proof}
If the two irreps are equivalent but
written in different bases, $\mat{D}^{(\nu)}=\mat{A}\mat{D}^{(\mu)}\mat{A}^{-1}$,
there is still one invariant, but it is $x\T\mat{A}^{-1}y$ rather than the scalar
product. Always write repeated irreps with the same matrices.

The number of independent bilinear invariants in $\Gamma_1\otimes\Gamma_2$ is the
multiplicity of the identity irrep, $m_{\irr{A}_1}(\Gamma_1\otimes\Gamma_2)=
\inner{\chi_{\Gamma_1}^{*}}{\chi_{\Gamma_2}}$, which by orthogonality equals $1$ if
$\Gamma_2\simeq\Gamma_1^{*}$ and $0$ otherwise (\cref{ch:ops2}).

\exercisesection

\begin{exercise}\label{ex:3:sigma3}
Use the column recipe to construct $\mat{V}(\sigma_{\mathrm{d}3})$ and
$\mat{V}(C_3^-)$ for $\Cv{3}$, and check them against \eqref{eq:VC3v}.
\end{exercise}

\begin{solution}
$\sigma_{\mathrm{d}3}$ is the mirror line at $240^\circ$ (equivalently $60^\circ$). A
direction at angle $\theta$ is reflected to $2\cdot240^\circ-\theta$: $\vec{e}_x$
($0^\circ$) goes to $480^\circ\equiv120^\circ$, $(-\tfrac12,\tfrac{\sqrt3}{2},0)$;
$\vec{e}_y$ ($90^\circ$) goes to $390^\circ\equiv30^\circ$, $(\tfrac{\sqrt3}{2},\tfrac12,0)$;
$\vec{e}_z$ is fixed. For $C_3^-$: $\vec{e}_x\to(-\tfrac12,-\tfrac{\sqrt3}{2},0)$ (angle
$-120^\circ$), $\vec{e}_y\to(\tfrac{\sqrt3}{2},-\tfrac12,0)$ (angle $-30^\circ$). Placing
these as columns reproduces \eqref{eq:VC3v}.
\end{solution}

\begin{exercise}\label{ex:3:orthsub}
Show that uniform translations and rigid rotations are orthogonal to each other and
to all pure vibrations with respect to the inner product
$\inner{\mat{u}}{\mat{v}}_M=\mat{u}\T\mat{M}\mat{v}$.
\end{exercise}

\begin{solution}
Translation $\vec{t}$ versus any $\mat{v}$:
$\sum_am_a\vec{t}\cdot\vec{v}_a=\vec{t}\cdot\sum_am_a\vec{v}_a$, which vanishes for rotations
($\sum_am_a\vec{\varphi}\times\vec{R}_a^0=\vec{\varphi}\times\sum_am_a\vec{R}_a^0=0$, the centre of
mass being the origin) and for vibrations (first constraint). Rotation versus vibration:
$\sum_am_a(\vec{\varphi}\times\vec{R}_a^0)\cdot\vec{u}_a=\vec{\varphi}\cdot\sum_am_a\vec{R}_a^0\times\vec{u}_a=0$
(second constraint).
\end{solution}

\begin{exercise}[conductivity tensor]\label{ex:3:cond}
A crystal has point group $\Cv{4}$ (a four-fold axis along $z$ and four vertical
mirrors). Using quadratic invariants, find the most general symmetric conductivity
tensor $\sigma_{ij}$. Compare with $\Cv{3}$ and $\Td$.
\end{exercise}

\begin{solution}
In $\Cv{4}$, $\GV=\irr{A}_1(z)\oplus\irr{E}(x,y)$, both real. The quadratic invariants are
$E_z^2$ and $E_x^2+E_y^2$, hence $\mat{\sigma}=\diag(\sigma_\perp,\sigma_\perp,\sigma_\parallel)$
with two constants, exactly as for $\Cv{3}$ (uniaxial crystals). For $\Td$ (cubic),
$\mat{\sigma}=\sigma\one$ is isotropic.
\end{solution}

\begin{exercise}\label{ex:3:NH3pot}
For $\mathrm{NH_3}$, $\Gvib=2\,\irr{A}_1\oplus2\,\irr{E}$ (\cref{ex:2:NH3}). How many
independent constants does the most general harmonic potential energy contain?
Describe the block structure of the dynamical matrix in symmetry-adapted coordinates.
\end{exercise}

\begin{solution}
$\Gvib=2\,\irr{A}_1\oplus2\,\irr{E}$. The $\irr{A}_1$ part gives a $2\times2$ symmetric block
($3$ constants: two squares and one cross term); the $\irr{E}$ part gives a $2\times2$ block
repeated twice ($3$ constants). Six constants in total, instead of the $21$ of a general
symmetric $6\times6$ matrix. Two non-degenerate and two doubly degenerate frequencies.
\end{solution}

\begin{exercise}\label{ex:3:prod}
For $\Cv{3}$ compute the character of $\irr{E}\otimes\irr{E}$ and decompose it. How
many independent bilinear invariants can be built from two $\irr{E}$ doublets
$(x_1,x_2)$ and $(y_1,y_2)$? Write them down.
\end{exercise}

\begin{solution}
$\chi_{\irr{E}\otimes\irr{E}}=(4,1,0)$, so
$m_{\irr{A}_1}=\tfrac16(4+2)=1$, $m_{\irr{A}_2}=\tfrac16(4+2)=1$, $m_{\irr{E}}=\tfrac16(8-2)=1$:
$\irr{E}\otimes\irr{E}=\irr{A}_1\oplus\irr{A}_2\oplus\irr{E}$. There is one bilinear invariant,
$x_1y_1+x_2y_2$ (both doublets written with the same real matrices). The $\irr{A}_2$
combination $x_1y_2-x_2y_1$ is invariant under $C_3$ but changes sign under the mirrors.
The $\irr{E}$ combination is $(x_1y_1-x_2y_2,\,-x_1y_2-x_2y_1)$.
\end{solution}

\chapter{Operations with representations II: spectroscopy and selection rules}\label{ch:ops2}
\chaptermark{Spectroscopy and selection rules}

\begin{watch}{\lec{4}}
Characters of product representations \ts{4}{123}{2:03} $\cdot$
microwave and infrared spectroscopy (blackboard) \ts{4}{632}{10:32} $\cdot$
counting bilinear invariants \ts{4}{1424}{23:44} $\cdot$
microwave activity \ts{4}{1735}{28:55} $\cdot$
infrared activity \ts{4}{1922}{32:02} $\cdot$
symmetry-adapted vs normal coordinates \ts{4}{2557}{42:37} $\cdot$
symmetric tensors and the symmetric square \ts{4}{2816}{46:56} $\cdot$
Raman scattering (blackboard) \ts{4}{3304}{55:04} $\cdot$
Raman activity \ts{4}{4570}{76:10}
\end{watch}

\section{Characters of product representations}

Let $\Gamma_1$ and $\Gamma_2$ have matrices $\mat{D}_1(g)$ and $\mat{D}_2(g)$. The
product representation acts on objects $t_{ij}$ with one index of each kind:
\[
t'_{ij}=\bigl[D_1(g)\bigr]_{ik}\bigl[D_2(g)\bigr]_{jl}\,t_{kl}
\equiv\bigl[\mat{D}_1\otimes\mat{D}_2\bigr]_{(ij),(kl)}\,t_{kl}.
\]
The trace sets the double index $(kl)$ equal to $(ij)$ and sums:
\[
\chi_{\Gamma_1\otimes\Gamma_2}(g)=\sum_{i,j}\bigl[D_1(g)\bigr]_{ii}\bigl[D_2(g)\bigr]_{jj}
=\chi_1(g)\,\chi_2(g).
\]
So, just as characters of direct sums add, characters of products multiply.

\begin{example}[$\GV\otimes\GV$ for $\Cv{3}$]
$\chiV=(3,0,1)$, hence $\chi_{\GV\otimes\GV}=(9,0,1)$ and
\[
m_{\irr{A}_1}=\tfrac16(9+0+3)=2,\quad m_{\irr{A}_2}=\tfrac16(9+0-3)=1,\quad
m_{\irr{E}}=\tfrac16(18+0+0)=3,
\]
so $\GV\otimes\GV=2\,\irr{A}_1\oplus\irr{A}_2\oplus3\,\irr{E}$ ($2+1+6=9$).
\end{example}

\section{Counting bilinear invariants}\label{sec:countinv}

A bilinear quantity $\sum b_{ij}x_iy_j$ built from $x\in\Gamma_1$ and $y\in\Gamma_2$
belongs to $\Gamma_1\otimes\Gamma_2$. It is invariant if it lies in a copy of the
identity representation $\irr{A}_1$ (all characters $1$). Therefore
\begin{equation}\label{eq:ninv}
\begin{aligned}
\#\{\text{independent bilinear invariants}\}
 &=m_{\irr{A}_1}(\Gamma_1\otimes\Gamma_2)\\
 &=\frac{1}{\abs{G}}\sum_g\chi_1(g)\chi_2(g)
 =\inner{\chi_1^{*}}{\chi_2}.
\end{aligned}
\end{equation}
If $\Gamma_1=\Gamma^{(\mu)}$ and $\Gamma_2=\Gamma^{(\nu)}$ are irreps, orthogonality
gives $m_{\irr{A}_1}=1$ when $\Gamma^{(\nu)}$ is equivalent to the complex conjugate
of $\Gamma^{(\mu)}$ and $0$ otherwise. For real irreps
($\Gamma^{(\mu)*}=\Gamma^{(\mu)}$): \emph{an invariant can only be formed from two
copies of the same irrep}, and then there is exactly one, the scalar product. This
proves the recipe of \cref{sec:invariants}.

\section{Symmetric and antisymmetric squares}\label{sec:symsq}

A \emph{symmetric} rank-2 tensor ($t_{ij}=t_{ji}$, like the polarizability) lives in
the 6-dimensional subspace of $\mathbb{R}^3\otimes\mathbb{R}^3$ spanned by the
symmetric products $x^2$, $y^2$, $z^2$, $xy$, $yz$, $zx$. It carries the \emph{symmetric square}
$[\Gamma^2]$. The antisymmetric tensors ($t_{ij}=-t_{ji}$, 3 components) carry the
\emph{antisymmetric square} $\{\Gamma^2\}$. Their characters are
\begin{equation}\label{eq:symsq}
\chi_{[\Gamma^2]}(g)=\tfrac12\bigl[\chi(g)^2+\chi(g^2)\bigr],\qquad
\chi_{\{\Gamma^2\}}(g)=\tfrac12\bigl[\chi(g)^2-\chi(g^2)\bigr].
\end{equation}
\emph{Proof.} Diagonalise $\mat{D}(g)$ (possible for a unitary matrix) with
eigenvalues $\lambda_i$. The symmetric products $e_ie_j$ ($i\le j$) are eigenvectors
with eigenvalues $\lambda_i\lambda_j$, so
$\chi_{[\Gamma^2]}=\sum_{i\le j}\lambda_i\lambda_j=\tfrac12\bigl[(\sum\lambda_i)^2+\sum\lambda_i^2\bigr]$,
and $\sum_i\lambda_i^2=\tr\mat{D}(g)^2=\chi(g^2)$. The antisymmetric case is the same
with $i<j$. \qed

To use \eqref{eq:symsq} you need, for each class, the class of $g^2$. For $\Cv{3}$:
$E^2=E$, $(C_3^{\pm})^2=C_3^{\mp}$ (same class), $\sigma^2=E$. Hence
\[
\begin{array}{l|ccc}
\Cv{3} & E & 2C_3 & 3\sd\\\hline
\chiV(g) & 3 & 0 & 1\\
\chiV(g^2) & 3 & 0 & 3\\
\chi_{[\GV^2]} & 6 & 0 & 2\\
\chi_{\{\GV^2\}} & 3 & 0 & -1
\end{array}
\qquad\Longrightarrow\qquad
[\GV^2]=2\,\irr{A}_1\oplus2\,\irr{E},\quad \{\GV^2\}=\irr{A}_2\oplus\irr{E}.
\]
Notice that $\{\GV^2\}=\GR$: an antisymmetric tensor $t_{ij}$ is equivalent to the
axial vector $\tfrac12\varepsilon_{ijk}t_{jk}$, just as $\vec{r}\times\vec{p}$ is built
from the antisymmetric products $x_ip_j-x_jp_i$ (\cref{ex:4:antisym}).

For $\Td$ ($C_2^2=E$, $S_4^2=C_2$, $\sd^2=E$) one finds
$[\GV^2]=\irr{A}_1\oplus\irr{E}\oplus\irr{T}_2$ and $\{\GV^2\}=\irr{T}_1=\GR$.

\section{A little spectroscopy}

We need only a qualitative picture of three kinds of molecular spectroscopy.

\paragraph{Electric-dipole coupling.} A molecule is a neutral cloud of charges. To
lowest order in the multipole expansion it couples to a (long-wavelength)
electromagnetic field through its dipole moment, $H_{\mathrm{int}}=-\vec{p}\cdot\vec{E}$.

\paragraph{Microwave (rotational) spectroscopy.} A molecule with a \emph{permanent}
dipole $\vec{p}_0$ that rotates produces an oscillating dipole and can absorb or emit
at its rotation frequencies. For a diatomic molecule $H=\hat{L}^2/2I$ gives levels
$\hbar^2l(l+1)/2I$ and the selection rule $\Delta l=\pm1$ produces equally spaced
lines in the microwave range. $\mathrm{HCl}$ is active; $\mathrm{H_2}$ has no
dipole and is inactive.

\paragraph{Infrared (vibrational) spectroscopy.} A vibration that makes the dipole
oscillate, $\delta\vec{p}(t)\neq0$, absorbs infrared light at its frequency. The
breathing mode of $\mathrm{BF_3}$ keeps the dipole zero and is inactive; a mode that
pushes the boron out of the plane against the fluorines creates an oscillating dipole
and is active.

\paragraph{Raman scattering.} Light of frequency $\omega_0$ far from any electronic
resonance induces a dipole $\vec{p}=\mat{\alpha}\vec{E}$ in the electron cloud, which
re-radiates at $\omega_0$ (Rayleigh scattering, the blue of the sky). By the
Born--Oppenheimer approximation the electron cloud follows the nuclei. If the molecule
rotates or vibrates, $\mat{\alpha}(t)$ oscillates at the rotation or vibration
frequency $\Omega$, and the product $\mat{\alpha}(t)\vec{E}_0\cos\omega_0t$ contains the
frequencies $\omega_0\pm\Omega$: weak side lines. Quantum mechanically the photon
loses energy to the molecule (Stokes line, $\omega_0-\Omega$) or gains it
(anti-Stokes, $\omega_0+\Omega$). The anti-Stokes line needs a thermally excited
vibration and is much weaker at room temperature.

\section{Selection rules from invariants}\label{sec:raman}

In each case we write the most general interaction energy allowed by symmetry and
ask whether it can be non-zero. The strategy is always the same: the energy must be an
\emph{invariant}, and \eqref{eq:ninv} tells when an invariant exists.

\paragraph{Microwave activity.} A permanent dipole is a property of the molecule at
rest, hence invariant under $G$. A non-zero invariant vector exists only if some
component of a vector is invariant:
\begin{keyresult}
A molecule is microwave active $\iff$ $\irr{A}_1\subset\GV$. The permanent dipole lies
in the $\irr{A}_1$ subspace.
\end{keyresult}
$\mathrm{CH_4}$: $\GV=\irr{T}_2$, inactive, with no need to know anything about its
charge distribution. $\mathrm{CH_3D}$: $\GV=\irr{A}_1\oplus\irr{E}$, active, with the
dipole along the three-fold axis.

\paragraph{Infrared activity.} Expand the dipole in the normal coordinates,
$\delta p_i=\sum_\alpha c_{i\alpha}Q_\alpha$. The interaction
$H_{\mathrm{IR}}=-\sum_{i,\alpha}c_{i\alpha}E_iQ_\alpha$ is a bilinear invariant in
$\GV\otimes\Gvib$. A mode transforming as $\Gamma^{(\mu)}$ appears in $H_{\mathrm{IR}}$
only if $\Gamma^{(\mu)}$ is also contained in $\GV$ (for real irreps):
\begin{keyresult}
A vibrational mode of symmetry $\Gamma^{(\mu)}$ is infrared active $\iff$
$\Gamma^{(\mu)}\subset\GV$.
\end{keyresult}

\paragraph{Rotational Raman activity.} A rotating molecule modulates $\mat{\alpha}$
only if $\mat{\alpha}$ is anisotropic, $\mat{\alpha}\neq a\one$. The number of
independent invariant symmetric tensors is $m_{\irr{A}_1}([\GV^2])$; one of them is
always $\one$. Hence:
\begin{keyresult}
Rotational Raman activity $\iff$ $m_{\irr{A}_1}([\GV^2])\geq2$. For point groups this
happens exactly when $\GV$ is reducible.
\end{keyresult}

\paragraph{Vibrational Raman activity.} To first order in the normal coordinates,
$\delta\alpha_{ij}=\sum_\alpha a_{ij\alpha}Q_\alpha$, so the Raman Hamiltonian
$H_{\mathrm{R}}=-\tfrac12\sum a_{ij\alpha}E_iE_jQ_\alpha$ is a trilinear invariant in
$[\GV^2]\otimes\Gvib$:
\begin{keyresult}
A vibrational mode of symmetry $\Gamma^{(\mu)}$ is Raman active $\iff$
$\Gamma^{(\mu)}\subset[\GV^2]$.
\end{keyresult}

\subsection{Methane and deuterated methane}
\[
\begin{array}{lcccc}
\toprule
 & \Gvib & \GV & [\GV^2] & \\
\midrule
\mathrm{CH_4}\ (\Td) & \irr{A}_1\oplus\irr{E}\oplus2\,\irr{T}_2 & \irr{T}_2 &
   \irr{A}_1\oplus\irr{E}\oplus\irr{T}_2 &\\
\mathrm{CH_3D}\ (\Cv{3}) & 3\,\irr{A}_1\oplus3\,\irr{E} & \irr{A}_1\oplus\irr{E} &
   2\,\irr{A}_1\oplus2\,\irr{E} &\\
\bottomrule
\end{array}
\]
\begin{itemize}
\item $\mathrm{CH_4}$: microwave inactive ($\irr{A}_1\not\subset\GV$); only the two
      $\irr{T}_2$ frequencies are infrared active; all four frequencies
      ($\irr{A}_1,\irr{E},\irr{T}_2,\irr{T}_2$) are Raman active; no rotational Raman
      spectrum ($\GV$ irreducible, $\mat{\alpha}$ isotropic).
\item $\mathrm{CH_3D}$: microwave active; all six frequencies are infrared and Raman
      active; rotational Raman active.
\end{itemize}
This completes the list of predictions announced in \cref{sec:methane}.

\begin{beyond}[Beyond the lectures: the general selection rule]
In quantum mechanics the same reasoning is phrased with matrix elements. Let
$\ket{\psi_i}$ belong to $\Gamma_i$, let the components of an operator $\hat{O}$ form a
set transforming as $\Gamma_O$, and let $\ket{\psi_f}$ belong to $\Gamma_f$. Then
\[
\bra{\psi_f}\hat{O}\ket{\psi_i}\neq0\quad\text{only if}\quad
\irr{A}_1\subset\Gamma_f^{*}\otimes\Gamma_O\otimes\Gamma_i
\quad\Longleftrightarrow\quad
\Gamma_f\subset\Gamma_O\otimes\Gamma_i .
\]
The matrix element is invariant under the group (it is a number), so it can only
pick up the identity component of the triple product. For electric-dipole transitions
$\Gamma_O=\GV$; for magnetic-dipole transitions $\Gamma_O=\GR$; for Raman and
quadrupole transitions $\Gamma_O=[\GV^2]$. The Wigner--Eckart theorem
(\cref{sec:wignereckart}) refines the statement: all non-zero matrix elements within a
pair of multiplets are fixed up to one constant per copy of $\Gamma_f$ in
$\Gamma_O\otimes\Gamma_i$.
\end{beyond}

\exercisesection

\begin{exercise}\label{ex:4:antisym}
Show from \eqref{eq:symsq} and \eqref{eq:chiV} that $\chi_{\{\GV^2\}}(g)=\chiR(g)$
for every rotation $C(\theta)$ and every rotoreflection $S(\theta)$. Conclude that
antisymmetric rank-2 tensors transform as axial vectors in every point group.
\end{exercise}

\begin{solution}
Rotation: $\chi=1+2c$ ($c=\cos\theta$) and $g^2=C(2\theta)$ has $\chi(g^2)=1+2\cos2\theta=4c^2-1$.
Then $\tfrac12[(1+2c)^2-(4c^2-1)]=1+2c=\chiV=\chiR$. Rotoreflection: $\chi=-1+2c$ and
$S(\theta)^2=C(2\theta)$, so $\tfrac12[(2c-1)^2-(4c^2-1)]=1-2c=-\chiV=\chiR$. The characters
agree on every element, so $\{\GV^2\}\cong\GR$ by \cref{thm:charequiv}.
\end{solution}

\begin{exercise}[$\mathrm{BF_3}$]\label{ex:4:BF3}
Using \cref{ex:2:BF3}, decide which vibrational frequencies of $\mathrm{BF_3}$ are
infrared active and which are Raman active. Is the molecule microwave active? Does it
show a rotational Raman spectrum?
\end{exercise}

\begin{solution}
The squares in $\Dh{3}$: $C_2'^2=\sh^2=\sv^2=E$, $C_3^2\in\{2C_3\}$, $S_3^2\in\{2C_3\}$. So
$\chiV(g^2)=(3,0,3,3,0,3)$, $\chiV^2=(9,0,1,1,4,1)$ and $\chi_{[\GV^2]}=(6,0,2,2,2,2)$,
giving $[\GV^2]=2\,\irr{A}_1'\oplus\irr{E}'\oplus\irr{E}''$. With
$\Gvib=\irr{A}_1'\oplus2\,\irr{E}'\oplus\irr{A}_2''$ and $\GV=\irr{E}'\oplus\irr{A}_2''$:
\begin{itemize}
\item $\irr{A}_1'$ (breathing): Raman only;
\item $\irr{A}_2''$ (out-of-plane): infrared only;
\item the two $\irr{E}'$ modes: infrared and Raman.
\end{itemize}
Microwave inactive ($\irr{A}_1'\not\subset\GV$: no permanent dipole). Rotational Raman
active ($\irr{A}_1'$ appears twice in $[\GV^2]$: anisotropic polarizability).
\end{solution}

\begin{exercise}[$\mathrm{H_2O}$]\label{ex:4:H2O}
With the axes of \cref{ex:2:H2O}, find the irreps of $x$, $y$, $z$ and of the six
quadratic functions. Which vibrations of water are infrared active, which are Raman
active?
\end{exercise}

\begin{solution}
$z\in\irr{A}_1$, $x\in\irr{B}_1$, $y\in\irr{B}_2$; $x^2,y^2,z^2\in\irr{A}_1$, $xy\in\irr{A}_2$,
$xz\in\irr{B}_1$, $yz\in\irr{B}_2$. So $\GV=\irr{A}_1\oplus\irr{B}_1\oplus\irr{B}_2$ and
$[\GV^2]=3\,\irr{A}_1\oplus\irr{A}_2\oplus\irr{B}_1\oplus\irr{B}_2$. All three modes
($2\,\irr{A}_1\oplus\irr{B}_2$) are both infrared and Raman active.
\end{solution}

\begin{exercise}[mutual exclusion]\label{ex:4:exclusion}
Suppose the point group contains the inversion $I$. Irreps are then even ($\mathrm{g}$,
$\chi(I)=+d$) or odd ($\mathrm{u}$, $\chi(I)=-d$). Show that $\GV$ contains only odd
irreps and $[\GV^2]$ only even ones. Conclude that no vibration of a centrosymmetric
molecule is both infrared and Raman active.
\end{exercise}

\begin{solution}
$\mat{V}(I)=-\one$, so every irrep contained in $\GV$ has $\mat{D}(I)=-\one$: odd. On symmetric
tensors $I$ acts as $(-1)(-1)=+1$, so every irrep in $[\GV^2]$ is even. An infrared-active
mode must be odd, a Raman-active one even; no mode is both. ($\mathrm{CO_2}$ is the textbook
example.)
\end{solution}

\begin{exercise}[electronic transitions]\label{ex:4:dipole}
An electronic state of a $\Cv{3}$ molecule has symmetry $\irr{A}_1$. To which
states ($\irr{A}_1$, $\irr{A}_2$, $\irr{E}$) can it be excited by electric-dipole
absorption, and with which light polarisation? Answer the same question for an
initial state of symmetry $\irr{E}$.
\end{exercise}

\begin{solution}
From $\irr{A}_1$: $\Gamma_f\subset\GV\otimes\irr{A}_1=\GV=\irr{A}_1\oplus\irr{E}$. Transitions to
$\irr{A}_1$ need $z$-polarised light, to $\irr{E}$ in-plane polarisation; $\irr{A}_1\to\irr{A}_2$ is
forbidden. From $\irr{E}$: $z$ polarisation gives $\irr{A}_1\otimes\irr{E}=\irr{E}$ (only
$\irr{E}\to\irr{E}$), in-plane polarisation gives $\irr{E}\otimes\irr{E}=\irr{A}_1\oplus\irr{A}_2\oplus\irr{E}$
(all final states allowed).
\end{solution}

\chapter{Spin, double groups and crystal-field splitting}\label{ch:spin}

\begin{watch}{\lec{5} (no captions; times read off the slides)}
Subduced representations, $\mathrm{CH_4}\to\mathrm{CH_3D}$: first part of the lecture $\cdot$
irreps of $\SO{3}$ \ts{5}{2940}{49:00} $\cdot$
splitting of atomic orbitals, weak spin--orbit \ts{5}{3790}{63:10} $\cdot$
splitting of multiplets, strong spin--orbit \ts{5}{4490}{74:50} $\cdot$
double groups \ts{5}{4910}{81:50} $\cdot$
reality of a representation \ts{5}{5620}{93:40}
\end{watch}

\section{Subgroups and subduction}\label{sec:subduction}

Deuterating one hydrogen of methane lowers the symmetry from $\Td$ to its subgroup
$\Cv{3}$. A representation $\Gamma$ of $G$ automatically gives a representation of
any subgroup $H\subset G$: just keep the matrices of the elements of $H$. This is the
\emph{subduced} (or restricted) representation $\Gamma\!\downarrow\!H$. Its character is
the restriction of $\chi_\Gamma$ to $H$. An irrep of $G$ is in general reducible
over $H$, and its decomposition is given by the reduction formula in $H$.

Take $\Cv{3}\subset\Td$ with the three-fold axis along $[111]$ and one of the six
mirrors of $\Td$ containing that axis. The classes $E$, $2C_3$, $3\sd$ of $\Cv{3}$ lie
in the classes $E$, $8C_3$, $6\sd$ of $\Td$. Reading the columns of \eqref{eq:Tdchar}:
\begin{equation}\label{eq:TdC3v}
\begin{array}{c|ccc|c}
\Td\downarrow\Cv{3} & E & 2C_3 & 3\sd & \\\hline
\irr{A}_1 & 1 & 1 & 1 & \irr{A}_1\\
\irr{A}_2 & 1 & 1 & -1 & \irr{A}_2\\
\irr{E} & 2 & -1 & 0 & \irr{E}\\
\irr{T}_1 & 3 & 0 & -1 & \irr{A}_2\oplus\irr{E}\\
\irr{T}_2 & 3 & 0 & 1 & \irr{A}_1\oplus\irr{E}
\end{array}
\end{equation}
Such tables are called \emph{correlation} or \emph{subduction} tables. Applied to methane:
\[
\Gvib(\mathrm{CH_4})\!\downarrow\!\Cv{3}=\irr{A}_1\oplus\irr{E}\oplus2(\irr{A}_1\oplus\irr{E})
=3\,\irr{A}_1\oplus3\,\irr{E}=\Gvib(\mathrm{CH_3D}).
\]
Each triply degenerate frequency splits into one non-degenerate and one doubly
degenerate frequency, with no new computation. The general lesson: \emph{lowering the
symmetry can only split levels}, and subduction tells how.

\section{The rotation group}

\subsection{Irreps of \texorpdfstring{$\SO{3}$}{SO(3)}}
The spherical harmonics $Y_l^m$, $m=-l,\dots,l$, span a $(2l+1)$-dimensional irrep
$D_l$ of $\SO{3}$ ($l=0,1,2,\dots$). Every rotation is conjugate to a rotation by the
same angle about $z$, so the character depends only on the angle $\theta$. Using
$\hat{L}_z\ket{l,m}=\hbar m\ket{l,m}$,
\begin{equation}\label{eq:chil}
\chi_l(\theta)=\sum_{m=-l}^{l}\ee^{-\ii m\theta}
=\frac{\ee^{\ii(l+1/2)\theta}-\ee^{-\ii(l+1/2)\theta}}{\ee^{\ii\theta/2}-\ee^{-\ii\theta/2}}
=\frac{\sin\bigl[(l+\tfrac12)\theta\bigr]}{\sin(\theta/2)},
\end{equation}
with $\chi_l(0)=2l+1$. (The sign convention of the exponent does not matter: the sum
is symmetric in $m\to-m$.) Useful values: $\chi_l(\cpi)=(-1)^l$;
$\chi_l(2\cpi/3)=1,0,-1$ for $l\equiv0,1,2\pmod3$.

\subsection{\texorpdfstring{$\Og{3}$}{O(3)} and parity}
The full orthogonal group is $\Og{3}=\SO{3}\times\{E,I\}$, and the inversion commutes
with everything. Each irrep of $\SO{3}$ gives two irreps of $\Og{3}$, $D_l^{+}$ and
$D_l^{-}$, with $\mat{D}(I)=\pm\one$. An improper element $I\,C(\theta)$ has character
$\pm\chi_l(\theta)$. Atomic orbitals have the parity of $Y_l^m$, $(-1)^l$:
$s=D_0^+$, $p=D_1^-$, $d=D_2^+$, $f=D_3^-$. A mirror is $\sigma=IC_2$, so
$\chi_{D_l^{(-1)^l}}(\sigma)=(-1)^l\chi_l(\cpi)=1$ for every orbital.

\subsection{Direct products\beyondmark}\label{sec:directproduct}
$\Og{3}=\SO{3}\times\{E,I\}$ is an instance of a general construction. If $G$ has two
normal subgroups (\cref{sec:normal}) $H$, $K$ with $H\cap K=\{e\}$ and $HK=G$, then
$G\cong H\times K$ and elements of $H$ commute with those of $K$. The irreps of
$H\times K$ are products of irreps of the factors, $\Gamma^{(\mu)}_H\otimes
\Gamma^{(\nu)}_K$, with characters $\chi^{(\mu)}(h)\chi^{(\nu)}(k)$. Other examples:
$\Oh=O\times\Ci$ (\cref{ex:5:Oh}) and $\Dh{6}=\Cv{6}\times\{E,\sh\}$ (\cref{ch:bands}).
Tables for such groups are often given only for one factor; the irreps of the full
group are the $\mathrm{g}$ (even) and $\mathrm{u}$ (odd) copies.

\section{Crystal-field splitting (weak spin--orbit coupling)}\label{sec:crystalfield}

An atom in a molecule or crystal no longer feels a spherically symmetric potential,
only one invariant under its \emph{local} point group $G$ (its site-symmetry group,
\cref{sec:sitesym}). If spin--orbit coupling is weak, an $l$ level splits according
to $D_l^{\pm}\!\downarrow\!G$. For $\Cv{3}$:
\begin{equation}\label{eq:C3vsplit}
\begin{array}{c|ccc|l}
\Cv{3} & E & 2C_3 & 3\sd & \\\hline
D_0^+ & 1 & 1 & 1 & s\to\irr{a}_1\\
D_1^- & 3 & 0 & 1 & p\to\irr{a}_1\oplus\irr{e}\\
D_2^+ & 5 & -1 & 1 & d\to\irr{a}_1\oplus2\,\irr{e}\\
D_3^- & 7 & 1 & 1 & f\to2\,\irr{a}_1\oplus\irr{a}_2\oplus2\,\irr{e}
\end{array}
\end{equation}
Lower-case labels are customary for one-electron orbitals. The most famous example is
a $d$ ion in an octahedral environment ($\Oh$): $d\to\irr{e}_{\mathrm{g}}\oplus
\irr{t}_{2\mathrm{g}}$, the splitting $\Delta_{\mathrm{o}}$ of ligand-field theory
(\cref{ex:5:Oh}). In a tetrahedral environment $d\to\irr{e}\oplus\irr{t}_2$.

\paragraph{Orbital mixing.} Once the symmetry is lowered, orbitals from different
$l$ that belong to the same irrep of $G$ may mix. In $\Cv{3}$, $s$, $p_z$ and $d_{z^2}$
are all $\irr{a}_1$. Strictly speaking there is no ``$p_z$ orbital'' in a crystal, only
an $\irr{a}_1$ orbital that resembles one. This matters when one builds tight-binding
models (\cref{sec:sitesym}).

\section{Half-integer angular momentum}\label{sec:doublevalued}

With spin, the relevant angular momenta $j$ can be half-integers. Formula
\eqref{eq:chil} still gives the character of $D_j$ with $l\to j$, but for half-integer
$j$ it has a disturbing property:
\[
\chi_j(\theta+2\cpi)=-\chi_j(\theta).
\]
A rotation by $2\cpi$ multiplies a spinor by $-1$. These are the irreps of $\SU{2}$,
which covers $\SO{3}$ twice (\cref{sec:su2}); as ``representations'' of $\SO{3}$ they
are \emph{double valued}.

Try anyway to split a $j=\tfrac12$ or $j=\tfrac32$ level in a $\Cv{3}$ environment
(strong spin--orbit coupling, so that $j$ is a good quantum number of the free atom).
For improper elements, the inversion acts trivially on spin and
$\sigma=IC_2$, so $\chi(\sigma)=\pm\chi_j(\cpi)=0$ for half-integer $j$:
\[
\begin{array}{c|ccc}
\Cv{3} & E & 2C_3 & 3\sd\\\hline
D_{1/2}^- & 2 & 1 & 0\\
D_{3/2}^- & 4 & -1 & 0
\end{array}
\qquad
\begin{aligned}
D_{1/2}^-:&\ m_{\irr{A}_1}=m_{\irr{A}_2}=\tfrac23,\ m_{\irr{E}}=\tfrac13 \quad(!)\\
D_{3/2}^-:&\ m_{\irr{A}_1}=m_{\irr{A}_2}=\tfrac13,\ m_{\irr{E}}=\tfrac53 \quad(!)
\end{aligned}
\]
Fractional multiplicities: the characters of spinors are not combinations of the
characters of $\Cv{3}$. The group is too small.

\section{The rotation group and \texorpdfstring{$\SU{2}$}{SU(2)}\beyondmark}\label{sec:su2}

The lecture only needs the characters of the rotation group. A theorist also needs the
relation between $\SO{3}$ and $\SU{2}$, the Clebsch--Gordan series and the
Wigner--Eckart theorem; they are summarised here.

\subsection{Rotations}
A rotation $R(\vec{n},\theta)$ by the angle $\theta\in[0,\cpi]$ about the unit axis $\vec{n}$
acts on vectors as
$R\vec{v}=\cos\theta\,\vec{v}+\sin\theta\,\vec{n}\times\vec{v}+(1-\cos\theta)(\vec{n}\cdot\vec{v})\vec{n}$.
Since $R'R(\vec{n},\theta)R'^{-1}=R(R'\vec{n},\theta)$, all rotations by the same angle
are conjugate: the classes of $\SO{3}$ are labelled by $\theta$, and characters are
functions of $\theta$ only. On states, rotations act as
$\hat{U}(\vec{n},\theta)=\exp(-\ii\theta\,\vec{n}\cdot\hat{\vec{J}}/\hbar)$, and the
generators satisfy $[\hat{J}_x,\hat{J}_y]=\ii\hbar\hat{J}_z$ and cyclic permutations.

\subsection{\texorpdfstring{$\SU{2}$}{SU(2)} covers \texorpdfstring{$\SO{3}$}{SO(3)} twice}
Every $\mat{U}\in\SU{2}$ can be written
\begin{equation}
\mat{U}(\vec{n},\theta)=\exp\bigl(-\tfrac{\ii}{2}\theta\,\vec{n}\cdot\vec{\sigma}\bigr)
=\cos\tfrac{\theta}{2}\,\one-\ii\sin\tfrac{\theta}{2}\,\vec{n}\cdot\vec{\sigma},
\qquad 0\le\theta\le2\cpi .
\end{equation}
Acting by conjugation on the traceless hermitian matrices $\vec{v}\cdot\vec{\sigma}$
(which form a copy of $\mathbb{R}^3$), $\mat{U}(\vec{v}\cdot\vec{\sigma})\mat{U}\herm=
(R\vec{v})\cdot\vec{\sigma}$ defines a rotation $R(\mat{U})$, with
$R_{ij}=\tfrac12\tr(\sigma_i\mat{U}\sigma_j\mat{U}\herm)$. The map $\mat{U}\mapsto R(\mat{U})$ is a
homomorphism onto $\SO{3}$ with kernel $\{\one,-\one\}$ (\cref{sec:normal}):
$\mat{U}(\vec{n},\theta)$ and $\mat{U}(\vec{n},\theta+2\cpi)=-\mat{U}(\vec{n},\theta)$ give
the same rotation. In particular
\begin{equation}
\mat{U}(\vec{n},2\cpi)=-\one:
\end{equation}
a rotation by $2\cpi$ multiplies a spinor by $-1$, and only a rotation by $4\cpi$ is the
identity. Topologically, $\SU{2}$ is the 3-sphere (simply connected), and
$\SO{3}=\SU{2}/\{\pm\one\}$ is doubly connected. A closed path of rotations from $0$ to
$2\cpi$ cannot be shrunk to a point; the same path traversed twice can. This is the
``belt trick''.

\subsection{Irreps and characters}
The irreps of $\SU{2}$ are $D_j$, $j=0,\tfrac12,1,\tfrac32,\dots$, of dimension $2j+1$,
with basis $\ket{j,m}$ and $D_j(\mat{U}(\vec{e}_z,\theta))\ket{j,m}=\ee^{-\ii m\theta}\ket{j,m}$.
Their characters are \eqref{eq:chil} with $l\to j$, and
\begin{equation}
\chi_j(\theta+2\cpi)=(-1)^{2j}\chi_j(\theta).
\end{equation}
For integer $j$, $D_j(-\one)=+\one$ and $D_j$ is a genuine representation of $\SO{3}$
(orbital angular momentum, $j=l$). For half-integer $j$, $D_j(-\one)=-\one$: a
representation of $\SU{2}$ but only a ``double-valued'' one of $\SO{3}$.

\subsection{Clebsch--Gordan series}
The product $D_{j_1}\otimes D_{j_2}$ has eigenvalues
$\ee^{-\ii(m_1+m_2)\theta}$ under $z$ rotations. Collecting the values of $m_1+m_2$ gives
\begin{equation}
\chi_{j_1}\chi_{j_2}=\sum_{j=\abs{j_1-j_2}}^{j_1+j_2}\chi_j,\qquad
D_{j_1}\otimes D_{j_2}=\bigoplus_{j=\abs{j_1-j_2}}^{j_1+j_2}D_j .
\end{equation}
For example, a $p$ electron with spin: $D_1\otimes D_{1/2}=D_{3/2}\oplus D_{1/2}$. This is
the origin of the $j=\tfrac32$, $j=\tfrac12$ multiplets split by spin--orbit coupling. The
change of basis from $\ket{j_1m_1}\ket{j_2m_2}$ to $\ket{jm}$ is given by the
Clebsch--Gordan coefficients $\braket{j_1m_1;j_2m_2}{jm}$.

\subsection{The Wigner--Eckart theorem}\label{sec:wignereckart}
A set of operators $T^{(k)}_q$, $q=-k,\dots,k$, transforming like $\ket{k,q}$ under
rotations is an \emph{irreducible tensor operator} of rank $k$. Then
\begin{equation}
\bra{j,m}T^{(k)}_q\ket{j',m'}=\braket{j'm';kq}{jm}\;
\frac{\langle j\Vert T^{(k)}\Vert j'\rangle}{\sqrt{2j+1}} .
\end{equation}
All $(2j+1)(2k+1)(2j'+1)$ matrix elements are fixed by a single reduced matrix element.
The Clebsch--Gordan coefficient gives the selection rules $m=m'+q$ and
$\abs{j-j'}\le k\le j+j'$. The dipole rule $\Delta l=\pm1$ follows from $k=1$ together
with parity.

For a finite group the statement is the same: matrix elements
$\bra{\Gamma_f,i}\hat{O}^{\Gamma_O}_a\ket{\Gamma_i,j}$ are products of coupling coefficients
of $G$ and reduced matrix elements, one reduced element for each copy of $\Gamma_f$ in
$\Gamma_O\otimes\Gamma_i$ (\cref{sec:raman}).

\section{Double groups}\label{sec:doublegroups}

The cure is to treat the rotation by $2\cpi$ as a group element of its own,
$\Ebar\equiv C(2\cpi)\neq E$, with $\Ebar^2=E$. The \emph{double group} $G^{\mathrm{D}}$
contains, for each $g\in G$, the two elements $g$ and $\Ebar g$, so
$\abs{G^{\mathrm{D}}}=2\abs{G}$. Concretely, $G^{\mathrm{D}}$ is the set of $\SU{2}$
matrices (times the trivial action of $I$) that project onto $G$. In it:
\begin{itemize}
\item $C_n^n=\Ebar$ rather than $E$ (a rotation by $2\cpi$);
\item $\sigma^2=\Ebar$, since $\sigma=IC_2$ and $C_2^2=\Ebar$;
\item $g$ and $\Ebar g$ usually fall in different classes (the precise rules, due to
      Opechowski, are in Bradley--Cracknell).
\end{itemize}
Its irreps are of two kinds: the \emph{single-valued} irreps of $G$ (with
$\mat{D}(\Ebar)=+\one$) and new \emph{double-valued} irreps with $\mat{D}(\Ebar)=-\one$,
which describe systems with an odd number of electrons (half-integer spin).

\begin{example}[The double group of $\Cv{3}$]
Order 12, six classes, six irreps ($1+1+4+1+1+4=12$):
\begin{equation}\label{eq:C3vdouble}
\begin{array}{c|cccccc|c}
\Cv{3}^{\mathrm{D}} & E & \Ebar & 2C_3 & 2\bar{C}_3 & 3\sd & 3\bar{\sigma}_{\mathrm{d}} & \text{type}\\\hline
\irr{A}_1 & 1 & 1 & 1 & 1 & 1 & 1 & 1\\
\irr{A}_2 & 1 & 1 & 1 & 1 & -1 & -1 & 1\\
\irr{E} & 2 & 2 & -1 & -1 & 0 & 0 & 1\\
{}^{1}\bar{\irr{E}} & 1 & -1 & -1 & 1 & \ii & -\ii & 2\\
{}^{2}\bar{\irr{E}} & 1 & -1 & -1 & 1 & -\ii & \ii & 2\\
\bar{\irr{E}}_1 & 2 & -2 & 1 & -1 & 0 & 0 & 3\\\hline
D_{1/2} & 2 & -2 & 1 & -1 & 0 & 0 & \\
D_{3/2} & 4 & -4 & -1 & 1 & 0 & 0 &
\end{array}
\end{equation}
Now the reduction works:
$D_{1/2}\!\downarrow\!\Cv{3}^{\mathrm{D}}=\bar{\irr{E}}_1$ and
$D_{3/2}\!\downarrow\!\Cv{3}^{\mathrm{D}}={}^{1}\bar{\irr{E}}\oplus{}^{2}\bar{\irr{E}}\oplus\bar{\irr{E}}_1$.
The last column is explained in the next section.
\end{example}

The one-dimensional characters are forced by the group structure: in the double group
$C_3^3=\Ebar$, so $\chi(C_3)^3=\chi(\Ebar)=-1$; and $\sigma^2=\Ebar$ gives
$\chi(\sigma)^2=-1$, $\chi(\sigma)=\pm\ii$.

\subsection{Construction, classes and irreps\beyondmark}\label{sec:doublegroupsformal}
For a point group $G\subset\Og{3}$, the double group $G^{\mathrm{D}}$ is the set of
pairs $(g,\pm\mat{U}_g)$, where $\mat{U}_g\in\SU{2}$ is a lift of the rotational part of $g$
(\cref{sec:su2}). The inversion acts trivially on spin, so a mirror $\sigma$ with normal
$\vec{n}$ is represented by the lift of $C_2(\vec{n})$, $-\ii\,\vec{n}\cdot\vec{\sigma}$, and
$\sigma^2=\Ebar$. The order is $2\abs{G}$, and $\Ebar=(E,-\one)$ is central. Useful facts:
\begin{itemize}
\item \emph{Classes (Opechowski's rules).} $g$ and $\Ebar g$ always lie in different
      classes, except when $g$ is a two-fold rotation or a reflection and $G$ contains a
      two-fold rotation or a reflection whose axis (or normal) is perpendicular to that of
      $g$. For example, in $\Cv{3}$ the mirrors $\sd$ and $\bar{\sigma}_{\mathrm{d}}$ are in
      different classes, while in $\Cv{2}$ and $\Cv{6}$ they are not.
\item \emph{Irreps.} Those with $\mat{D}(\Ebar)=+\one$ are the ordinary irreps of $G$. The
      others, with $\mat{D}(\Ebar)=-\one$, are the double-valued irreps, and
      $\sum_{\text{double-valued}}d^2=\abs{G}$.
\item \emph{Electrons.} With spin--orbit coupling the states of an odd number of electrons
      carry double-valued irreps only. For weak spin--orbit coupling, obtain them as
      $\Gamma_{\mathrm{orb}}\otimes D_{1/2}$ (\cref{ex:5:SO}).
\end{itemize}

\section{Reality of a representation}\label{sec:reality}

Given a representation $\mat{D}$, the complex conjugate matrices $\mat{D}^{*}$ also form
a representation, with character $\chi^{*}$. For an irrep there are three
possibilities:
\begin{enumerate}
\item \emph{real} (type 1): $\mat{D}$ is equivalent to a representation by real
      matrices;
\item \emph{complex} (type 2): $\mat{D}^{*}$ is not equivalent to $\mat{D}$
      ($\chi$ is not real);
\item \emph{pseudoreal} (type 3): $\mat{D}^{*}$ is equivalent to $\mat{D}$ ($\chi$ is real),
      yet no basis makes all matrices real.
\end{enumerate}
The \emph{Frobenius--Schur indicator} distinguishes them from the character alone:
\begin{equation}\label{eq:FS}
\frac{1}{\abs{G}}\sum_{g\in G}\chi(g^2)=
\begin{cases}+1&\text{real},\\ 0&\text{complex},\\ -1&\text{pseudoreal.}\end{cases}
\end{equation}
In the double group of $\Cv{3}$: $\irr{A}_1,\irr{A}_2,\irr{E}$ are real; ${}^{1}\bar{\irr{E}}$
and ${}^{2}\bar{\irr{E}}$ are complex and conjugate to each other; $\bar{\irr{E}}_1$ is
pseudoreal (\cref{ex:5:FS}). The spin-$\tfrac12$ representation of $\SU{2}$ itself is the
prototypical pseudoreal representation: $\sigma_y\mat{U}^{*}\sigma_y=\mat{U}$ for every
$\mat{U}\in\SU{2}$.

Why this matters: time reversal is antiunitary and effectively relates a
representation to its complex conjugate. It glues complex-conjugate pairs of irreps
into degenerate \emph{physically irreducible} representations. For
${}^{1}\bar{\irr{E}}\oplus{}^{2}\bar{\irr{E}}$ this gives the Kramers degeneracy of the
$j=\tfrac32$ level, which in a $\Cv{3}$ field splits only into two Kramers doublets. It
also decides when the bilinear invariants of \cref{sec:invariants} come in pairs
(\cref{sec:complexinv}). See \cref{sec:TR}.

\section{Time reversal\beyondmark}\label{sec:TR}

Time reversal appears in the lectures in two places: in the pairing of
complex-conjugate irreps (\cref{sec:reality}, and \cref{sec:complexinv} in the next
chapter) and in band topology (\cref{ch:topology}). This section gives the minimum
theory for an atom or a molecule; \cref{sec:TRcrystal} adds what changes in a crystal.

\subsection{An antiunitary symmetry}
By Wigner's theorem on symmetries, a symmetry of quantum mechanics is represented by an
operator that is either unitary or \emph{antiunitary}: antilinear,
$\TR(a\psi+b\phi)=a^{*}\TR\psi+b^{*}\TR\phi$, with
$\braket{\TR\psi}{\TR\phi}=\braket{\psi}{\phi}^{*}$. Time reversal must reverse momenta
and angular momenta but not positions, and it must preserve $[\hat{x},\hat{p}]=\ii\hbar$.
This forces antiunitarity. Every antiunitary operator can be written $\TR=\hat{U}K$ with
$\hat{U}$ unitary and $K$ complex conjugation in a fixed basis.
\begin{itemize}
\item \emph{Spinless particles}: $\TR=K$ in the position representation,
      $(\TR\psi)(\vec{r})=\psi(\vec{r})^{*}$, and $\TR^2=+1$.
\item \emph{Spin $\tfrac12$}: $\TR=-\ii\sigma_yK$, which reverses $\vec{\sigma}$
      ($\TR\vec{\sigma}\TR^{-1}=-\vec{\sigma}$), and $\TR^2=-1$.
\item $N$ electrons: $\TR^2=(-1)^N$.
\end{itemize}

\subsection{Kramers' theorem}\label{sec:kramers}
\begin{theorem}[Kramers]
If $[\hat{H},\TR]=0$ and $\TR^2=-1$, every energy level is at least doubly degenerate.
\end{theorem}
\begin{proof}
$\TR\psi$ has the same energy as $\psi$. If they were the same state,
$\TR\psi=c\psi$, then $\TR^2\psi=\TR(c\psi)=c^{*}\TR\psi=\abs{c}^2\psi=\psi$,
contradicting $\TR^2=-1$. In fact the two states are orthogonal: antiunitarity gives
$\braket{a}{b}=\braket{\TR b}{\TR a}$, and with $a=\psi$, $b=\TR\psi$,
$\braket{\psi}{\TR\psi}=\braket{\TR^2\psi}{\TR\psi}=-\braket{\psi}{\TR\psi}$, so
$\braket{\psi}{\TR\psi}=0$.
\end{proof}

\subsection{Extra degeneracies}\label{sec:herring}
Time reversal commutes with the point-group operations, so it maps the states of a
level onto states of the same energy that transform under the complex-conjugate
representation. It can therefore force two irreps to stick together, or double a
single one. For a point group $G$ the answer depends on the reality type of the irrep
(\cref{sec:reality}), measured by the Frobenius--Schur indicator
$s=\frac{1}{\abs{G}}\sum_g\chi(g^2)$, and on $\TR^2$:
\begin{equation}
\begin{array}{lcc}
\toprule
\text{reality type of }\Gamma & \TR^2=+1\ \text{(no spin)} & \TR^2=-1\ \text{(spin }\tfrac12)\\
\midrule
\text{real } (s=+1) & \text{no extra degeneracy} & \text{doubled: } \Gamma\oplus\Gamma\\
\text{complex } (s=0) & \Gamma\oplus\Gamma^{*}\ \text{degenerate} & \Gamma\oplus\Gamma^{*}\ \text{degenerate}\\
\text{pseudoreal } (s=-1) & \text{doubled: } \Gamma\oplus\Gamma & \text{no extra degeneracy}\\
\bottomrule
\end{array}
\end{equation}
Examples. For spinless electrons in $C_4$ (\cref{sec:complexinv}), the complex pair
${}^{1}\irr{E}$, ${}^{2}\irr{E}$ becomes a degenerate doublet, the physically
irreducible representation $\irr{E}$. For spin-$\tfrac12$ electrons in $\Cv{3}$, the
complex pair ${}^{1}\bar{\irr{E}}$, ${}^{2}\bar{\irr{E}}$ becomes a Kramers doublet, while
the pseudoreal $\bar{\irr{E}}_1$ is already a doublet and is not doubled again. So a
$j=\tfrac32$ level at a $\Cv{3}$ site splits into two Kramers doublets.

\exercisesection

\begin{exercise}\label{ex:5:Cs}
Let $\Cs=\{E,\sigma\}$ be the subgroup of $\Cv{3}$ generated by one mirror, with irreps
$\irr{A}'$ (even) and $\irr{A}''$ (odd). Construct the subduction table
$\Cv{3}\to\Cs$.
\end{exercise}

\begin{solution}
Keep the columns $E$ and $\sigma$: $\irr{A}_1\to(1,1)=\irr{A}'$,
$\irr{A}_2\to(1,-1)=\irr{A}''$, $\irr{E}\to(2,0)=\irr{A}'\oplus\irr{A}''$.
\end{solution}

\begin{exercise}\label{ex:5:Td}
Show that in a tetrahedral environment $p\to\irr{t}_2$, $d\to\irr{e}\oplus\irr{t}_2$
and $f\to\irr{a}_1\oplus\irr{t}_1\oplus\irr{t}_2$. Remember that $S_4=\sh C_4$ is
improper and use $\chi(IC(\theta))=\pm\chi_l(\theta)$.
\end{exercise}

\begin{solution}
Write $S_4=\sh C_4=IC(270^\circ)$ and $\sd=IC_2$. Then
$\chi(S_4)=\pm\chi_l(90^\circ)$ and $\chi(\sd)=\pm\chi_l(180^\circ)$, with the parity sign
$(-1)^l$. Using \eqref{eq:chil}:
\[
\begin{array}{c|ccccc|l}
 & E & 8C_3 & 3C_2 & 6S_4 & 6\sd & \\\hline
p\ (l=1) & 3 & 0 & -1 & -1 & 1 & \irr{T}_2\\
d\ (l=2) & 5 & -1 & 1 & -1 & 1 & \irr{E}\oplus\irr{T}_2\\
f\ (l=3) & 7 & 1 & -1 & 1 & 1 & \irr{A}_1\oplus\irr{T}_1\oplus\irr{T}_2
\end{array}
\]
For instance for $d$: $m_{\irr{E}}=\tfrac1{24}(10+8+6+0+0)=1$ and
$m_{\irr{T}_2}=\tfrac1{24}(15+0-3+6+6)=1$.
\end{solution}

\begin{exercise}[octahedral field]\label{ex:5:Oh}
Since $d$ orbitals are even under inversion, their splitting in $\Oh=O\times\Ci$ is
decided by the rotation group $O$, with classes $E$, $8C_3$, $3C_2$, $6C_4$, $6C_2'$
and character table
\[
\begin{array}{c|ccccc}
O & E & 8C_3 & 3C_2 & 6C_4 & 6C_2'\\\hline
\irr{A}_1&1&1&1&1&1\\ \irr{A}_2&1&1&1&-1&-1\\ \irr{E}&2&-1&2&0&0\\
\irr{T}_1&3&0&-1&1&-1\\ \irr{T}_2&3&0&-1&-1&1
\end{array}
\]
Show that $d\to\irr{e}\oplus\irr{t}_2$ (in $\Oh$: $\irr{e}_{\mathrm{g}}\oplus\irr{t}_{2\mathrm{g}}$).
Which of the two is spanned by $d_{xy},d_{yz},d_{zx}$?
\end{exercise}

\begin{solution}
$\chi_2$: $E\to5$, $C_3\to-1$, $C_2\to1$, $C_4\to\sin(5\cpi/4)/\sin(\cpi/4)=-1$,
$C_2'\to1$. Then $m_{\irr{E}}=\tfrac1{24}(10+8+6)=1$, $m_{\irr{T}_2}=\tfrac1{24}(15-3+6+6)=1$,
all others $0$: $d\to\irr{e}\oplus\irr{t}_2$, which are $\irr{e}_{\mathrm{g}}\oplus\irr{t}_{2\mathrm{g}}$
in $\Oh$. The three orbitals $xy,yz,zx$ are permuted among themselves by the cubic
rotations and span $\irr{t}_{2\mathrm{g}}$; $x^2-y^2$ and $3z^2-r^2$ span $\irr{e}_{\mathrm{g}}$.
\end{solution}

\begin{exercise}[weak vs strong spin--orbit]\label{ex:5:SO}
A $p$ electron sits at a $\Cv{3}$ site. (a) Weak spin--orbit coupling: split the
orbitals first, $p\to\irr{a}_1\oplus\irr{e}$, then multiply by the spin representation
$D_{1/2}=\bar{\irr{E}}_1$. (b) Strong spin--orbit coupling: split $j=\tfrac12$ and
$j=\tfrac32$ directly. Show that both orders give the same set of double-group irreps,
$2\bar{\irr{E}}_1\oplus{}^{1}\bar{\irr{E}}\oplus{}^{2}\bar{\irr{E}}$.
\end{exercise}

\begin{solution}
(a) $\irr{A}_1\otimes\bar{\irr{E}}_1=\bar{\irr{E}}_1$. The product $\irr{E}\otimes\bar{\irr{E}}_1$
has characters
\[
\bigl(2\cdot2,\ 2\cdot(-2),\ (-1)(1),\ (-1)(-1),\ 0,\ 0\bigr)=(4,-4,-1,1,0,0),
\]
equal to ${}^{1}\bar{\irr{E}}\oplus{}^{2}\bar{\irr{E}}\oplus\bar{\irr{E}}_1$. Total:
$2\bar{\irr{E}}_1\oplus{}^{1}\bar{\irr{E}}\oplus{}^{2}\bar{\irr{E}}$.
(b) $D_{1/2}\to\bar{\irr{E}}_1$ and $D_{3/2}\to{}^{1}\bar{\irr{E}}\oplus{}^{2}\bar{\irr{E}}\oplus\bar{\irr{E}}_1$:
the same set. The final labels do not depend on the order in which the perturbations
are applied; only the energies do.
\end{solution}

\begin{exercise}\label{ex:5:FS}
Verify the ``type'' column of \eqref{eq:C3vdouble} with the Frobenius--Schur indicator.
You need the classes of the squares: $(C_3)^2\in\{2\bar{C}_3\}$,
$(\bar{C}_3)^2\in\{2\bar{C}_3\}$, $\sigma^2=\bar{\sigma}^2=\Ebar$.
\end{exercise}

\begin{solution}
Sum $\chi(g^2)$ over the 12 elements, grouped by class (sizes $1,1,2,2,3,3$), with
$g^2$ equal to $E$, $E$, $\bar C_3$, $\bar C_3$, $\Ebar$, $\Ebar$ respectively:
\begin{align*}
\irr{E}:&\ 2+2+2(-1)+2(-1)+3\cdot2+3\cdot2=12 &&\Rightarrow s=+1,\\
{}^{1}\bar{\irr{E}}:&\ 1+1+2\cdot1+2\cdot1+3(-1)+3(-1)=0 &&\Rightarrow s=0,\\
\bar{\irr{E}}_1:&\ 2+2+2(-1)+2(-1)+3(-2)+3(-2)=-12 &&\Rightarrow s=-1.
\end{align*}
$\irr{A}_1$ and $\irr{A}_2$ give $s=+1$ trivially.
\end{solution}

\chapter{Space groups and their representations}\label{ch:space}

\begin{watch}{\lec{6}}
Complex irreps and bilinear invariants \ts{6}{3}{0:03} $\cdot$
$\Cv{4}$ versus $C_4$ \ts{6}{314}{5:14} $\cdot$
circular polarisation and time reversal \ts{6}{1081}{18:01} $\cdot$
from atoms to crystals; the bands of graphene \ts{6}{1519}{25:19} $\cdot$
symmetry-enforced vs accidental degeneracies \ts{6}{1775}{29:35} $\cdot$
irreps of the translation group \ts{6}{2724}{45:24} $\cdot$
the Brillouin zone as the set of irreps \ts{6}{3232}{53:52} $\cdot$
Bloch's theorem as Wigner's theorem \ts{6}{3295}{54:55} $\cdot$
Bloch sums from orbitals \ts{6}{3832}{63:52} $\cdot$
space groups, symmorphic and non-symmorphic \ts{6}{4217}{70:17} $\cdot$
star and little group \ts{6}{4661}{77:41} $\cdot$
small representations \ts{6}{5175}{86:15}
\end{watch}

\section{Complex representations and bilinear invariants}\label{sec:complexinv}

So far all irreps were real. When an irrep is complex, the invariant scalar product
of two copies must be the \emph{hermitian} one, $\sum_iz_i^{*}w_i$, because the matrices
are unitary rather than orthogonal. Two consequences:
\begin{itemize}
\item a bilinear invariant pairs an irrep with its complex conjugate, as in
      \eqref{eq:ninv};
\item $z^{*}w$ is complex: its real and imaginary parts are two independent real
      invariants.
\end{itemize}

Compare $\Cv{4}$ with its rotation subgroup $C_4=\{E,C_4^+,C_2,C_4^-\}$ (rotations about
$z$ only). $C_4$ is Abelian, every element is a class, so there are four irreps, all
one-dimensional (the only solution of $\sum d_\mu^2=4$ with four terms):
\begin{equation}\label{eq:C4table}
\begin{array}{c|cccc|l}
C_4 & E & C_4^+ & C_2 & C_4^- & \text{components}\\\hline
\irr{A} & 1&1&1&1 & E_z\\
\irr{B} & 1&-1&1&-1 & \\
{}^{1}\irr{E} & 1&\ii&-1&-\ii & E_+=E_x+\ii E_y\\
{}^{2}\irr{E} & 1&-\ii&-1&\ii & E_-=E_x-\ii E_y
\end{array}
\end{equation}
\begin{quote}\small
\emph{Convention.} Here we list vector \emph{components}: the rotation $C_4^+$ maps
$(E_x,E_y)\mapsto(-E_y,E_x)$, hence $E_+\mapsto\ii E_+$. If instead one lists
\emph{basis functions} with $(\hat{O}_gf)(\vec{r})=f(g^{-1}\vec{r})$, the function $x+\ii y$
picks up $-\ii$. The two conventions are complex conjugates of each other; tables
differ on this point, so check before using one.
\end{quote}
Every irrep of an Abelian group is one-dimensional. This will be crucial for the
translation group.

Physically, $E_\pm$ are circular polarisations. A rotation by $90^\circ$ shifts their
phase by $\pm90^\circ$. Time reversal reverses the sense of rotation and exchanges
them. The pair ${}^{1}\irr{E}\oplus{}^{2}\irr{E}$ is therefore glued by time reversal
into a two-dimensional \emph{physically irreducible} representation, which is why it
carries the letter $\irr{E}$.

\paragraph{Invariants.} In $\Cv{4}$, $\GV=\irr{A}_1\oplus\irr{E}$ as in $\Cv{3}$, and
the bilinear invariants of two vectors $\vec{E}$, $\vec{P}$ are $E_zP_z$ and
$E_xP_x+E_yP_y$. In $C_4$, $\GV=\irr{A}\oplus{}^{1}\irr{E}\oplus{}^{2}\irr{E}$, and
\[
E_+^{*}P_+=(E_xP_x+E_yP_y)+\ii\,(E_xP_y-E_yP_x)
\]
provides, besides $E_zP_z$, \emph{two} invariants: the in-plane scalar product and
$(\vec{E}\times\vec{P})_z$. The second is allowed because $C_4$ has no mirrors.
Groups made only of proper rotations are called \emph{chiral}; they allow such
``vector-product'' couplings, with physical consequences such as optical activity.

\paragraph{Checking symmetry-adapted coordinates.} To verify that given coordinates
belong to an irrep, transform them with the \emph{generators} only (for $C_4$, just
$C_4^+$; for $\Cv{4}$, $C_4^+$ and one mirror); every other element is a product of
generators. For a multidimensional irrep compare with the tabulated matrices, since
the character is not enough.

\section{From molecules to crystals}

A crystal is a ``molecule'' (a set of atoms, the basis) repeated at the points of a
Bravais lattice
\[
\vec{t}=n_1\vec{a}_1+n_2\vec{a}_2+n_3\vec{a}_3,\qquad n_i\in\mathbb{Z}.
\]
For graphene the basis consists of two carbon atoms, A and B, identical but with
different environments (\cref{fig:honeycomb}). The crystal is invariant under all
lattice translations, which form the Abelian \emph{translation group} $\Tgroup$.
Taken literally, a perfect crystal has infinitely many degrees of freedom and its
representations are infinite dimensional. These are the \emph{band representations}
of \cref{ch:topology}. The translation group lets us break them into finite pieces.

\begin{figure}[ht]
\centering
\begin{tikzpicture}[scale=0.95]
  \def\s{1.7320508}
  % hexagon centres at t = n1 a1 + n2 a2 with a1=(1,0), a2=(1/2,sqrt3/2) (scale a=1.6)
  \begin{scope}[scale=1.6]
  \foreach \n in {-2,...,2}{\foreach \m in {-2,...,2}{
    \pgfmathsetmacro{\cx}{\n+0.5*\m}\pgfmathsetmacro{\cy}{0.8660254*\m}
    \pgfmathtruncatemacro{\inside}{(abs(\cx)<2.3) && (abs(\cy)<1.6) ? 1 : 0}
    \ifnum\inside=1
      \foreach \k in {0,...,5}{
        \draw[gray!60] ($(\cx,\cy)+(30+60*\k:0.57735)$)--($(\cx,\cy)+(90+60*\k:0.57735)$);}
      \fill[red!70!black] ($(\cx,\cy)+(30:0.57735)$) circle (0.045);
      \fill[red!70!black] ($(\cx,\cy)+(150:0.57735)$) circle (0.045);
      \fill[red!70!black] ($(\cx,\cy)+(270:0.57735)$) circle (0.045);
      \fill[green!50!black] ($(\cx,\cy)+(90:0.57735)$) circle (0.045);
      \fill[green!50!black] ($(\cx,\cy)+(210:0.57735)$) circle (0.045);
      \fill[green!50!black] ($(\cx,\cy)+(330:0.57735)$) circle (0.045);
    \fi}}
  \draw[-{Stealth},thick] (0,0)--(1,0) node[below right]{$\vec{a}_1$};
  \draw[-{Stealth},thick] (0,0)--(0.5,0.8660254) node[above]{$\vec{a}_2$};
  \fill (0,0) circle (0.035) node[below left,font=\scriptsize]{$O$};
  \node[red!70!black,font=\scriptsize,right] at ($(30:0.57735)+(0.03,0)$) {A};
  \node[green!50!black,font=\scriptsize,above] at ($(90:0.57735)+(0,0.02)$) {B};
  \end{scope}
\end{tikzpicture}
\caption{The honeycomb lattice. The origin $O$ is the centre of a hexagon (Wyckoff
position $1a$, site symmetry $\Cv{6}$). The atoms of sublattice A (red) sit at
$\vec{q}_{\mathrm{A}}=\tfrac13(\vec{a}_1+\vec{a}_2)$ and those of sublattice B (green) at
$\vec{q}_{\mathrm{B}}=\tfrac23(\vec{a}_1+\vec{a}_2)$ plus lattice vectors (Wyckoff position
$2b$, site symmetry $\Cv{3}$). Lattice constant $a=\abs{\vec{a}_1}$.}
\label{fig:honeycomb}
\end{figure}

\section{Irreducible representations of the translation group}

Since $\Tgroup$ is Abelian, its irreps are one-dimensional: to each translation a
unimodular complex number (a phase), with $D(\vec{t}_1)D(\vec{t}_2)=D(\vec{t}_1+\vec{t}_2)$.
The only continuous functions turning sums into products are exponentials, so
\begin{equation}\label{eq:Tirrep}
D^{\vec{k}}(\vec{t})=\ee^{-\ii\vec{k}\cdot\vec{t}}
\end{equation}
for some wave vector $\vec{k}$. (Imposing Born--von K\'arm\'an periodic boundary
conditions on a crystal of $N$ cells makes $\vec{k}$ discrete, with $N$ allowed values;
nothing else changes.)

Introduce the reciprocal basis, $\vec{a}_i\cdot\vec{b}_j=2\cpi\,\delta_{ij}$, and write
$\vec{k}=\sum_ik_i\vec{b}_i$. Then $\vec{k}\cdot\vec{t}=2\cpi\sum_ik_in_i$. If
$\vec{G}=\sum_im_i\vec{b}_i$ with integer $m_i$ (a reciprocal-lattice vector),
$\ee^{-\ii\vec{G}\cdot\vec{t}}=1$ for every $\vec{t}$. Hence $\vec{k}$ and $\vec{k}+\vec{G}$
label the same irrep: they are \emph{equivalent}, written $\vec{k}\equiv\vec{k}+\vec{G}$.

\begin{keyresult}[The Brillouin zone]
The inequivalent irreps of $\Tgroup$ are in one-to-one correspondence with the points
of any primitive cell of the reciprocal lattice. The first Brillouin zone is the
most symmetric such cell. \emph{The Brillouin zone is the set of irreps of the
translation group}, and $\vec{k}$ (``crystal momentum'') is a label of an irrep, just as
$l$ labels an irrep of $\SO{3}$.
\end{keyresult}

\section{Bloch's theorem is Wigner's theorem}

Let the translation operators act on functions as $(\hat{T}_{\vec{t}}\psi)(\vec{r})=
\psi(\vec{r}-\vec{t})$ (active convention: the function is moved by $+\vec{t}$). A
function belongs to the irrep $\vec{k}$ if
$\hat{T}_{\vec{t}}\psi=\ee^{-\ii\vec{k}\cdot\vec{t}}\psi$, i.e.\
$\psi(\vec{r}+\vec{t})=\ee^{\ii\vec{k}\cdot\vec{t}}\psi(\vec{r})$. The general solution is
\begin{equation}
\psi_{\vec{k}}(\vec{r})=\ee^{\ii\vec{k}\cdot\vec{r}}\,u_{\vec{k}}(\vec{r}),\qquad
u_{\vec{k}}(\vec{r}+\vec{t})=u_{\vec{k}}(\vec{r}).
\end{equation}
Since the Hamiltonian of a crystal commutes with $\Tgroup$, its eigenfunctions can be
chosen to belong to irreps of $\Tgroup$ (Wigner). That is all Bloch's theorem says.
(Bradley and Cracknell use the passive convention, with the opposite sign of the
exponent in \eqref{eq:Tirrep}.)

\subsection{Bloch sums}
Place an orbital $\varphi_\alpha$ at position $\vec{q}_\alpha$ of every cell and form
\begin{equation}\label{eq:blochsum}
\Phi_{\alpha\vec{k}}(\vec{r})=\sum_{\vec{t}}\ee^{\ii\vec{k}\cdot\vec{t}}\,
\varphi_\alpha(\vec{r}-\vec{t}-\vec{q}_\alpha).
\end{equation}
Relabelling the sum shows $\hat{T}_{\vec{t}'}\Phi_{\alpha\vec{k}}=\ee^{-\ii\vec{k}\cdot\vec{t}'}
\Phi_{\alpha\vec{k}}$: the Bloch sum belongs to the irrep $\vec{k}$. With this
convention $\Phi_{\alpha,\vec{k}+\vec{G}}=\Phi_{\alpha\vec{k}}$. Bloch sums are
``Bloch-like'' waves: eigenfunctions of translations but not in general of $\hat{H}$.

Because $\hat{H}$ is invariant, it has no matrix elements between inequivalent irreps
of $\Tgroup$ (\cref{sec:countinv}). With $N_{\mathrm{o}}$ orbitals per cell, the
$\infty\times\infty$ tight-binding problem in real space splits into one
$N_{\mathrm{o}}\times N_{\mathrm{o}}$ problem at each $\vec{k}$, giving $N_{\mathrm{o}}$ bands.
Different points of the Brillouin zone do not talk to each other. This is the only
use of symmetry you can make without group theory. To go further we need the rest of
the space group.

\section{Space groups}

A space-group element combines a point operation $R$ (rotation, reflection,
rotoinversion) with a translation $\vec{v}$. In Seitz notation
\begin{equation}
\seitz{R}{\vec{v}}\,\vec{r}=R\vec{r}+\vec{v},
\end{equation}
first $R$, then the translation. From this,
\begin{equation}\label{eq:seitz}
\seitz{R_1}{\vec{v}_1}\seitz{R_2}{\vec{v}_2}=\seitz{R_1R_2}{\vec{v}_1+R_1\vec{v}_2},\qquad
\seitz{R}{\vec{v}}^{-1}=\seitz{R^{-1}}{-R^{-1}\vec{v}}.
\end{equation}
The pure translations $\seitz{E}{\vec{t}}$ form a normal subgroup $\Tgroup$, and the set
of rotational parts forms the \emph{point group} $\bar{G}\cong G/\Tgroup$
(\cref{sec:normal,sec:semidirect}).

\begin{definition}
A space group is \emph{symmorphic} if, for a suitable origin, every element has the
form $\seitz{R}{\vec{t}}$ with $\vec{t}\in\Tgroup$; equivalently, $G$ contains the point
group $\bar{G}$ as a subgroup, $G=\Tgroup\rtimes\bar{G}$. Otherwise it is
\emph{non-symmorphic}: some operations come only with a fractional translation that
no choice of origin removes. These are \emph{glide planes} (reflection plus a half
translation parallel to the plane) and \emph{screw axes} (rotation $C_n$ plus a
translation $\tfrac{p}{n}\vec{t}$ along the axis).
\end{definition}
In three dimensions there are 230 space groups, of which 73 are symmorphic. In two
dimensions there are 17 plane (wallpaper) groups, 13 of them symmorphic; the
honeycomb lattice has the symmorphic plane group $p6mm$. Its three-dimensional
analogue, used in the lecture and in the Bilbao tables, is $P6mm$ (No.~183). For
symmorphic groups the representation theory reduces to that of point groups; for
non-symmorphic groups it is subtler but fully tabulated.

\subsection{Semidirect products\beyondmark}\label{sec:semidirect}
The translation subgroup $\Tgroup$ is normal (\cref{sec:normal}):
$\seitz{R}{\vec{v}}\seitz{E}{\vec{t}}\seitz{R}{\vec{v}}^{-1}=\seitz{E}{R\vec{t}}$, and
$R\vec{t}$ is again a lattice vector (\cref{ex:6:seitz}). The quotient $G/\Tgroup$ is
isomorphic to the point group $\bar{G}$: the coset of $\seitz{R}{\vec{v}}$ is determined
by $R$ alone.

If $G$ has a normal subgroup $N$ and a subgroup $H$ with $N\cap H=\{e\}$ and $NH=G$,
then $G=N\rtimes H$ is a \emph{semidirect product}: $H$ acts on $N$ by conjugation.
(When $H$ is normal too, this is the direct product of \cref{sec:directproduct}.)
Examples: $\Cv{3}=C_3\rtimes\{E,\sd\}$; a symmorphic space group is
$\Tgroup\rtimes\bar{G}$, the point group acting on the lattice by rotation.

Non-symmorphic space groups are \emph{not} semidirect products: $\bar{G}$ is still the
quotient $G/\Tgroup$ but is not a subgroup of $G$, because some rotations only occur
combined with fractional translations. This is the structural reason why their
representation theory needs projective representations (see the box on non-symmorphic
groups in \cref{sec:star}).

\subsection{Space-group operations move \texorpdfstring{$\vec{k}$}{k}}
Let $g=\seitz{R}{\vec{v}}$ act on functions by $(\hat{O}_g\psi)(\vec{r})=\psi(g^{-1}\vec{r})$.
From \eqref{eq:seitz}, $g^{-1}\seitz{E}{\vec{t}}g=\seitz{E}{R^{-1}\vec{t}}$. Hence, if
$\psi$ belongs to $\vec{k}$,
\[
\hat{T}_{\vec{t}}\,\hat{O}_g\psi=\hat{O}_g\hat{T}_{R^{-1}\vec{t}}\psi
=\ee^{-\ii\vec{k}\cdot R^{-1}\vec{t}}\,\hat{O}_g\psi
=\ee^{-\ii(R\vec{k})\cdot\vec{t}}\,\hat{O}_g\psi :
\]
\emph{$\hat{O}_g$ maps the irrep $\vec{k}$ of $\Tgroup$ onto the irrep $R\vec{k}$.}

\section{Star, little group and small representations}\label{sec:star}

\begin{definition}
\begin{itemize}
\item The \emph{star} of $\vec{k}$ is the set of inequivalent vectors
      $\{R\vec{k}\,:\,R\in\bar{G}\}$.
\item The \emph{little group} $G_{\vec{k}}$ is the set of $g=\seitz{R}{\vec{v}}\in G$
      with $R\vec{k}\equiv\vec{k}$, i.e.\ $R\vec{k}=\vec{k}+\vec{G}$ for some reciprocal
      vector $\vec{G}$. Its point parts form the \emph{little co-group}
      $\bar{G}_{\vec{k}}\subset\bar{G}$.
\item A \emph{small representation} is an irrep of $G_{\vec{k}}$ in which the
      translations are represented by $\ee^{-\ii\vec{k}\cdot\vec{t}}\,\one$.
\end{itemize}
\end{definition}
The equivalence ``up to $\vec{G}$'' is essential: for graphene the corner $K$ of the
hexagonal zone is carried by $C_3$ to another corner that differs from it by a
reciprocal-lattice vector, so $C_3$ belongs to the little co-group of $K$
(\cref{ch:bands}).

Why the little group matters: the Bloch states at a given $\vec{k}$ are mapped onto
themselves by $G_{\vec{k}}$, and the $N_{\mathrm{o}}\times N_{\mathrm{o}}$ Hamiltonian
$\mat{H}(\vec{k})$ commutes with the small representation. By Wigner's theorem the
energy levels at $\vec{k}$ are labelled by small irreps, and their degeneracies are the
dimensions of those irreps. For bilinear quantities (Hamiltonians, optical matrix
elements between states at the same $\vec{k}$) only small representations are needed.
Full space-group irreps are needed for electron--phonon processes, superstructures
and similar problems. They are obtained by \emph{induction} from small irreps and
have dimension (size of the star)$\times$(dimension of the small irrep).

\begin{keyresult}[Small irreps of symmorphic groups]
If $G$ is symmorphic, the small irreps at $\vec{k}$ are
\begin{equation}\label{eq:smallsym}
D^{\vec{k},\mu}\bigl(\seitz{R}{\vec{t}}\bigr)=\ee^{-\ii\vec{k}\cdot\vec{t}}\,
\mat{D}^{(\mu)}(R),\qquad R\in\bar{G}_{\vec{k}},\ \vec{t}\in\Tgroup,
\end{equation}
where $\mat{D}^{(\mu)}$ runs over the irreps of the little co-group (a point group).
They are labelled like point-group irreps, or as $\Gamma_1,\Gamma_2,\dots$,
$K_1,K_2,\dots$ in the Bilbao convention.
\end{keyresult}

\begin{beyond}[Beyond the lectures: non-symmorphic groups]
For $g=\seitz{R}{\vec{v}}\in G_{\vec{k}}$ with fractional $\vec{v}$, try
$\mat{\Delta}(g)=\ee^{-\ii\vec{k}\cdot\vec{v}}\mat{D}(R)$. Since
$g_1g_2=\seitz{R_1R_2}{\vec{v}_1+R_1\vec{v}_2}$, the product
$\mat{\Delta}(g_1)\mat{\Delta}(g_2)=\ee^{-\ii\vec{k}\cdot(\vec{v}_1+\vec{v}_2)}\mat{D}(R_1)\mat{D}(R_2)$
must equal $\ee^{-\ii\vec{k}\cdot(\vec{v}_1+R_1\vec{v}_2)}\mat{D}(R_1R_2)$. The two phases
differ by $\ee^{-\ii\vec{k}\cdot\vec{v}_2+\ii\vec{k}\cdot R_1\vec{v}_2}=
\ee^{\ii(R_1^{-1}\vec{k}-\vec{k})\cdot\vec{v}_2}$. If every $R\in\bar{G}_{\vec{k}}$ fixes
$\vec{k}$ exactly (always the case inside the zone), this factor is $1$ and $\mat{D}$
is an ordinary irrep of $\bar{G}_{\vec{k}}$. On the zone boundary,
$R_1^{-1}\vec{k}=\vec{k}+\vec{G}_1$ and the factor $\ee^{\ii\vec{G}_1\cdot\vec{v}_2}$ can
differ from $1$ when $\vec{v}_2$ is fractional. Then $\mat{D}$ must be a
\emph{projective} representation of the little co-group. This is how glides and screws
force bands to stick together on the zone boundary (\cref{ex:6:glide}).
\end{beyond}

\section{Degeneracies: enforced or accidental?}\label{sec:accidental}

Look at a first-principles band structure of graphene (\cref{fig:grapheneTB} shows the
$\uppi$ bands). Its bands carry labels of small irreps at $\Gamma$, $K$, $M$ and along the
lines joining them. Two kinds of degeneracy appear:
\begin{itemize}
\item \emph{Symmetry-enforced}: at $K$ two bands meet because they carry a
      two-dimensional small irrep. Any Hamiltonian with the symmetry of graphene must
      have a degeneracy there: the Dirac point.
\item \emph{Accidental}: along a symmetry line two bands with \emph{different}
      one-dimensional irreps cross. Group theory permits the crossing (the
      Hamiltonian cannot couple different irreps) but does not require it.
      Moving the bands changes where they cross, or removes the crossing.
\end{itemize}
Two further observations. At an isolated high-symmetry point an accidental
degeneracy requires fine-tuning of the energies to infinite precision, so it has
probability zero. If you find one numerically, either the precision is insufficient
or you have missed a symmetry. Along a line, $\vec{k}$ is a continuous parameter, and
crossings of different irreps happen generically. Two bands of the \emph{same} irrep
generically repel (avoided crossing), because they can be coupled.

Near a degeneracy the bands disperse linearly unless a symmetry forbids the linear
term. Group theory also decides this, through the method of invariants
(\cref{sec:kp}).

\section{Time reversal in crystals\beyondmark}\label{sec:TRcrystal}

\Cref{sec:TR} treated time reversal for an atom or a molecule. In a crystal it also
acts on the wave vector.

\subsection{Time reversal and Bloch states}
Complex conjugation turns $\ee^{\ii\vec{k}\cdot\vec{r}}u_{\vec{k}}$ into
$\ee^{-\ii\vec{k}\cdot\vec{r}}u^{*}_{\vec{k}}$: time reversal maps $\vec{k}$ to $-\vec{k}$.
Hence
\begin{itemize}
\item $E_n(\vec{k})=E_n(-\vec{k})$ (without spin), or $E_{n\uparrow}(\vec{k})=E_{n\downarrow}(-\vec{k})$
      and in general Kramers pairs at $\vec{k}$ and $-\vec{k}$ (with spin);
\item at the \emph{time-reversal-invariant momenta} (TRIM), $\vec{k}\equiv-\vec{k}$ (in
      graphene: $\Gamma$ and the three $M$ points), time reversal maps the states at
      $\vec{k}$ onto themselves. With $\TR^2=-1$ every band is then Kramers degenerate;
      in the presence of inversion symmetry as well, every band is doubly degenerate
      at every $\vec{k}$.
\end{itemize}
$K$ and $K'=-K$ in graphene are exchanged by time reversal; $K$ is not a TRIM.

At a TRIM of a symmorphic space group, the test of \cref{sec:herring} applies to the
little co-group. For general $\vec{k}$ and non-symmorphic groups the same question is
answered by Herring's criterion, a sum of $\chi\bigl((\TR g)^2\bigr)$ over the elements
$g$ that map $\vec{k}$ to $-\vec{k}$; see Bradley--Cracknell.

\subsection{Magnetic groups}\label{sec:magnetic}
In a magnetically ordered crystal time reversal alone is broken, but a combination
$\TR g$ with a spatial operation $g$ may survive. For example, an antiferromagnet can be
invariant under $\TR$ combined with a translation or with inversion. The groups that
contain such antiunitary elements are the \emph{magnetic} (Shubnikov) groups: 1651
magnetic space groups in three dimensions, of four types (no $\TR$ at all; $\TR$ itself,
the ``grey'' groups; $\TR$ only combined with point operations; $\TR$ combined with a
translation). Their representation theory uses Wigner's \emph{corepresentations}, whose
degeneracies are again decided by tests of the type above. The Bilbao server tabulates
their irreducible corepresentations and band representations (magnetic topological
quantum chemistry).

\exercisesection

\begin{exercise}\label{ex:6:C4v}
Show that $(\vec{E}\times\vec{P})_z=E_xP_y-E_yP_x$ is invariant under $C_4$ but changes
sign under the mirrors of $\Cv{4}$. To which irrep of $\Cv{4}$ does it belong?
\end{exercise}

\begin{solution}
Under $C_4^+$, $(E_x,E_y)\to(-E_y,E_x)$ and likewise for $\vec{P}$:
$E_x'P_y'-E_y'P_x'=(-E_y)P_x-E_x(-P_y)=E_xP_y-E_yP_x$. Under the mirror
$x\to-x$: $(-E_x)P_y-E_y(-P_x)=-(E_xP_y-E_yP_x)$. So it is invariant under the rotations
and odd under the mirrors: $\irr{A}_2$ of $\Cv{4}$, like $R_z$.
\end{solution}

\begin{exercise}\label{ex:6:recip}
For the hexagonal lattice $\vec{a}_1=a(1,0)$, $\vec{a}_2=a(\tfrac12,\tfrac{\sqrt3}{2})$:
(a) find $\vec{b}_1,\vec{b}_2$; (b) show that $K=\tfrac23\vec{b}_1+\tfrac13\vec{b}_2$ is a
corner of the Brillouin zone with $\abs{K}=4\cpi/(3a)$, and that $M=\tfrac12\vec{b}_1$ is
the midpoint of an edge; (c) show that the six corners fall into two classes of three
equivalent points, $K$ and $K'\equiv-K$, and that $K\not\equiv K'$.
\end{exercise}

\begin{solution}
(a) $\vec{b}_1=\tfrac{2\cpi}{a}(1,-\tfrac{1}{\sqrt3})$, $\vec{b}_2=\tfrac{2\cpi}{a}(0,\tfrac{2}{\sqrt3})$,
$\abs{\vec{b}_i}=\tfrac{4\cpi}{\sqrt3a}$.
(b) $K=\tfrac23\vec{b}_1+\tfrac13\vec{b}_2=(\tfrac{4\cpi}{3a},0)$. The zone is a hexagon with
inradius $\abs{\vec{b}}/2=\tfrac{2\cpi}{\sqrt3a}$ and circumradius
$\tfrac{2\cpi}{\sqrt3a}/\cos30^\circ=\tfrac{4\cpi}{3a}=\abs{K}$: $K$ is a corner. $M=\vec{b}_1/2$
has $\abs{M}$ equal to the inradius and lies on the edge bisector: the midpoint of an
edge.
(c) $C_3^+K-K=\tfrac{4\cpi}{3a}(-\tfrac32,\tfrac{\sqrt3}{2})=-\vec{b}_1$, so $C_3K\equiv K$; the
corners at $0^\circ,120^\circ,240^\circ$ are equivalent, and so are those at
$60^\circ,180^\circ,300^\circ$ (equivalent to $K'=-K$). But $K'-K=-2K=(-\tfrac{8\cpi}{3a},0)$
would need $m_1=-\tfrac43$ in $m_1\vec{b}_1+m_2\vec{b}_2$: not a reciprocal vector, so
$K\not\equiv K'$.
\end{solution}

\begin{exercise}\label{ex:6:seitz}
Prove the composition and inverse rules \eqref{eq:seitz}. Show that
$g\seitz{E}{\vec{t}}g^{-1}=\seitz{E}{R\vec{t}}$ for $g=\seitz{R}{\vec{v}}$, so that
$\Tgroup$ is a normal subgroup.
\end{exercise}

\begin{solution}
$\seitz{R_1}{\vec{v}_1}\seitz{R_2}{\vec{v}_2}\vec{r}=R_1(R_2\vec{r}+\vec{v}_2)+\vec{v}_1
=R_1R_2\vec{r}+\vec{v}_1+R_1\vec{v}_2$.
$\seitz{R^{-1}}{-R^{-1}\vec{v}}\seitz{R}{\vec{v}}\vec{r}=R^{-1}(R\vec{r}+\vec{v})-R^{-1}\vec{v}=\vec{r}$.
Finally
$\seitz{R}{\vec{v}}\seitz{E}{\vec{t}}\seitz{R^{-1}}{-R^{-1}\vec{v}}
=\seitz{R}{\vec{v}+R\vec{t}}\seitz{R^{-1}}{-R^{-1}\vec{v}}=\seitz{E}{R\vec{t}}$, and $R\vec{t}$ is a
lattice vector because $R$ maps the lattice onto itself.
\end{solution}

\begin{exercise}\label{ex:6:bloch}
Verify from \eqref{eq:blochsum} that $\hat{T}_{\vec{t}'}\Phi_{\alpha\vec{k}}=
\ee^{-\ii\vec{k}\cdot\vec{t}'}\Phi_{\alpha\vec{k}}$ and $\Phi_{\alpha,\vec{k}+\vec{G}}=
\Phi_{\alpha\vec{k}}$. What changes if the phase $\ee^{\ii\vec{k}\cdot(\vec{t}+\vec{q}_\alpha)}$
is used instead?
\end{exercise}

\begin{solution}
$\hat{T}_{\vec{t}'}\Phi_{\alpha\vec{k}}(\vec{r})=\sum_{\vec{t}}\ee^{\ii\vec{k}\cdot\vec{t}}
\varphi_\alpha(\vec{r}-\vec{t}'-\vec{t}-\vec{q}_\alpha)=\sum_{\vec{s}}\ee^{\ii\vec{k}\cdot(\vec{s}-\vec{t}')}
\varphi_\alpha(\vec{r}-\vec{s}-\vec{q}_\alpha)=\ee^{-\ii\vec{k}\cdot\vec{t}'}\Phi_{\alpha\vec{k}}$
with $\vec{s}=\vec{t}+\vec{t}'$. Since $\ee^{\ii\vec{G}\cdot\vec{t}}=1$, $\Phi_{\alpha,\vec{k}+\vec{G}}=\Phi_{\alpha\vec{k}}$.
With the phase $\ee^{\ii\vec{k}\cdot(\vec{t}+\vec{q}_\alpha)}$ each Bloch sum is multiplied by
$\ee^{\ii\vec{k}\cdot\vec{q}_\alpha}$: the translation property is unchanged, but now
$\Phi_{\alpha,\vec{k}+\vec{G}}=\ee^{\ii\vec{G}\cdot\vec{q}_\alpha}\Phi_{\alpha\vec{k}}$, so $\mat{H}(\vec{k})$ is
periodic only up to a diagonal unitary transformation. Both conventions are in use.
The one with $\vec{q}_\alpha$ makes $\mat{H}(\vec{k})$ reflect the positions of the orbitals
(useful for Berry phases); the other makes it periodic.
\end{solution}

\begin{exercise}[a glide line]\label{ex:6:glide}
The plane group $pg$ (rectangular lattice $a\vec{e}_x$, $b\vec{e}_y$) is generated by
translations and the glide $g=\seitz{m_x}{\tfrac{b}{2}\vec{e}_y}$, where
$m_x\colon(x,y)\mapsto(-x,y)$. (a) Show that $g^2=\seitz{E}{b\vec{e}_y}$. (b) On the line
$k_x=0$, Bloch states can be chosen as eigenstates of $\hat{O}_g$; find the possible
eigenvalues as functions of $k_y$. (c) Follow an eigenvalue continuously from
$k_y=-\cpi/b$ to $k_y=\cpi/b$ and show that the two eigenvalue branches are exchanged:
bands on this line come in connected pairs.
\end{exercise}

\begin{solution}
(a) $g^2=\seitz{m_x^2}{\tfrac{b}{2}\vec{e}_y+m_x\tfrac{b}{2}\vec{e}_y}=\seitz{E}{b\vec{e}_y}$, since
$m_x$ leaves $\vec{e}_y$ unchanged.
(b) On $k_x=0$, $m_x\vec{k}=\vec{k}$, so $g$ is in the little group and $\hat{O}_g$ can be
diagonalised together with $\hat{H}$. Its eigenvalues satisfy
$\lambda^2=\ee^{-\ii k_yb}$ (the eigenvalue of $\hat{T}_{b\vec{e}_y}$), so
$\lambda_\pm(k_y)=\pm\ee^{-\ii k_yb/2}$.
(c) $\lambda_+(-\cpi/b)=\ee^{\ii\cpi/2}=\ii$ and $\lambda_+(\cpi/b)=\ee^{-\ii\cpi/2}=-\ii=\lambda_-(-\cpi/b)$.
Since $k_y=\cpi/b$ and $k_y=-\cpi/b$ are the same point, a band that starts in the
branch $\lambda_+$ at $-\cpi/b$ arrives in the branch $\lambda_-$: it cannot close on itself
and must connect to a partner. Bands along this line come in connected pairs, a
consequence of the fractional translation.
\end{solution}

\chapter{The group theory of electron bands}\label{ch:bands}

\begin{watch}{\lec{7}}
Generators, graphically \ts{7}{196}{3:16} $\cdot$
tight binding for graphene \ts{7}{580}{9:40} $\cdot$
the point group and $\sh$ parity: $\upsigma$ and $\uppi$ bands \ts{7}{1083}{18:03} $\cdot$
little groups \ts{7}{1476}{24:36} $\cdot$
high-symmetry points and lines; why $K$ has $\Cv{3}$ \ts{7}{1728}{28:48} $\cdot$
the site-symmetry group \ts{7}{2104}{35:04} $\cdot$
orbital mixing and distorted orbitals \ts{7}{2421}{40:21} $\cdot$
the character formula for band states \ts{7}{3174}{52:54} $\cdot$
$\Gamma$ point \ts{7}{3551}{59:11} $\cdot$
$K$ point and its phases \ts{7}{3867}{64:27} $\cdot$
compatibility relations \ts{7}{4381}{73:01} $\cdot$
what symmetry cannot decide \ts{7}{5191}{86:31}
\end{watch}

This chapter applies the machinery to one concrete crystal. We label every $\uppi$ band
of graphene by small irreps, prove that the Dirac point is enforced by symmetry, and
introduce compatibility relations. They are the bridge to topology in
\cref{ch:topology}.

\section{Generators}\label{sec:generators}

A set of elements \emph{generates} $G$ if every element is a product of them and their
inverses. $C_n$ is generated by $C_n^+$; $\Cv{n}$ by $C_n^+$ and one mirror; $\Td$ by
$S_4$ and $C_3$ (or $S_4$ and a mirror that does not contain the $S_4$ axis). To check
that a set of functions transforms correctly under a representation, or that a matrix
commutes with a representation, it is enough to check the generators.

The lecture opens with a graphical check: mark a generic point on a stereographic
projection, apply the candidate generators repeatedly, and count the distinct images.
If their number equals $\abs{G}$, the set generates the group, and each element can be
read off as a word in the generators.

\section{A tight-binding model for graphene}

Carbon has the configuration $1s^2\,2s^2\,2p^2$. The $1s$ electrons are inert, so a
minimal tight-binding model uses the four valence orbitals $2s,2p_x,2p_y,2p_z$ on each
of the two atoms of the cell: $8\times8$ matrices $\mat{H}(\vec{k})$ and eight bands.
With nearest- and next-nearest-neighbour hopping such a model already reproduces
first-principles bands well near the Fermi level. Adding empty $3s$ orbitals improves
the upper bands.

\paragraph{The point group.} With the origin at a hexagon centre (\cref{fig:honeycomb}),
the honeycomb is invariant under rotations by multiples of $60^\circ$ and under six
vertical mirrors. They fall into two classes (\cref{fig:mirrors}):
\begin{itemize}
\item $\sv$: the three mirrors that \emph{contain} C--C bonds (lines at $30^\circ$,
      $90^\circ$, $150^\circ$ with our axes); they map each sublattice onto itself;
\item $\sd$: the three mirrors that \emph{bisect} C--C bonds (lines at $0^\circ$,
      $60^\circ$, $120^\circ$); they exchange sublattices A and B.
\end{itemize}
The in-plane point group is $\Cv{6}$ (plane group $p6mm$). Graphene is also invariant
under the reflection $\sh$ in its own plane, which doubles the group to $\Dh{6}$.
Since $\sh$ commutes with everything, we can use it once and for all:
$2s,2p_x,2p_y$ are even under $\sh$ and $2p_z$ is odd. The Hamiltonian cannot couple
states of opposite $\sh$ parity, so the $8\times8$ problem splits into the six
$\upsigma$ bands (bonding) and the two $\uppi$ bands from the $p_z$ orbitals, which contain
the Fermi level and the Dirac point. From now on we study the $\uppi$ bands with the group
$\Cv{6}$.

\begin{equation}\label{eq:C6vtable}
\begin{array}{c|cccccc}
\Cv{6} & E & C_2 & 2C_3 & 2C_6 & 3\sv & 3\sd\\\hline
\irr{A}_1 & 1&1&1&1&1&1\\
\irr{A}_2 & 1&1&1&1&-1&-1\\
\irr{B}_1 & 1&-1&1&-1&1&-1\\
\irr{B}_2 & 1&-1&1&-1&-1&1\\
\irr{E}_1 & 2&-2&-1&1&0&0\\
\irr{E}_2 & 2&2&-1&-1&0&0
\end{array}
\end{equation}
Which of the two $\irr{B}$ irreps is called $\irr{B}_1$ depends on which mirrors are
called $\sv$; with our choice $\irr{B}_1$ is even under the bond-containing mirrors.

\section{Little groups of graphene}

\begin{figure}[ht]
\centering
\begin{tikzpicture}[scale=1.25,>=Stealth]
  \draw[thick] (0:2)--(60:2)--(120:2)--(180:2)--(240:2)--(300:2)--cycle;
  \foreach \a in {30,90,150}{\draw[blue!60!black,dashed] (\a:-2.4)--(\a:2.4);}
  \foreach \a in {0,60,120}{\draw[red!60!black,dotted,thick] (\a:-2.4)--(\a:2.4);}
  \node[blue!60!black,font=\scriptsize] at (90:2.6) {$\sv$};
  \node[red!60!black,font=\scriptsize] at (60:2.65) {$\sd$};
  \fill (0,0) circle (0.05) node[above left,font=\small]{$\Gamma$};
  \fill (0:2) circle (0.05) node[right,font=\small]{$K$};
  \fill (-30:1.732) circle (0.05) node[below right,font=\small]{$M$};
  \draw[very thick,orange!80!black] (0,0)--(0:2) node[midway,above,font=\small]{$T$};
  \draw[very thick,orange!80!black] (0,0)--(-30:1.732) node[midway,below left,font=\small]{$\Sigma$};
  \draw[very thick,orange!80!black] (0:2)--(-30:1.732) node[midway,right,font=\small]{$T'$};
\end{tikzpicture}
\caption{The Brillouin zone of graphene with its high-symmetry points and lines. Blue
dashed: mirrors $\sv$ (containing C--C bonds in real space). Red dotted: mirrors $\sd$
(bisecting bonds). With $\vec{a}_1=a(1,0)$, the point $K=\tfrac23\vec{b}_1+\tfrac13\vec{b}_2
=(\tfrac{4\cpi}{3a},0)$ and $M=\tfrac12\vec{b}_1$.}
\label{fig:mirrors}
\end{figure}

The little co-group of each point or line is the set of point operations that map it
to an equivalent point:
\begin{equation}\label{eq:littlegroups}
\begin{array}{llll}
\toprule
\vec{k} & \bar{G}_{\vec{k}} & \text{elements} & \\
\midrule
\Gamma & \Cv{6} & \text{all 12} &\\
T\ (\Gamma K) & \Cs & E,\ \sd(0^\circ) &\\
K & \Cv{3} & E,\ C_3^\pm,\ 3\sd &\\
T'\ (KM) & \Cs & E,\ \sd(60^\circ) &\\
M & \Cv{2} & E,\ C_2,\ \sv(150^\circ),\ \sd(60^\circ) &\\
\Sigma\ (M\Gamma) & \Cs & E,\ \sv(150^\circ) &\\
\bottomrule
\end{array}
\end{equation}
Two remarks.
\begin{itemize}
\item At $K$ the rotation $C_3$ carries $K$ to another corner of the hexagon. The two
      corners differ by a reciprocal-lattice vector, so they label the \emph{same} irrep
      of $\Tgroup$ and $C_3\in\bar{G}_K$. For the same reason the six corners are only
      two inequivalent points, $K$ and $K'=-K$. This is why graphene has two Dirac
      points, not six.
\item The little co-group of a line is a subgroup of the little co-groups of its end
      points. This is what makes compatibility relations possible
      (\cref{sec:compat}).
\end{itemize}

\section{The site-symmetry group}\label{sec:sitesym}

\begin{definition}
The \emph{site-symmetry group} $G_{\vec{q}}$ of a point $\vec{q}$ is the set of
space-group elements that leave $\vec{q}$ fixed: $g\vec{q}=\vec{q}$ exactly (not up to a
lattice vector). It is always isomorphic to a point group. Points related by the space
group form an \emph{orbit}. Orbits whose site-symmetry groups are conjugate form a
\emph{Wyckoff position}.
\end{definition}
The crystal as a whole is invariant under $\Cv{6}$ about a hexagon centre, but an atom
does not see six-fold symmetry: along one direction it has a nearest neighbour, along
the opposite direction a third-nearest one. The site-symmetry group of a carbon atom
is $\Cv{3}$, generated by the three-fold rotation about the atom and the three $\sv$
mirrors through it. For $p6mm$ the Wyckoff positions with non-trivial site symmetry
are
\[
1a\ \text{(hexagon centres, }\Cv{6}),\qquad
2b\ \text{(the atoms, }\Cv{3}),\qquad
3c\ \text{(bond centres, }\Cv{2}).
\]

\paragraph{Atomic orbitals become irreps of the site group.} In the crystal the atomic
orbitals are no longer those of a free atom; they transform under irreps of
$G_{\vec{q}}$, obtained by subduction from $\Og{3}$ (\cref{sec:crystalfield}). At a
$\Cv{3}$ site, $s\to\irr{a}_1$ and $p\to\irr{a}_1\oplus\irr{e}$ with $p_z\in\irr{a}_1$
and $(p_x,p_y)\in\irr{e}$. The ``$p_z$ orbital'' of a tight-binding model of graphene
is really an $\irr{a}_1$ orbital of $\Cv{3}$: a $p_z$ orbital distorted by the crystal
field and mixed with other $\irr{a}_1$ orbitals.

This has practical consequences. A model built as if the orbitals were spherically
symmetric can miss couplings that are forbidden for free-atom orbitals but allowed by
the true, lower site symmetry. For example, the hoppings to two neighbours at equal
distance but in inequivalent directions need not be equal. Such an over-restricted
model may display spurious degeneracies or miss phases. Knowing the site-symmetry
irreps guarantees that the model is the most general one allowed by symmetry.

\section{Characters of band states}

Place one $\irr{a}_1$ orbital at each atom and form the two Bloch sums
$\Phi_{\mathrm{A}\vec{k}}$, $\Phi_{\mathrm{B}\vec{k}}$ of \eqref{eq:blochsum}. At each
$\vec{k}$ they span a two-dimensional representation of the little group, and
$\mat{H}(\vec{k})$ is a $2\times2$ matrix commuting with it. To find its small-irrep
content we need its character.

Let $g=\seitz{R}{\vec{0}}$ be in the little group of $\vec{k}$. If $g$ maps the atom at
$\vec{q}_\alpha$ onto an atom of the same sublattice, $g\vec{q}_\alpha=\vec{q}_\alpha+\vec{t}_\alpha$
with $\vec{t}_\alpha\in\Tgroup$, the Bloch sum is mapped onto itself up to a phase.
Otherwise it is mapped onto the other sublattice and does not contribute to the trace.
A short computation with \eqref{eq:blochsum} gives

\begin{keyresult}[Character of a band representation at $\vec{k}$]
\begin{equation}\label{eq:bandchar}
\chi_{\vec{k}}(g)=\sum_{\alpha\,:\,g\vec{q}_\alpha\equiv\vec{q}_\alpha}
\ee^{-\ii\vec{k}\cdot\vec{t}_\alpha}\;\chi_\rho(g_\alpha),
\qquad \vec{t}_\alpha=g\vec{q}_\alpha-\vec{q}_\alpha,
\end{equation}
where the sum runs over the atoms of the cell left invariant up to a lattice
translation, $g_\alpha=\seitz{E}{-\vec{t}_\alpha}g\in G_{\vec{q}_\alpha}$ is the
corresponding site-symmetry operation, and $\chi_\rho$ is the character of the orbital
irrep $\rho$ of the site group.
\end{keyresult}
This is the crystal version of $\chiM=\Ninv\chiV$ (\cref{sec:mechanical}): count invariant
atoms, weight each by the orbital character, and add the Bloch phase. For our
$\irr{a}_1$ orbitals $\chi_\rho=1$ and only the phases remain.

\subsection{The \texorpdfstring{$\Gamma$}{Gamma} point}
At $\vec{k}=0$ all phases are $1$ and $\chi_\Gamma(g)$ is the number of sublattices
mapped onto themselves: $2$ for $E$, $C_3$, $\sv$; $0$ for $C_2$, $C_6$, $\sd$ (these
exchange A and B). From \eqref{eq:C6vtable},
\[
\chi_\Gamma=(2,0,2,0,2,0)\qquad\Longrightarrow\qquad
\Gamma=\irr{A}_1\oplus\irr{B}_1
\]
(in the Bilbao labels used in the lecture, $\Gamma_1\oplus\Gamma_3$). Two
non-degenerate levels: the bonding combination $\Phi_{\mathrm{A}}+\Phi_{\mathrm{B}}$
($\irr{A}_1$) and the antibonding one $\Phi_{\mathrm{A}}-\Phi_{\mathrm{B}}$ ($\irr{B}_1$,
odd under the operations that exchange sublattices).

\subsection{The \texorpdfstring{$K$}{K} point}\label{sec:Kpoint}
Now phases matter. Only $E$ and $C_3^\pm$ keep the sublattices; the three $\sd$ in
$\bar{G}_K$ exchange them. Take $C_3^+$ about the origin. It maps
$\vec{q}_{\mathrm{A}}=\tfrac13(\vec{a}_1+\vec{a}_2)$ to $\vec{q}_{\mathrm{A}}-\vec{a}_1$ and
$\vec{q}_{\mathrm{B}}=\tfrac23(\vec{a}_1+\vec{a}_2)$ to $\vec{q}_{\mathrm{B}}-2\vec{a}_1$. With
$K=\tfrac23\vec{b}_1+\tfrac13\vec{b}_2$ we have $K\cdot\vec{a}_1=\tfrac{4\cpi}{3}$, so the
phases are
\[
\ee^{\ii K\cdot\vec{a}_1}=\ee^{4\cpi\ii/3}=\omega^{*},\qquad
\ee^{2\ii K\cdot\vec{a}_1}=\ee^{8\cpi\ii/3}=\omega,\qquad
\omega=\ee^{2\cpi\ii/3},
\]
and $\chi_K(C_3)=\omega+\omega^{*}=-1$. Physically: under the rotation, the Bloch sum
on sublattice A is multiplied by $\omega^{*}$, the one on B by $\omega$. Thus
\[
\chi_K=(2,-1,0)\ \text{on}\ (E,2C_3,3\sd)\qquad\Longrightarrow\qquad K=\irr{E}\quad(K_3).
\]
\begin{keyresult}[The Dirac point is enforced by symmetry]
At $K$ the two $\uppi$ states form the two-dimensional irrep $\irr{E}$ of $\Cv{3}$, so they
are degenerate for \emph{every} Hamiltonian with the symmetry of graphene. Whether the
dispersion near the degeneracy is linear is a further question, answered in
\cref{sec:kp}: it is.
\end{keyresult}

\subsection{The \texorpdfstring{$M$}{M} point}
$\bar{G}_M=\Cv{2}=\{E,C_2,\sv,\sd\}$. The $C_2$ and $\sd$ exchange sublattices, and
$\sv(150^\circ)$ keeps them with trivial phases, so $\chi_M=(2,0,2,0)$ and
$M=\irr{A}_1\oplus\irr{B}_1$, where in $\Cv{2}$ we call $\irr{B}_1$ the irrep that is
odd under $C_2$ and even under $\sv$.

\section{\texorpdfstring{$\vec{k}\cdot\vec{p}$}{k.p} Hamiltonians and the method of invariants\beyondmark}\label{sec:kp}
\sectionmark{The method of invariants}

Group theory tells us which bands must be degenerate at a high-symmetry point. The
\emph{method of invariants} tells us how they disperse nearby: it constructs the most
general effective Hamiltonian allowed by symmetry, order by order in the distance from
the point. It is the same idea as the invariant polarizability of \cref{sec:invariants}
and the invariant dynamical matrix of \cref{sec:wigner}.

\subsection{The effective Hamiltonian}
Let $\vec{k}_0$ be a high-symmetry point and let $\{\ket{n}\}$, $n=1,\dots,d$, be the states
of the (possibly degenerate) band multiplet of interest at $\vec{k}_0$, carrying a
representation $\mat{D}(g)$ of the little co-group $\bar{G}_{\vec{k}_0}$ (usually one small
irrep). For $\vec{k}=\vec{k}_0+\vec{q}$, projecting the Hamiltonian onto these states
(L\"owdin partitioning, which includes the other bands perturbatively) gives a
$d\times d$ matrix $\mat{H}(\vec{q})$, expanded in powers of $\vec{q}$:
\[
\mat{H}(\vec{q})=\mat{H}_0+\sum_iq_i\mat{H}_i+\sum_{ij}q_iq_j\mat{H}_{ij}+\dots
\]
At first order $\mat{H}_i\propto\bra{n}\hat{p}_i\ket{n'}$, hence the name
$\vec{k}\cdot\vec{p}$.

\subsection{The symmetry constraint}
For every $g$ in the little co-group, with rotational part $R_g$,
\begin{equation}\label{eq:kpconstraint}
\mat{D}(g)\,\mat{H}(R_g^{-1}\vec{q})\,\mat{D}(g)^{-1}=\mat{H}(\vec{q}).
\end{equation}
Expand $\mat{H}(\vec{q})=\sum_ac_a\,\mat{X}_a\,f_a(\vec{q})$, where the hermitian matrices
$\mat{X}_a$ span the $d^2$-dimensional space of $d\times d$ hermitian matrices and
$f_a$ are polynomials in $\vec{q}$. Both sets carry representations of the little
co-group: the matrices transform as $\mat{D}\otimes\mat{D}^{*}$ (by conjugation), the
polynomials of degree $n$ as the symmetric power $[\GV^n]$ restricted to the relevant
components of $\vec{q}$. Then \eqref{eq:kpconstraint} says that $\mat{H}$ is an
invariant, bilinear in matrices and polynomials. Hence
\begin{keyresult}[Method of invariants]
Decompose the matrices and the polynomials of each degree into irreps (with matched
bases). The most general $\mat{H}$ is a sum of one invariant, with a free real
coefficient, for each pair (matrix irrep, polynomial irrep) related as in
\cref{sec:countinv}. The number of free parameters at order $n$ equals the multiplicity
of $\irr{A}_1$ in $(\mat{D}\otimes\mat{D}^{*})\otimes[\GV^n]$.
\end{keyresult}
Additional antiunitary symmetries (time reversal, or time reversal combined with a
spatial operation that maps $\vec{k}_0$ to itself) add further constraints of the same
type.

\subsection{Example: the Dirac cone of graphene}
At $K$ the $\uppi$ states carry the two-dimensional irrep $\irr{E}$ of $\Cv{3}$
(\cref{sec:Kpoint}). Choose the band basis so that the matrices of $\irr{E}$ are those of
in-plane vectors, $\mat{D}(C_3^+)=$ rotation by $120^\circ$ and $\mat{D}(\sd)=\diag(1,-1)$
for the mirror that reflects $q_y\to-q_y$. The four hermitian $2\times2$ matrices
decompose as
\[
\mat{D}\otimes\mat{D}^{*}=\irr{E}\otimes\irr{E}=\irr{A}_1\oplus\irr{A}_2\oplus\irr{E}:
\qquad \one\ (\irr{A}_1),\quad \sigma_y\ (\irr{A}_2),\quad (\sigma_z,\sigma_x)\ (\irr{E}).
\]
Conjugation by a rotation by $\varphi$ rotates the pair $(\sigma_z,\sigma_x)$ by
$2\varphi$. For $\varphi=\pm120^\circ$ this is the same as rotating by $\mp120^\circ$, so
$(\sigma_z,-\sigma_x)$ transforms exactly like $(q_x,q_y)$. The polynomials decompose as:
degree $0$: $1$ ($\irr{A}_1$); degree $1$: $(q_x,q_y)$ ($\irr{E}$); degree $2$:
$q_x^2+q_y^2$ ($\irr{A}_1$) and $(q_x^2-q_y^2,\,2q_xq_y)$ ($\irr{E}$, rotating by $2\varphi$
like $(\sigma_z,\sigma_x)$). Pairing equal irreps gives, up to second order,
\begin{equation}\label{eq:dirac}
\mat{H}(\vec{q})=\varepsilon_0\one
+v\,(q_x\sigma_z-q_y\sigma_x)
+\alpha\,q^2\one
+\beta\bigl[(q_x^2-q_y^2)\sigma_z+2q_xq_y\sigma_x\bigr]+O(q^3).
\end{equation}
Consequences, all from symmetry:
\begin{itemize}
\item No constant $\sigma_y$ term: the mass term transforms as $\irr{A}_2$ and there is no
      $\irr{A}_2$ polynomial of degree $0$. The degeneracy at $K$ cannot be lifted without
      breaking the symmetry.
\item One linear invariant: the dispersion is linear and isotropic,
      $E-\varepsilon_0=\pm v\abs{\vec{q}}+O(q^2)$, a Dirac cone. The matrix
      $q_x\sigma_z-q_y\sigma_x$ is unitarily equivalent to the familiar
      $q_x\sigma_x+q_y\sigma_y$.
\item At second order, the $\beta$ term produces the three-fold \emph{trigonal warping}
      of the cone.
\item If the sublattices become inequivalent (as in hexagonal boron nitride), the
      relevant irrep at $K$ splits and a constant mass term, and a gap, become allowed.
\end{itemize}
The linear term exists because $\irr{E}\otimes\irr{E}$ contains the irrep of $\vec{q}$. When
that fails (for other groups or irreps), the crossing is quadratic or of higher order.
This is exactly the question asked in \lec{6}. Complete catalogues of such
$\vec{k}\cdot\vec{p}$ models for all space and magnetic groups have been built with this
method.

\paragraph{Further reading.}
G.\,L.\ Bir and G.\,E.\ Pikus, \emph{Symmetry and Strain-Induced Effects in Semiconductors}
(Wiley, 1974), the classic source of the method; R.\ Winkler, \emph{Spin--Orbit Coupling
Effects in Two-Dimensional Electron and Hole Systems} (Springer, 2003), with tables of
invariants; Dresselhaus, Dresselhaus and Jorio, the chapter on $\vec{k}\cdot\vec{p}$
perturbation theory.

\section{Compatibility relations}\label{sec:compat}

Along a line $\vec{k}$ varies continuously, and the little co-group is a subgroup of
those of the end points. A state at $\Gamma$ carrying an irrep $\Gamma^{(\mu)}$ evolves
continuously into the line, where it must carry the subduced representation
$\Gamma^{(\mu)}\!\downarrow\!\bar{G}_{\text{line}}$. These subductions are the
\emph{compatibility relations}. On every line $\bar{G}=\Cs$ with irreps $\irr{A}'$
(even under the mirror) and $\irr{A}''$ (odd). Reading the character of the
relevant mirror:
\[
\begin{array}{llll}
\toprule
\text{line (mirror)} & \text{from } \Gamma & \text{from } K & \text{from } M\\
\midrule
T\ (\sd) & \irr{A}_1\to\irr{A}',\ \irr{B}_1\to\irr{A}'' & \irr{E}\to\irr{A}'\oplus\irr{A}'' & \\
\Sigma\ (\sv) & \irr{A}_1\to\irr{A}',\ \irr{B}_1\to\irr{A}' & & \irr{A}_1\to\irr{A}',\ \irr{B}_1\to\irr{A}'\\
T'\ (\sd) & & \irr{E}\to\irr{A}'\oplus\irr{A}'' & \irr{A}_1\to\irr{A}',\ \irr{B}_1\to\irr{A}''\\
\bottomrule
\end{array}
\]
Each line can be reached from both ends, and the two subductions must agree. On $T$:
from $\Gamma$, $\irr{A}'\oplus\irr{A}''$; from $K$, $\irr{A}'\oplus\irr{A}''$. The same holds
on $\Sigma$ and $T'$. Such consistency checks catch almost every mistake. In practice
one never computes characters along lines; subduction from the end points gives them
for free.

\begin{figure}[ht]
\centering
\begin{tikzpicture}
\begin{axis}[width=11cm,height=7cm,xmin=0,xmax=9.9108,ymin=-4.2,ymax=3.2,
  xtick={0,4.1888,6.2832,9.9108},xticklabels={$\Gamma$,$K$,$M$,$\Gamma$},
  ylabel={$E/t$},xmajorgrids,grid style={gray!40},
  every axis plot/.append style={thick}]
\addplot[blue!70!black] table[x=x,y=lo]{graphene_pi.dat};
\addplot[red!70!black] table[x=x,y=hi]{graphene_pi.dat};
% labels at the high-symmetry points sit on the vertical grid lines: white background
\tikzset{pt/.style={font=\scriptsize,fill=white,inner sep=1.5pt}}
\node[pt,anchor=west]  at (axis cs:0.08,-3.95) {$\irr{A}_1$};
\node[pt,anchor=west]  at (axis cs:0.08,2.75) {$\irr{B}_1$};
\node[pt,anchor=south] at (axis cs:4.1888,0.6) {$\irr{E}$};
\node[pt,anchor=north] at (axis cs:6.2832,-1.0) {$\irr{B}_1$};
\node[pt,anchor=south] at (axis cs:6.2832,1.4) {$\irr{A}_1$};
% labels on the lines: each sits on the free side of its (monotonic) curve
\tikzset{ln/.style={font=\scriptsize,inner sep=2pt}}
\node[ln,blue!70!black,anchor=north west] at (axis cs:2.13,-2.06) {$\irr{A}'$};
\node[ln,red!70!black,anchor=south west]  at (axis cs:2.13,1.88)  {$\irr{A}''$};
\node[ln,blue!70!black,anchor=north east] at (axis cs:5.27,-0.51) {$\irr{A}''$};
\node[ln,red!70!black,anchor=south east]  at (axis cs:5.27,0.99)  {$\irr{A}'$};
\node[ln,blue!70!black,anchor=north east] at (axis cs:8.13,-2.48) {$\irr{A}'$};
\node[ln,red!70!black,anchor=south east]  at (axis cs:8.13,2.05)  {$\irr{A}'$};
\end{axis}
\end{tikzpicture}
\caption{The $\uppi$ bands of graphene in a tight-binding model with nearest- ($t$) and
next-nearest-neighbour ($t'=0.1\,t$) hopping, labelled by small irreps of $\Cv{6}$ and
its subgroups. The labels at $\Gamma$ and $M$ were obtained from the eigenvectors;
symmetry alone fixes only which pair of labels appears at each point.}
\label{fig:grapheneTB}
\end{figure}

\section{What symmetry decides and what it does not}

Symmetry fixes the \emph{set} of irreps at each point: $\{\irr{A}_1,\irr{B}_1\}$ at
$\Gamma$, $\{\irr{E}\}$ at $K$, $\{\irr{A}_1,\irr{B}_1\}$ at $M$. It does not fix their
energetic order. In \cref{fig:grapheneTB} the $\irr{A}_1$ state is below at $\Gamma$ but
above at $M$. The model produced this, not symmetry, and different hoppings could
reorder them. The compatibility relations then decide how the labels are connected.
The order of the irreps at high-symmetry points, combined with the compatibility
relations, is exactly the information that topological quantum chemistry uses.

Irrep labels are also what selection rules need: optical transitions between bands at
$\vec{k}$ are allowed or forbidden according to the small irreps involved
(\cref{sec:raman}).

\exercisesection

\begin{exercise}\label{ex:7:littlegroups}
Verify the little co-groups \eqref{eq:littlegroups}. In particular show that
$\sd(60^\circ)$ maps $K$ to an equivalent point and that $\sv(150^\circ)$ fixes $M$.
\end{exercise}

\begin{solution}
$\sd(60^\circ)$ reflects the direction $0^\circ$ of $K$ to $120^\circ$, i.e.\ to $C_3^+K$,
which is equivalent to $K$ (\cref{ex:6:recip}). $M$ lies at $-30^\circ$, on the mirror line
at $150^\circ\equiv-30^\circ$, so $\sv(150^\circ)$ fixes it exactly. By contrast
$\sv(30^\circ)$ sends $M$ to $90^\circ$, i.e.\ to $\vec{b}_2/2$, and
$\vec{b}_2/2-\vec{b}_1/2$ is not a reciprocal vector. The lines inherit the mirror that
contains them: $\sd(0^\circ)$ for $T$, $\sv(150^\circ)$ for $\Sigma$, and $\sd(60^\circ)$
for $T'$ (it maps the segment $KM$ onto an equivalent one).
\end{solution}

\begin{exercise}[the $\upsigma$ bands]\label{ex:7:sigma}
The $\upsigma$ bands come from $s$ ($\irr{a}_1$) and $(p_x,p_y)$ ($\irr{e}$) orbitals at the
$\Cv{3}$ sites. Use \eqref{eq:bandchar} to show that the six $\upsigma$ states decompose as
\[
\Gamma:\ \irr{A}_1\oplus\irr{B}_1\oplus\irr{E}_1\oplus\irr{E}_2,\qquad
K:\ \irr{A}_1\oplus\irr{A}_2\oplus2\,\irr{E},\qquad
M:\ 2\,\irr{A}_1\oplus\irr{A}_2\oplus2\,\irr{B}_1\oplus\irr{B}_2 .
\]
How many symmetry-enforced degeneracies are there at $\Gamma$ and at $K$?
\end{exercise}

\begin{solution}
Each atom carries $\irr{a}_1\oplus\irr{e}$ with site characters
$\chi_\rho(E)=3$, $\chi_\rho(C_3)=1-1=0$, $\chi_\rho(\sv)=1+0=1$.
$\Gamma$: sublattices kept by $E,C_3,\sv$ with trivial phases,
$\chi_\Gamma=(6,0,0,0,2,0)$ on $(E,C_2,2C_3,2C_6,3\sv,3\sd)$. The reduction gives
$m_{\irr{A}_1}=m_{\irr{B}_1}=\tfrac1{12}(6+6)=1$, $m_{\irr{E}_1}=m_{\irr{E}_2}=\tfrac1{12}\,12=1$,
others $0$.
$K$: $\chi_K(C_3)=(\omega+\omega^{*})\cdot0=0$, $\sd$ exchanges sublattices, so
$\chi_K=(6,0,0)$ and $K=\irr{A}_1\oplus\irr{A}_2\oplus2\,\irr{E}$.
$M$: $\chi_M=(6,0,2,0)$ on $(E,C_2,\sv,\sd)$, giving
$2\,\irr{A}_1\oplus\irr{A}_2\oplus2\,\irr{B}_1\oplus\irr{B}_2$.
Enforced degeneracies: two doubly degenerate levels at $\Gamma$ ($\irr{E}_1$, $\irr{E}_2$) and
two at $K$ ($2\,\irr{E}$); none at $M$.
\end{solution}

\begin{exercise}[nearest-neighbour model]\label{ex:7:NN}
With only nearest-neighbour hopping $t$, $\mat{H}(\vec{k})=-t\begin{pmatrix}0&f\\f^{*}&0\end{pmatrix}$
with $f(\vec{k})=1+\ee^{-\ii\vec{k}\cdot\vec{a}_1}+\ee^{-\ii\vec{k}\cdot\vec{a}_2}$ (Bloch sums
\eqref{eq:blochsum}). Compute the energies at $\Gamma$, $K$ and $M$ and identify the
irrep of each eigenvector by applying $C_2$.
\end{exercise}

\begin{solution}
$\Gamma$: $f=3$, $E=\mp3t$. The lower state $(1,1)/\sqrt2$ is invariant under all
operations: $\irr{A}_1$; the upper one $(1,-1)/\sqrt2$ is odd under the operations
exchanging sublattices: $\irr{B}_1$.
$K$: $K\cdot\vec{a}_1=\tfrac{4\cpi}{3}$, $K\cdot\vec{a}_2=\tfrac{2\cpi}{3}$, so
$f=1+\omega+\omega^{*}=0$: $E=0$, doubly degenerate ($\irr{E}$).
$M$: $M\cdot\vec{a}_1=\cpi$, $M\cdot\vec{a}_2=0$, $f=1-1+1=1$, $E=\mp t$. $C_2$ maps
$\vec{q}_{\mathrm{A}}\to-\vec{q}_{\mathrm{A}}=\vec{q}_{\mathrm{B}}-(\vec{a}_1+\vec{a}_2)$, so
$\hat{O}_{C_2}\Phi_{\mathrm{A}}=\ee^{\ii M\cdot(\vec{a}_1+\vec{a}_2)}\Phi_{\mathrm{B}}=-\Phi_{\mathrm{B}}$ and
likewise $\hat{O}_{C_2}\Phi_{\mathrm{B}}=-\Phi_{\mathrm{A}}$. The lower state $(1,1)$ has $C_2$
eigenvalue $-1$ ($\irr{B}_1$), the upper $(1,-1)$ has $+1$ ($\irr{A}_1$), as in
\cref{fig:grapheneTB}.
\end{solution}

\begin{exercise}\label{ex:7:compat}
Write the compatibility relations for the $\upsigma$ bands of \cref{ex:7:sigma} along
the line $T$, and check that the subductions from $\Gamma$ and from $K$ agree.
\end{exercise}

\begin{solution}
On $T$ the mirror is $\sd$. From $\Gamma$: $\irr{A}_1\to\irr{A}'$, $\irr{B}_1\to\irr{A}''$,
$\irr{E}_1\to\irr{A}'\oplus\irr{A}''$, $\irr{E}_2\to\irr{A}'\oplus\irr{A}''$: in total
$3\irr{A}'\oplus3\irr{A}''$. From $K$: $\irr{A}_1\to\irr{A}'$, $\irr{A}_2\to\irr{A}''$,
$2\,\irr{E}\to2(\irr{A}'\oplus\irr{A}'')$: again $3\irr{A}'\oplus3\irr{A}''$. The two subductions
agree.
\end{solution}

\chapter{Group theory and topology}\label{ch:topology}

\begin{watch}{\lec{8} (no captions; times read off the slides)}
Band topology \ts{8}{600}{$\approx$10:00} $\cdot$
Wannier representability \ts{8}{1100}{$\approx$18:20} $\cdot$
Wannier representability and TQC \ts{8}{1500}{$\approx$25:00} $\cdot$
classification of band representations \ts{8}{2100}{$\approx$35:00} $\cdot$
EBRs from Wyckoff positions \ts{8}{2700}{$\approx$45:00} $\cdot$
irreps from EBRs \ts{8}{3300}{$\approx$55:00} $\cdot$
double-valued EBRs of $p6mm$ \ts{8}{3900}{$\approx$65:00} $\cdot$
spinful $p_z$ orbitals in graphene \ts{8}{4400}{$\approx$73:20} $\cdot$
EBR analysis \ts{8}{4800}{$\approx$80:00} $\cdot$
what TQC can and cannot diagnose \ts{8}{5150}{$\approx$85:50}
\end{watch}

Topological quantum chemistry (TQC) connects the symmetry labels of \cref{ch:bands}
with the global, topological properties of bands. Its central idea fits in one
sentence: \emph{a set of bands is topologically trivial exactly when it can be
described by localised, symmetric orbitals in real space; the irreps at
high-symmetry points often reveal when this is impossible.}

\section{Local versus global}

A cylinder and a M\"obius strip are made by gluing two opposite edges of a rectangle,
with or without a twist. Every small piece of either looks like a piece of rectangle:
locally they are identical. Globally they differ, and no continuous deformation turns
one into the other. Topology studies such global properties.

Bands carry the same kind of structure. The Bloch Hamiltonian is periodic in
reciprocal space, $\mat{H}(\vec{k}+\vec{G})=\mat{H}(\vec{k})$, so the Brillouin zone is
a torus. At each $\vec{k}$ an isolated set of $n$ bands defines an $n$-dimensional
subspace of states. As $\vec{k}$ runs over the torus these subspaces form a
\emph{vector bundle}. Locally one can always choose smooth basis vectors (a gauge).
Whether one can choose them smoothly and periodically over the \emph{whole} zone is a
global question, and the answer can be no. The Chern number of a two-dimensional band
is the simplest obstruction.

\section{Wannier functions}

From Bloch states $\psi_{n\vec{k}}$ of an isolated set of bands one builds
\emph{Wannier functions}
\begin{equation}
w_{n\vec{t}}(\vec{r})=\frac{V_{\mathrm{cell}}}{(2\cpi)^d}\int_{\mathrm{BZ}}\dd^d k\;
\ee^{-\ii\vec{k}\cdot\vec{t}}\sum_m U_{mn}(\vec{k})\,\psi_{m\vec{k}}(\vec{r}),
\end{equation}
labelled by a band index and a lattice vector, and related by translations,
$w_{n\vec{t}}(\vec{r})=w_{n\vec{0}}(\vec{r}-\vec{t})$. The unitary matrix
$\mat{U}(\vec{k})$ is the gauge freedom. Wannier functions are the real-space,
orbital-like description of the bands, and their localisation depends on the gauge:
a smooth, periodic gauge gives exponentially localised functions.

\begin{keyresult}[Wannier representability]
The Wannier functions of an isolated set of bands can be chosen
\emph{exponentially localised and transforming among themselves under all the
symmetries of the crystal} if and only if the bands are topologically trivial. Such a
set of bands is called \emph{Wannier representable}.
\end{keyresult}
In TQC this equivalence \emph{defines} ``trivial''. For two-dimensional insulators
without extra symmetry it contains the theorem that exponentially localised Wannier
functions exist precisely when all Chern numbers vanish. The global property in
reciprocal space (topology) is thus tied to a local property in real space
(localisation of symmetric orbitals).

\section{Band representations}

\subsection{Definition}
Take a Wyckoff position with sites $\vec{q}_1,\dots,\vec{q}_n$ in the cell,
site-symmetry group $G_{\vec{q}_1}$, and a set of orbitals at $\vec{q}_1$ transforming as
a representation $\rho$ of $G_{\vec{q}_1}$. Copy them to the other sites with the
space-group operations (coset representatives: the cosets of $G_{\vec{q}_1}$ in $G$
label the sites of the orbit, \cref{sec:cosets}) and to all cells with translations.
The infinite set of orbitals so obtained is mapped onto itself by the space group:
it carries a representation of $G$, the \emph{band representation}
\begin{equation}
\rho\uparrow G\qquad\text{(the representation \emph{induced} from $\rho$)}.
\end{equation}
In the Bilbao notation, $\rho@w$ ($w$ being the Wyckoff label); for instance
$\irr{A}_1@2b$ is our graphene $p_z$ example. Fourier transforming to Bloch sums, the
band representation becomes, at each $\vec{k}$, a finite representation of the little
group, of dimension $n\dim\rho$, whose character is exactly
\eqref{eq:bandchar}. Its small-irrep content at the high-symmetry points is
the ``irrep vector'' of the band representation.

A set of bands is Wannier representable if and only if it carries (is equivalent to)
a band representation. The Wannier functions are then the orbitals from which the
band representation was induced.

\subsection{Elementary band representations}
Band representations can be added. A band representation that cannot be written as a
sum of other band representations is \emph{elementary} (EBR). EBRs are the building
blocks of all atomic-limit band structures.
\begin{itemize}
\item Every EBR is induced from an irrep of the site-symmetry group of a \emph{maximal}
      Wyckoff position (one whose site-symmetry group is not contained in that of
      another position). Orbitals at non-maximal positions can be moved continuously
      to a maximal one without breaking symmetry, so they give composite band
      representations. (A few irreps at maximal positions also give composite band
      representations; these exceptions are listed in the tables.)
\item For each space group the number of EBRs is finite. They are tabulated, with
      their irrep vectors, in the \texttt{BANDREP} application of the Bilbao
      Crystallographic Server, for all 230 space groups (with and without spin--orbit
      coupling) and for the magnetic groups.
\end{itemize}

For $p6mm$ the maximal positions are $1a$ ($\Cv{6}$), $2b$ ($\Cv{3}$) and $3c$
($\Cv{2}$). Evaluating \eqref{eq:bandchar} for every irrep of each site group gives
the spinless EBRs:
\begin{equation}\label{eq:ebrp6mm}
\begin{array}{llccc}
\toprule
\rho@w & \text{bands} & \Gamma\ (\Cv{6}) & K\ (\Cv{3}) & M\ (\Cv{2})\\
\midrule
\irr{A}_1@1a & 1 & \irr{A}_1 & \irr{A}_1 & \irr{A}_1\\
\irr{A}_2@1a & 1 & \irr{A}_2 & \irr{A}_2 & \irr{A}_2\\
\irr{B}_1@1a & 1 & \irr{B}_1 & \irr{A}_2 & \irr{B}_1\\
\irr{B}_2@1a & 1 & \irr{B}_2 & \irr{A}_1 & \irr{B}_2\\
\irr{E}_1@1a & 2 & \irr{E}_1 & \irr{E} & \irr{B}_1\oplus\irr{B}_2\\
\irr{E}_2@1a & 2 & \irr{E}_2 & \irr{E} & \irr{A}_1\oplus\irr{A}_2\\
\midrule
\irr{A}_1@2b & 2 & \irr{A}_1\oplus\irr{B}_1 & \irr{E} & \irr{A}_1\oplus\irr{B}_1\\
\irr{A}_2@2b & 2 & \irr{A}_2\oplus\irr{B}_2 & \irr{E} & \irr{A}_2\oplus\irr{B}_2\\
\irr{E}@2b & 4 & \irr{E}_1\oplus\irr{E}_2 & \irr{A}_1\oplus\irr{A}_2\oplus\irr{E} &
   \irr{A}_1\oplus\irr{A}_2\oplus\irr{B}_1\oplus\irr{B}_2\\
\midrule
\irr{A}_1@3c & 3 & \irr{A}_1\oplus\irr{E}_2 & \irr{A}_1\oplus\irr{E} & \irr{A}_1\oplus\irr{B}_1\oplus\irr{B}_2\\
\irr{A}_2@3c & 3 & \irr{A}_2\oplus\irr{E}_2 & \irr{A}_2\oplus\irr{E} & \irr{A}_2\oplus\irr{B}_1\oplus\irr{B}_2\\
\irr{B}_1@3c & 3 & \irr{B}_1\oplus\irr{E}_1 & \irr{A}_2\oplus\irr{E} & \irr{A}_1\oplus\irr{A}_2\oplus\irr{B}_1\\
\irr{B}_2@3c & 3 & \irr{B}_2\oplus\irr{E}_1 & \irr{A}_1\oplus\irr{E} & \irr{A}_1\oplus\irr{A}_2\oplus\irr{B}_2\\
\bottomrule
\end{array}
\end{equation}
Conventions: the $\Cv{6}$ labels are those of \eqref{eq:C6vtable}. At $1a$ the site
irrep $\irr{B}_1$ is even under the bond-containing mirrors $\sv$. At $3c$ the site
irrep $\irr{B}_1$ is even under the mirror containing the bond, i.e.\ it is a
$p$ orbital pointing along the bond. At $M$, $\irr{B}_1$ is odd under $C_2$ and even
under $\sv$. The row $\irr{A}_1@2b$ is exactly the $\uppi$-band result of \cref{ch:bands}.
In the Bilbao labels the same information appears as vectors of multiplicities, one
entry per small irrep; always check the conventions of the table you use.

\section{The TQC criterion}\label{sec:tqc}

Every band representation has an irrep vector that is a sum of EBR vectors with
non-negative integer coefficients. Consequently:
\begin{keyresult}[Topological quantum chemistry]
Take an isolated set of bands (separated by gaps from the rest) and list its small
irreps at the maximal $\vec{k}$ points.
\begin{enumerate}
\item If this irrep vector \emph{cannot} be written as a sum of EBR vectors with
      non-negative integer coefficients, the set is \emph{necessarily} topological:
      no false positives.
\item If it can be written only with some negative coefficients, the set is
      \emph{fragile} topological: it becomes trivial when certain trivial bands are
      added.
\item If it can be written with non-negative coefficients, the set \emph{may} still
      be topological: false negatives exist, because not all topology is
      \emph{symmetry indicated}. Two sets of bands can have identical irreps
      everywhere and different topology (a Chern insulator and a trivial insulator,
      for example).
\end{enumerate}
\end{keyresult}
A second, purely structural statement is often decisive.
\begin{theorem}[Disconnected EBRs]
If the bands of an EBR split into two or more groups separated by gaps, at least one
of the groups is topological.
\end{theorem}
Indeed, if every group were Wannier representable, each would carry a band
representation, and the EBR would be their sum, contradicting its elementarity.

The compatibility relations of \cref{sec:compat} are what make band structures
connected or disconnected. Along each line the subduced irreps must match. If there is
no way of connecting the irreps at the high-symmetry points into two separated groups
that satisfy all compatibility relations, the EBR is forced to be connected. A
disconnection, when allowed, is a signal of topology.

\section{Graphene without spin--orbit coupling}

For the spinless $\uppi$ bands, $\irr{A}_1@2b$ has the two-dimensional irrep $\irr{E}$ at $K$.
The two bands are forced to touch there, so the EBR cannot split: graphene is a
semimetal with symmetry-enforced Dirac points. No gap, no topological insulator.

\section{Graphene with spin--orbit coupling}

With spin, each $p_z$ orbital becomes a Kramers doublet. At the $\Cv{3}$ site,
$\irr{a}_1\otimes D_{1/2}=\bar{\irr{E}}_1$ (\cref{ch:spin}), so the four spinful $\pi$
bands carry the double-valued EBR $\bar{\irr{E}}_1@2b$. Using the double-group
characters of $\Cv{6}$, $\Cv{3}$ and $\Cv{2}$ (\cref{ex:8:spinful}) one finds
\begin{equation}\label{eq:spinfulgraphene}
\bar{\irr{E}}_1@2b:\qquad
\Gamma:\ \bar{\irr{E}}_1\oplus\bar{\irr{E}}_2,\qquad
K:\ \bar{\irr{E}}_1\oplus{}^{1}\bar{\irr{E}}\oplus{}^{2}\bar{\irr{E}},\qquad
M:\ 2\,\bar{\irr{E}} .
\end{equation}
In the Bilbao labels shown in the lecture these are
$\Gamma_7\oplus\Gamma_8$, $K_4\oplus K_5\oplus K_6$ and $2M_5$, with $K_4$ the
two-dimensional irrep at $K$.

With spin--orbit coupling the degeneracy at $K$ need not be fourfold: the four bands can
separate into two isolated pairs. In the lecture's example (a Kane--Mele-type model),
\begin{itemize}
\item one pair carries $\{\Gamma_7\,|\,K_4\,|\,M_5\}$. This is exactly the irrep vector of
      the EBR $\bar{\irr{E}}_1@1a$ (spinful orbitals at the hexagon centres). By irreps
      alone it \emph{could} be trivial: an ``obstructed atomic limit'' with Wannier
      centres in the empty hexagons rather than on the atoms;
\item the other pair carries $\{\Gamma_8\,|\,K_5\oplus K_6\,|\,M_5\}$. The two-band EBRs
      of this group are $\bar{\irr{E}}_1@1a=\{\Gamma_7|K_4|M_5\}$,
      $\bar{\irr{E}}_2@1a=\{\Gamma_8|K_4|M_5\}$ and $\bar{\irr{E}}_3@1a=\{\Gamma_9|K_5\oplus K_6|M_5\}$.
      None of them matches, so this pair is necessarily topological.
\end{itemize}
This is consistent with the theorem on disconnected EBRs. In the Kane--Mele model the
topological pair is the $\mathbb{Z}_2$ quantum spin Hall insulator. There, in fact,
\emph{both} pairs are $\mathbb{Z}_2$ non-trivial. The pair whose irreps mimic an atomic
limit is a false negative of the irrep criterion, a concrete instance of point 3 above.

\begin{beyond}[Beyond the lectures: where to go next]
\begin{itemize}
\item The \texttt{BANDREP}, \texttt{Check Topological Mat.}\ and Topological Materials
      Database tools of the Bilbao server apply this analysis to first-principles band
      structures; programs such as \texttt{IrRep} extract the irreps from DFT output.
\item Symmetry indicators (Po, Vishwanath, Watanabe) are the group of irrep vectors
      modulo band representations; they classify what irreps can diagnose.
\item Magnetic topological quantum chemistry extends everything to the 1651 magnetic
      space groups, where time reversal combines with spatial operations
      (\cref{sec:magnetic}).
\item Review: J.\ Cano and B.\ Bradlyn, Annu.\ Rev.\ Condens.\ Matter Phys.\ 12,
      225 (2021).
\end{itemize}
\end{beyond}

\exercisesection

\begin{exercise}\label{ex:8:A1at1a}
Derive the rows $\irr{A}_1@1a$ and $\irr{E}_1@1a$ of \eqref{eq:ebrp6mm}. For orbitals
at the origin all phases in \eqref{eq:bandchar} are trivial. Why does $\irr{E}_1@1a$
split at $M$?
\end{exercise}

\begin{solution}
An orbital at the origin is mapped onto itself by every point operation, with
$\vec{t}_\alpha=0$, so $\chi_{\vec{k}}(g)=\chi_\rho(g)$: the band representation at $\vec{k}$ is
simply $\rho$ subduced to $\bar{G}_{\vec{k}}$. For $\irr{A}_1$ this is the identity irrep
everywhere. For $\irr{E}_1$: at $\Gamma$, $\irr{E}_1$; at $K$ the characters
$(2,-1,0)$ on $(E,C_3,\sd)$ give $\irr{E}$; at $M$ the characters $(2,-2,0,0)$ on
$(E,C_2,\sv,\sd)$ give $\irr{B}_1\oplus\irr{B}_2$. It splits because $\Cv{2}$ has only
one-dimensional irreps.
\end{solution}

\begin{exercise}\label{ex:8:notsum}
Show that $\irr{A}_1@2b$ is not the sum of EBRs induced from $1a$. Then explain why the
two spinless $\uppi$ bands can never be separated into two isolated bands without
breaking a symmetry of $p6mm$.
\end{exercise}

\begin{solution}
One-band EBRs from $1a$ have one-dimensional irreps at $K$, so any sum of two of them
has two one-dimensional irreps at $K$, not $\irr{E}$. The two-band EBRs from $1a$ have
$\irr{E}_1$ or $\irr{E}_2$ at $\Gamma$, not $\irr{A}_1\oplus\irr{B}_1$. Hence $\irr{A}_1@2b$ is not a
sum of $1a$ EBRs (the Wannier functions of the $\uppi$ bands sit on the atoms). The two
$\uppi$ bands can never be separated: at $K$ they form the two-dimensional irrep $\irr{E}$,
so they touch there for every symmetric Hamiltonian.
\end{solution}

\begin{exercise}\label{ex:8:obstructed}
A hypothetical isolated band of a $p6mm$ crystal has irreps $\irr{B}_1$ at $\Gamma$,
$\irr{A}_2$ at $K$ and $\irr{B}_1$ at $M$. Is it Wannier representable according to the
irrep criterion? Where would its Wannier centres be?
\end{exercise}

\begin{solution}
$\{\irr{B}_1|\irr{A}_2|\irr{B}_1\}$ is exactly the irrep vector of $\irr{B}_1@1a$
\eqref{eq:ebrp6mm}. By the irrep criterion the band is compatible with a Wannier
function centred at the hexagon centres ($1a$), with the symmetry of an
$f$-like orbital (odd under the bond-bisecting mirrors). If the atoms sit at $2b$, such a
band is an \emph{obstructed atomic limit}: trivial, but with Wannier centres away from
the atoms. Irreps cannot exclude hidden topology (false negatives), but nothing in them
requires it.
\end{solution}

\begin{exercise}[spinful $p_z$ orbitals]\label{ex:8:spinful}
The double group of $\Cv{6}$ has three two-dimensional double-valued irreps with
characters (classes $E$, $\Ebar$, $C_2$, $C_3$, $\bar{C}_3$, $C_6$, $\bar{C}_6$, $\sv$, $\sd$)
\[
\begin{array}{c|ccccccccc}
 & E & \Ebar & C_2 & C_3 & \bar{C}_3 & C_6 & \bar{C}_6 & \sv & \sd\\\hline
\bar{\irr{E}}_1 & 2&-2&0&1&-1&\sqrt3&-\sqrt3&0&0\\
\bar{\irr{E}}_2 & 2&-2&0&1&-1&-\sqrt3&\sqrt3&0&0\\
\bar{\irr{E}}_3 & 2&-2&0&-2&2&0&0&0&0
\end{array}
\]
(class sizes $1,1,2,2,2,2,2,6,6$; $C_2$ and $\bar{C}_2$, and each mirror and its barred
partner, fall in the same class).
Using \eqref{eq:bandchar} with the spin character $\chi_{1/2}(\theta)=2\cos(\theta/2)$
for rotations and $0$ for mirrors, derive \eqref{eq:spinfulgraphene}.
\end{exercise}

\begin{solution}
Each site carries $\irr{a}_1\otimes D_{1/2}$; its site character is the spin character.
At $\Gamma$: $E\to4$, $\Ebar\to-4$; $C_2$, $C_6$, $\sd$ exchange sublattices $\to0$; $C_3$
keeps them with trivial phases, $\to2\chi_{1/2}(120^\circ)=2$, and $\bar{C}_3\to-2$;
$\sv\to2\cdot0=0$. The reduction in the double group of $\Cv{6}$ (order 24, class sizes
$1,1,2,2,2,2,2,6,6$) gives
$m_{\bar{\irr{E}}_1}=m_{\bar{\irr{E}}_2}=\tfrac1{24}(8+8+4+4)=1$ and
$m_{\bar{\irr{E}}_3}=\tfrac1{24}(8+8-8-8)=0$.
At $K$: $C_3$ gives $(\omega+\omega^{*})\cdot1=-1$, $\bar{C}_3$ gives $+1$, the mirrors
$0$, so $\chi_K=(4,-4,-1,1,0,0)=\chi_{D_{3/2}}$, i.e.\
$\bar{\irr{E}}_1\oplus{}^{1}\bar{\irr{E}}\oplus{}^{2}\bar{\irr{E}}$ by \eqref{eq:C3vdouble}.
At $M$: $(4,-4,0,0,0)$ on the five classes of the double group of $\Cv{2}$ (order 8),
whose only double-valued irrep $\bar{\irr{E}}$ has $(2,-2,0,0,0)$: $2\,\bar{\irr{E}}$.
\end{solution}


\end{document}
